% !TeX TXS-program:compile = txs:///arara
% arara: pdflatex: {shell: no, synctex: no, interaction: batchmode}
% arara: pdflatex: {shell: no, synctex: no, interaction: batchmode}

\documentclass{article}
\usepackage[margin=1in]{geometry}
\usepackage{bera}
\usepackage{inconsolata}
\usepackage[T1]{fontenc}
\usepackage[scale=0.875]{cabin}
\usepackage{game-icons}
\usepackage{xcolor}
\pagecolor{yellow!10}
\usepackage{longtable}
\usepackage{booktabs}
\usepackage{hyperref}
\usepackage{fancyvrb}
\usepackage{shortvrb}
\MakeShortVerb{\|}
\setlength{\parindent}{0pt}

\newenvironment{gameiconcase}
{%
  \footnotesize
  \begin{longtable}{cllr}%icon + name + macro + page
    \cmidrule[\heavyrulewidth]{1-4}% \toprule
    \bfseries Icon & \bfseries Name & \bfseries Direct command & \bfseries PDF page \\
    \cmidrule{1-4}
    \endhead
}%
{%
    \cmidrule[\heavyrulewidth]{1-4}
  \end{longtable}%
}

\newenvironment{gameiconbadgecase}
{%
  \footnotesize
  \begin{longtable}{cclr}%badge + badge on an icon + name + page
    \cmidrule[\heavyrulewidth]{1-4}
    \bfseries Badge & \bfseries On an icon & \bfseries Name & \bfseries PDF page \\
    \cmidrule{1-4}
    \endhead
}%
{%
    \cmidrule[\heavyrulewidth]{1-4}
  \end{longtable}%
}

\NewDocumentCommand{\showcasegameiconsbadge}{mm}{%name + page
  \textcolor{violet}{\Large\gameiconsbadge{#1}} & \textcolor{violet}{\Large\gameicons[badge=#1]{meeple}} & \ttfamily #1 & \ttfamily #2 \\
}

\NewDocumentCommand{\showcasegameicons}{mm}{%name + page
  \textcolor{violet}{\Large\gameicons{#1}} & \itshape #1 & \ttfamily \textbackslash gameicons\{#1\} & \ttfamily #2 \\
}

\begin{document}

\title{The \textsf{game-icons} package (0.1a) -- 28/10/2026}
\author{%
  Cédric Pierquet\\%
  \url{https://game-icons.net} \\%
  \url{https://github.com/cpierquet/latex-packages/tree/main/game-icons} (\LaTeX{} package)
}
\date{cpierquet -- at -- outlook . fr}\maketitle

This package provides \LaTeX{} support for the \textsf{Game-icons.net} project (\url{https://game-icons.net}).

\medskip

\hspace*{5mm}\fbox{\begin{minipage}{0.75\linewidth}{\emph{A set of 4133 free CC BY 3.0 licensed SVG/PDF icons (and 59 badges), by Lorc, Delapouite \& contributors, for use in your games and projects.}}\end{minipage}}

\medskip

To use \textsf{game-icons} in your document, include the package with |\usepackage{game-icons}|.

An icon can be accessed using the icon name. A list of all included icons with their respective \textit{page} can be found at the end of this document.

\medskip

The individual icons are provided as pages of a PDF and are ultimately loaded using the |\includegraphics| command from the |graphicx| package.

\medskip

The icons are black by default and take the current color, so |\textcolor|, |\color|, etc. work as expected.

\section{Example}

\begin{Verbatim}[frame=single]
...
\usepackage{game-icons}
...
\begin{document}
...
Inline \gameicons{meeple} version, and colored \textcolor{red}{\gameicons{meeple}}
...
\end{document}
\end{Verbatim}

Result: Inline \gameicons{meeple} version, and \textcolor{red}{\gameicons{meeple}} colored version

\pagebreak

\section{Usage}

\subsection{Generic macro}

\begin{Verbatim}[frame=single]
\gameicons[height=...,(d)strut=...]<includegraphics options>{name}
\end{Verbatim}

Key \texttt{[height=...]} can be:

\begin{itemize}
  \item \texttt{dauto} (by default): inline insertion with height/depth;
  \item \texttt{auto}: inline insertion with height/nodepth;
  \item \texttt{a length}: normal insertion with fixed height;
  \item \texttt{manual height/depth}: normal insertion with fixed height/depth.
\end{itemize}

The keys \texttt{[(d)strut=...]} can be used to set the \textit{reference} box for dimension calculations.

\medskip

The given height is the height of the \textit{full} box of the icon (the icon may include a small border).

\medskip

Example: Inline{\setlength\fboxsep{0pt}\fbox{\gameicons{meeple}}}version or inline {\setlength\fboxsep{0pt}\fbox{\gameicons{meeple}}} version

\pagebreak

\subsection{Examples}

\begin{Verbatim}[frame=single]
%normal version: inline with automatic height/depth (from qX characters)
\gameicons[height=dauto]{meeple} %or \gameicons{meeple}
\end{Verbatim}

{\Huge q\gameicons{meeple}X}

\begin{Verbatim}[frame=single]
%normal version: inline with automatic height/depth (from specified characters)
(q\gameicons[height=dauto,dstrut=(qB)]{meeple}B)
\end{Verbatim}

{\Huge (q\gameicons[height=dauto,dstrut=(qB)]{meeple}B)}

\begin{Verbatim}[frame=single]
%normal version: inline with automatic height (from X character)
\gameicons[height=auto]{meeple}
\end{Verbatim}

{\Huge X\gameicons[height=auto]{meeple}X}

\begin{Verbatim}[frame=single]
%normal version: inline with manual height/depth
\gameicons[height={0.8em/-0.4em}]{meeple}
\end{Verbatim}

{\Huge qb\gameicons[height={0.8em/-0.4em}]{meeple}Xz}

\begin{Verbatim}[frame=single]
%fixed height
\gameicons[height=2in]{meeple}
\end{Verbatim}

\gameicons[height=1.5in]{meeple}

\subsection{Badges}

The \textsf{Game-icons.net} badges (59 small symbols) can be added to an icon, as in the Studio of the website, with the key \texttt{[badge=...]}:

\begin{Verbatim}[frame=single]
\gameicons[badge=bolt]{wizard-staff}
\gameicons[badge=number-3,badgepos=tr]{shield}
\gameicons[badge=skull,badgepos=tl,badgecolor=red,badgescale=0.5]{anvil}
\end{Verbatim}

{\Huge q\gameicons{wizard-staff}X\quad q\gameicons[badge=bolt]{wizard-staff}X\quad q\gameicons[badge=number-3,badgepos=tr]{shield}X\quad q\gameicons[badge=skull,badgepos=tl,badgecolor=red,badgescale=0.5]{anvil}X}

\begin{itemize}
  \item \texttt{badge=...}: name of the badge (see the list at the end of this document);
  \item \texttt{badgepos=...}: position on the icon, \texttt{br} (by default), \texttt{bl}, \texttt{tr}, \texttt{tl} (corners) or \texttt{c} (center);
  \item \texttt{badgescale=...}: size of the badge relative to the size of the icon (\texttt{0.4} by default);
  \item \texttt{badgecolor=...}: color of the badge (needs the \texttt{xcolor} or \texttt{color} package), the icon keeps the current color.
\end{itemize}

The height keys (\texttt{height}, \texttt{strut}, \texttt{dstrut}) apply to the icon, and the badge follows its size. A badge alone can also be used with \verb|\gameiconsbadge[height=...]<includegraphics options>{name}|: \gameiconsbadge{bolt}, \textcolor{red}{\gameiconsbadge{skull}}.

\subsection{Used packages}

\texttt{graphicx} is loaded and used by the package (the key handling relies on \texttt{expl3}, which is part of the \LaTeX{} kernel since October 2020).

\subsection{License and attribution}

The package is released under the \LaTeX{} Project Public License v1.3c or later.

\medskip

The icons are provided by \textsf{Game-icons.net} (\url{https://game-icons.net}) under the \textsf{CC BY 3.0} license (\url{https://creativecommons.org/licenses/by/3.0/}), except the icons of \textsf{Viscious Speed} and \textsf{Zeromancer}, which are \textsf{CC0}. The attribution is required: the authors must be credited when the icons are used, for example \textit{``Icons made by \{author\}. Available on \url{https://game-icons.net}''}, as requested in the \texttt{license.txt} file of the icons repository (\url{https://github.com/game-icons/icons}).

\medskip

Each icon was drawn by one of the following contributors (one folder per contributor in the repository):

\medskip

{\small
Lorc, Delapouite, John Colburn, Felbrigg, John Redman, Carl Olsen, Sbed, PriorBlue, Willdabeast, Viscious Speed (CC0), Lord Berandas, Irongamer, HeavenlyDog, Lucas, Faithtoken, Skoll, Andy Meneely, Cathelineau, Kier Heyl, Aussiesim, Sparker, Zeromancer (CC0), Rihlsul, Quoting, Guard13007, DarkZaitzev, SpencerDub, GeneralAce135, Zajkonur, Catsu, Starseeker, Pepijn Poolman, Pierre Leducq, Caro Asercion, SeregaCthtuf.\par}

\medskip

The \textsf{badges} are also provided by \textsf{Game-icons.net}. The author of each icon can be found on \url{https://game-icons.net}.

\subsection{Icons with the same name}

The repository contains 4180 icons, but only 4133 different names: 47 names exist twice, drawn by two different contributors (for example \texttt{castle}, \texttt{bat}, \texttt{fairy}, \texttt{mouse} or \texttt{screw}). This package provides \textbf{one icon per name}: when a name exists in several folders, the icon of the first contributor in alphabetical order of the folder names is kept.

\subsection{Bugs}

For bug reports and feature requests, report on the GitHub repository \url{https://github.com/cpierquet/latex-packages/issues}.

\newpage

\section{List of icons}

The \textsf{PDF} file is \texttt{gameicons-all}.

Icons are colored in \textcolor{violet}{violet}.

\begin{gameiconcase}
\showcasegameicons{3d-glasses}{1}
\showcasegameicons{3d-hammer}{2}
\showcasegameicons{3d-meeple}{3}
\showcasegameicons{3d-stairs}{4}
\showcasegameicons{abacus}{5}
\showcasegameicons{abbot-meeple}{6}
\showcasegameicons{abdominal-armor}{7}
\showcasegameicons{abstract-001}{8}
\showcasegameicons{abstract-002}{9}
\showcasegameicons{abstract-003}{10}
\showcasegameicons{abstract-004}{11}
\showcasegameicons{abstract-005}{12}
\showcasegameicons{abstract-006}{13}
\showcasegameicons{abstract-007}{14}
\showcasegameicons{abstract-008}{15}
\showcasegameicons{abstract-009}{16}
\showcasegameicons{abstract-010}{17}
\showcasegameicons{abstract-011}{18}
\showcasegameicons{abstract-012}{19}
\showcasegameicons{abstract-013}{20}
\showcasegameicons{abstract-014}{21}
\showcasegameicons{abstract-015}{22}
\showcasegameicons{abstract-016}{23}
\showcasegameicons{abstract-017}{24}
\showcasegameicons{abstract-018}{25}
\showcasegameicons{abstract-019}{26}
\showcasegameicons{abstract-020}{27}
\showcasegameicons{abstract-021}{28}
\showcasegameicons{abstract-022}{29}
\showcasegameicons{abstract-023}{30}
\showcasegameicons{abstract-024}{31}
\showcasegameicons{abstract-025}{32}
\showcasegameicons{abstract-026}{33}
\showcasegameicons{abstract-027}{34}
\showcasegameicons{abstract-028}{35}
\showcasegameicons{abstract-029}{36}
\showcasegameicons{abstract-030}{37}
\showcasegameicons{abstract-031}{38}
\showcasegameicons{abstract-032}{39}
\showcasegameicons{abstract-033}{40}
\showcasegameicons{abstract-034}{41}
\showcasegameicons{abstract-035}{42}
\showcasegameicons{abstract-036}{43}
\showcasegameicons{abstract-037}{44}
\showcasegameicons{abstract-038}{45}
\showcasegameicons{abstract-039}{46}
\showcasegameicons{abstract-040}{47}
\showcasegameicons{abstract-041}{48}
\showcasegameicons{abstract-042}{49}
\showcasegameicons{abstract-043}{50}
\showcasegameicons{abstract-044}{51}
\showcasegameicons{abstract-045}{52}
\showcasegameicons{abstract-046}{53}
\showcasegameicons{abstract-047}{54}
\showcasegameicons{abstract-048}{55}
\showcasegameicons{abstract-049}{56}
\showcasegameicons{abstract-050}{57}
\showcasegameicons{abstract-051}{58}
\showcasegameicons{abstract-052}{59}
\showcasegameicons{abstract-053}{60}
\showcasegameicons{abstract-054}{61}
\showcasegameicons{abstract-055}{62}
\showcasegameicons{abstract-056}{63}
\showcasegameicons{abstract-057}{64}
\showcasegameicons{abstract-058}{65}
\showcasegameicons{abstract-059}{66}
\showcasegameicons{abstract-060}{67}
\showcasegameicons{abstract-061}{68}
\showcasegameicons{abstract-062}{69}
\showcasegameicons{abstract-063}{70}
\showcasegameicons{abstract-064}{71}
\showcasegameicons{abstract-065}{72}
\showcasegameicons{abstract-066}{73}
\showcasegameicons{abstract-067}{74}
\showcasegameicons{abstract-068}{75}
\showcasegameicons{abstract-069}{76}
\showcasegameicons{abstract-070}{77}
\showcasegameicons{abstract-071}{78}
\showcasegameicons{abstract-072}{79}
\showcasegameicons{abstract-073}{80}
\showcasegameicons{abstract-074}{81}
\showcasegameicons{abstract-075}{82}
\showcasegameicons{abstract-076}{83}
\showcasegameicons{abstract-077}{84}
\showcasegameicons{abstract-078}{85}
\showcasegameicons{abstract-079}{86}
\showcasegameicons{abstract-080}{87}
\showcasegameicons{abstract-081}{88}
\showcasegameicons{abstract-082}{89}
\showcasegameicons{abstract-083}{90}
\showcasegameicons{abstract-084}{91}
\showcasegameicons{abstract-085}{92}
\showcasegameicons{abstract-086}{93}
\showcasegameicons{abstract-087}{94}
\showcasegameicons{abstract-088}{95}
\showcasegameicons{abstract-089}{96}
\showcasegameicons{abstract-090}{97}
\showcasegameicons{abstract-091}{98}
\showcasegameicons{abstract-092}{99}
\showcasegameicons{abstract-093}{100}
\showcasegameicons{abstract-094}{101}
\showcasegameicons{abstract-095}{102}
\showcasegameicons{abstract-096}{103}
\showcasegameicons{abstract-097}{104}
\showcasegameicons{abstract-098}{105}
\showcasegameicons{abstract-099}{106}
\showcasegameicons{abstract-100}{107}
\showcasegameicons{abstract-101}{108}
\showcasegameicons{abstract-102}{109}
\showcasegameicons{abstract-103}{110}
\showcasegameicons{abstract-104}{111}
\showcasegameicons{abstract-105}{112}
\showcasegameicons{abstract-106}{113}
\showcasegameicons{abstract-107}{114}
\showcasegameicons{abstract-108}{115}
\showcasegameicons{abstract-109}{116}
\showcasegameicons{abstract-110}{117}
\showcasegameicons{abstract-111}{118}
\showcasegameicons{abstract-112}{119}
\showcasegameicons{abstract-113}{120}
\showcasegameicons{abstract-114}{121}
\showcasegameicons{abstract-115}{122}
\showcasegameicons{abstract-116}{123}
\showcasegameicons{abstract-117}{124}
\showcasegameicons{abstract-118}{125}
\showcasegameicons{abstract-119}{126}
\showcasegameicons{abstract-120}{127}
\showcasegameicons{abstract-121}{128}
\showcasegameicons{accordion}{129}
\showcasegameicons{ace}{130}
\showcasegameicons{achievement}{131}
\showcasegameicons{achilles-heel}{132}
\showcasegameicons{acid-blob}{133}
\showcasegameicons{acid-shield}{134}
\showcasegameicons{acid-tube}{135}
\showcasegameicons{acid}{136}
\showcasegameicons{acorn}{137}
\showcasegameicons{acoustic-megaphone}{138}
\showcasegameicons{acrobatic}{139}
\showcasegameicons{aerial-signal}{140}
\showcasegameicons{aero-bike}{141}
\showcasegameicons{aerodynamic-harpoon}{142}
\showcasegameicons{aerosol}{143}
\showcasegameicons{africa}{144}
\showcasegameicons{afterburn}{145}
\showcasegameicons{agave}{146}
\showcasegameicons{ages}{147}
\showcasegameicons{air-balloon}{148}
\showcasegameicons{air-force}{149}
\showcasegameicons{air-man}{150}
\showcasegameicons{air-zigzag}{151}
\showcasegameicons{airplane-arrival}{152}
\showcasegameicons{airplane-departure}{153}
\showcasegameicons{airplane}{154}
\showcasegameicons{airtight-hatch}{155}
\showcasegameicons{ak47}{156}
\showcasegameicons{ak47u}{157}
\showcasegameicons{akhet}{158}
\showcasegameicons{alarm-clock}{159}
\showcasegameicons{algae}{160}
\showcasegameicons{algeria}{161}
\showcasegameicons{alien-bug}{162}
\showcasegameicons{alien-egg}{163}
\showcasegameicons{alien-fire}{164}
\showcasegameicons{alien-skull}{165}
\showcasegameicons{alien-stare}{166}
\showcasegameicons{all-for-one}{167}
\showcasegameicons{all-seeing-eye}{168}
\showcasegameicons{allied-star}{169}
\showcasegameicons{alligator-clip}{170}
\showcasegameicons{almond}{171}
\showcasegameicons{amber-mosquito}{172}
\showcasegameicons{ambulance}{173}
\showcasegameicons{american-football-ball}{174}
\showcasegameicons{american-football-helmet}{175}
\showcasegameicons{american-football-player}{176}
\showcasegameicons{american-shield}{177}
\showcasegameicons{amethyst}{178}
\showcasegameicons{ammo-box}{179}
\showcasegameicons{ammonite-fossil}{180}
\showcasegameicons{ammonite}{181}
\showcasegameicons{amphora}{182}
\showcasegameicons{ample-dress}{183}
\showcasegameicons{amplitude}{184}
\showcasegameicons{amputation}{185}
\showcasegameicons{anarchy}{186}
\showcasegameicons{anatomy}{187}
\showcasegameicons{anchor}{188}
\showcasegameicons{ancient-columns}{189}
\showcasegameicons{ancient-ruins}{190}
\showcasegameicons{ancient-screw}{191}
\showcasegameicons{ancient-sword}{192}
\showcasegameicons{android-mask}{193}
\showcasegameicons{andromeda-chain}{194}
\showcasegameicons{angel-outfit}{195}
\showcasegameicons{angel-wings}{196}
\showcasegameicons{angler-fish}{197}
\showcasegameicons{angola}{198}
\showcasegameicons{angry-eyes}{199}
\showcasegameicons{angular-spider}{200}
\showcasegameicons{animal-hide}{201}
\showcasegameicons{animal-skull}{202}
\showcasegameicons{ankh}{203}
\showcasegameicons{annexation}{204}
\showcasegameicons{ant}{205}
\showcasegameicons{antarctica}{206}
\showcasegameicons{anteater}{207}
\showcasegameicons{anthem}{208}
\showcasegameicons{anti-aircraft-gun}{209}
\showcasegameicons{antibody}{210}
\showcasegameicons{anticlockwise-rotation}{211}
\showcasegameicons{ants}{212}
\showcasegameicons{anubis}{213}
\showcasegameicons{anvil-impact}{214}
\showcasegameicons{anvil}{215}
\showcasegameicons{apc}{216}
\showcasegameicons{apollo-capsule}{217}
\showcasegameicons{apothecary}{218}
\showcasegameicons{apple-core}{219}
\showcasegameicons{apple-maggot}{220}
\showcasegameicons{apple-seeds}{221}
\showcasegameicons{aquarium}{222}
\showcasegameicons{aquarius}{223}
\showcasegameicons{aqueduct}{224}
\showcasegameicons{arabic-door}{225}
\showcasegameicons{arc-triomphe}{226}
\showcasegameicons{arch-bridge}{227}
\showcasegameicons{archaeopteryx-fossil}{228}
\showcasegameicons{archer}{229}
\showcasegameicons{archery-target}{230}
\showcasegameicons{architect-mask}{231}
\showcasegameicons{archive-register}{232}
\showcasegameicons{archive-research}{233}
\showcasegameicons{arcing-bolt}{234}
\showcasegameicons{arena}{235}
\showcasegameicons{aries}{236}
\showcasegameicons{arm-bandage}{237}
\showcasegameicons{arm-sling}{238}
\showcasegameicons{arm}{239}
\showcasegameicons{armadillo-tail}{240}
\showcasegameicons{armadillo}{241}
\showcasegameicons{armchair}{242}
\showcasegameicons{armor-blueprint}{243}
\showcasegameicons{armor-cuisses}{244}
\showcasegameicons{armor-downgrade}{245}
\showcasegameicons{armor-punch}{246}
\showcasegameicons{armor-upgrade}{247}
\showcasegameicons{armor-vest}{248}
\showcasegameicons{armored-boomerang}{249}
\showcasegameicons{armored-pants}{250}
\showcasegameicons{armoured-shell}{251}
\showcasegameicons{arrest}{252}
\showcasegameicons{arrow-cluster}{253}
\showcasegameicons{arrow-cursor}{254}
\showcasegameicons{arrow-dunk}{255}
\showcasegameicons{arrow-flights}{256}
\showcasegameicons{arrow-scope}{257}
\showcasegameicons{arrow-wings}{258}
\showcasegameicons{arrowed}{259}
\showcasegameicons{arrowhead}{260}
\showcasegameicons{arrows-shield}{261}
\showcasegameicons{arson}{262}
\showcasegameicons{artichoke}{263}
\showcasegameicons{artificial-hive}{264}
\showcasegameicons{artificial-intelligence}{265}
\showcasegameicons{artillery-shell}{266}
\showcasegameicons{ascending-block}{267}
\showcasegameicons{asian-lantern}{268}
\showcasegameicons{asparagus}{269}
\showcasegameicons{aspergillum}{270}
\showcasegameicons{assassin-pocket}{271}
\showcasegameicons{asteroid}{272}
\showcasegameicons{astrolabe}{273}
\showcasegameicons{astronaut-helmet}{274}
\showcasegameicons{at-sea}{275}
\showcasegameicons{atlas}{276}
\showcasegameicons{atom-core}{277}
\showcasegameicons{atom}{278}
\showcasegameicons{atomic-slashes}{279}
\showcasegameicons{attached-shield}{280}
\showcasegameicons{attack-gauge}{281}
\showcasegameicons{aubergine}{282}
\showcasegameicons{audio-cassette}{283}
\showcasegameicons{aura}{284}
\showcasegameicons{australia}{285}
\showcasegameicons{auto-repair}{286}
\showcasegameicons{autogun}{287}
\showcasegameicons{automatic-sas}{288}
\showcasegameicons{avocado}{289}
\showcasegameicons{avoidance}{290}
\showcasegameicons{awareness}{291}
\showcasegameicons{axe-in-log}{292}
\showcasegameicons{axe-in-stump}{293}
\showcasegameicons{axe-swing}{294}
\showcasegameicons{axe-sword}{295}
\showcasegameicons{axolotl}{296}
\showcasegameicons{aztec-calendar-sun}{297}
\showcasegameicons{azul-flake}{298}
\showcasegameicons{baby-bottle}{299}
\showcasegameicons{baby-face}{300}
\showcasegameicons{babyfoot-players}{301}
\showcasegameicons{back-forth}{302}
\showcasegameicons{back-pain}{303}
\showcasegameicons{backbone-shell}{304}
\showcasegameicons{backgammon}{305}
\showcasegameicons{backpack}{306}
\showcasegameicons{backstab}{307}
\showcasegameicons{backup}{308}
\showcasegameicons{backward-time}{309}
\showcasegameicons{bacon}{310}
\showcasegameicons{bad-breath}{311}
\showcasegameicons{bad-gnome}{312}
\showcasegameicons{badger}{313}
\showcasegameicons{bagpipes}{314}
\showcasegameicons{balaclava}{315}
\showcasegameicons{balkenkreuz}{316}
\showcasegameicons{ball-glow}{317}
\showcasegameicons{ball-heart}{318}
\showcasegameicons{ball-pyramid}{319}
\showcasegameicons{ballerina-shoes}{320}
\showcasegameicons{ballista}{321}
\showcasegameicons{balloon-dog}{322}
\showcasegameicons{balloons}{323}
\showcasegameicons{bamboo-fountain}{324}
\showcasegameicons{bamboo}{325}
\showcasegameicons{banana-bunch}{326}
\showcasegameicons{banana-peel}{327}
\showcasegameicons{banana-peeled}{328}
\showcasegameicons{banana}{329}
\showcasegameicons{bandage-roll}{330}
\showcasegameicons{bandaged}{331}
\showcasegameicons{bandana}{332}
\showcasegameicons{bandit}{333}
\showcasegameicons{banging-gavel}{334}
\showcasegameicons{banjo}{335}
\showcasegameicons{bank}{336}
\showcasegameicons{banknote}{337}
\showcasegameicons{baobab}{338}
\showcasegameicons{bar-stool}{339}
\showcasegameicons{barbarian}{340}
\showcasegameicons{barbecue}{341}
\showcasegameicons{barbed-arrow}{342}
\showcasegameicons{barbed-coil}{343}
\showcasegameicons{barbed-nails}{344}
\showcasegameicons{barbed-spear}{345}
\showcasegameicons{barbed-star}{346}
\showcasegameicons{barbed-sun}{347}
\showcasegameicons{barbed-wire}{348}
\showcasegameicons{barbute}{349}
\showcasegameicons{barefoot}{350}
\showcasegameicons{barn-owl}{351}
\showcasegameicons{barn}{352}
\showcasegameicons{barracks-tent}{353}
\showcasegameicons{barracks}{354}
\showcasegameicons{barrel-leak}{355}
\showcasegameicons{barrel}{356}
\showcasegameicons{barricade}{357}
\showcasegameicons{barrier}{358}
\showcasegameicons{base-dome}{359}
\showcasegameicons{baseball-bat}{360}
\showcasegameicons{baseball-glove}{361}
\showcasegameicons{basket}{362}
\showcasegameicons{basketball-ball}{363}
\showcasegameicons{basketball-basket}{364}
\showcasegameicons{basketball-jersey}{365}
\showcasegameicons{basset-hound-head}{366}
\showcasegameicons{bassoon}{367}
\showcasegameicons{bastet}{368}
\showcasegameicons{bat-blade}{369}
\showcasegameicons{bat-leth}{370}
\showcasegameicons{bat-mask}{371}
\showcasegameicons{bat-wing}{372}
\showcasegameicons{bat}{373}
\showcasegameicons{bathtub}{374}
\showcasegameicons{baton}{375}
\showcasegameicons{battered-axe}{376}
\showcasegameicons{batteries}{377}
\showcasegameicons{battery-0}{378}
\showcasegameicons{battery-25}{379}
\showcasegameicons{battery-50}{380}
\showcasegameicons{battery-75}{381}
\showcasegameicons{battery-100}{382}
\showcasegameicons{battery-minus}{383}
\showcasegameicons{battery-pack-alt}{384}
\showcasegameicons{battery-pack}{385}
\showcasegameicons{battery-plus}{386}
\showcasegameicons{battle-axe}{387}
\showcasegameicons{battle-gear}{388}
\showcasegameicons{battle-mech}{389}
\showcasegameicons{battle-tank}{390}
\showcasegameicons{battleship}{391}
\showcasegameicons{batwing-emblem}{392}
\showcasegameicons{bayonet}{393}
\showcasegameicons{beach-bag}{394}
\showcasegameicons{beach-ball}{395}
\showcasegameicons{beach-bucket}{396}
\showcasegameicons{beam-satellite}{397}
\showcasegameicons{beam-wake}{398}
\showcasegameicons{beams-aura}{399}
\showcasegameicons{beanstalk}{400}
\showcasegameicons{bear-face}{401}
\showcasegameicons{bear-head}{402}
\showcasegameicons{beard}{403}
\showcasegameicons{beast-eye}{404}
\showcasegameicons{beaver}{405}
\showcasegameicons{bed-lamp}{406}
\showcasegameicons{bed}{407}
\showcasegameicons{bee}{408}
\showcasegameicons{beech}{409}
\showcasegameicons{beehive}{410}
\showcasegameicons{beer-bottle}{411}
\showcasegameicons{beer-horn}{412}
\showcasegameicons{beer-stein}{413}
\showcasegameicons{beet}{414}
\showcasegameicons{beetle-shell}{415}
\showcasegameicons{behold}{416}
\showcasegameicons{belgium}{417}
\showcasegameicons{bell-pepper}{418}
\showcasegameicons{bell-shield}{419}
\showcasegameicons{bellows}{420}
\showcasegameicons{belt-armor}{421}
\showcasegameicons{belt-buckles}{422}
\showcasegameicons{belt}{423}
\showcasegameicons{berries-bowl}{424}
\showcasegameicons{berry-bush}{425}
\showcasegameicons{bestial-fangs}{426}
\showcasegameicons{beveled-star}{427}
\showcasegameicons{biceps}{428}
\showcasegameicons{big-diamond-ring}{429}
\showcasegameicons{big-egg}{430}
\showcasegameicons{big-gear}{431}
\showcasegameicons{big-wave}{432}
\showcasegameicons{billed-cap}{433}
\showcasegameicons{bindle}{434}
\showcasegameicons{binoculars}{435}
\showcasegameicons{biohazard}{436}
\showcasegameicons{biplane}{437}
\showcasegameicons{birch-trees}{438}
\showcasegameicons{bird-cage}{439}
\showcasegameicons{bird-claw}{440}
\showcasegameicons{bird-house}{441}
\showcasegameicons{bird-limb}{442}
\showcasegameicons{bird-mask}{443}
\showcasegameicons{bird-scepter}{444}
\showcasegameicons{bird-twitter}{445}
\showcasegameicons{bison}{446}
\showcasegameicons{black-bar}{447}
\showcasegameicons{black-belt}{448}
\showcasegameicons{black-book}{449}
\showcasegameicons{black-bridge}{450}
\showcasegameicons{black-flag}{451}
\showcasegameicons{black-hand-shield}{452}
\showcasegameicons{black-hole-bolas}{453}
\showcasegameicons{black-knight-helm}{454}
\showcasegameicons{black-sea}{455}
\showcasegameicons{blackball}{456}
\showcasegameicons{blackcurrant}{457}
\showcasegameicons{blacksmith}{458}
\showcasegameicons{blade-bite}{459}
\showcasegameicons{blade-drag}{460}
\showcasegameicons{blade-fall}{461}
\showcasegameicons{blanket}{462}
\showcasegameicons{blast}{463}
\showcasegameicons{blaster}{464}
\showcasegameicons{bleeding-eye}{465}
\showcasegameicons{bleeding-heart}{466}
\showcasegameicons{bleeding-wound}{467}
\showcasegameicons{blender}{468}
\showcasegameicons{blindfold}{469}
\showcasegameicons{block-house}{470}
\showcasegameicons{blood}{471}
\showcasegameicons{bloody-stash}{472}
\showcasegameicons{bloody-sword}{473}
\showcasegameicons{blunderbuss}{474}
\showcasegameicons{bo}{475}
\showcasegameicons{boar-ensign}{476}
\showcasegameicons{boar-tusks}{477}
\showcasegameicons{boar}{478}
\showcasegameicons{boarding-pass}{479}
\showcasegameicons{boat-engine}{480}
\showcasegameicons{boat-fishing}{481}
\showcasegameicons{boat-horizon}{482}
\showcasegameicons{boat-propeller}{483}
\showcasegameicons{boba}{484}
\showcasegameicons{body-balance}{485}
\showcasegameicons{body-height}{486}
\showcasegameicons{body-swapping}{487}
\showcasegameicons{boiling-bubbles}{488}
\showcasegameicons{bok-choy}{489}
\showcasegameicons{bolas}{490}
\showcasegameicons{bolivia}{491}
\showcasegameicons{bolt-bomb}{492}
\showcasegameicons{bolt-cutter}{493}
\showcasegameicons{bolt-drop}{494}
\showcasegameicons{bolt-eye}{495}
\showcasegameicons{bolt-saw}{496}
\showcasegameicons{bolt-shield}{497}
\showcasegameicons{bolt-spell-cast}{498}
\showcasegameicons{bolter-gun}{499}
\showcasegameicons{bomber}{500}
\showcasegameicons{bombing-run}{501}
\showcasegameicons{bone-gnawer}{502}
\showcasegameicons{bone-knife}{503}
\showcasegameicons{bone-mace}{504}
\showcasegameicons{bonsai-tree}{505}
\showcasegameicons{book-aura}{506}
\showcasegameicons{book-cover}{507}
\showcasegameicons{book-pile}{508}
\showcasegameicons{book-storm}{509}
\showcasegameicons{bookmark}{510}
\showcasegameicons{bookmarklet}{511}
\showcasegameicons{bookshelf}{512}
\showcasegameicons{boombox}{513}
\showcasegameicons{boomerang-cross}{514}
\showcasegameicons{boomerang-sun}{515}
\showcasegameicons{boomerang}{516}
\showcasegameicons{boot-kick}{517}
\showcasegameicons{boot-prints}{518}
\showcasegameicons{boot-stomp}{519}
\showcasegameicons{boots}{520}
\showcasegameicons{booze}{521}
\showcasegameicons{bordered-shield}{522}
\showcasegameicons{boss-key}{523}
\showcasegameicons{bottle-cap}{524}
\showcasegameicons{bottle-vapors}{525}
\showcasegameicons{bottled-bolt}{526}
\showcasegameicons{bottled-shadow}{527}
\showcasegameicons{bottom-right-3d-arrow}{528}
\showcasegameicons{boulder-dash}{529}
\showcasegameicons{bouncing-spring}{530}
\showcasegameicons{bouncing-sword}{531}
\showcasegameicons{bow-arrow}{532}
\showcasegameicons{bow-string}{533}
\showcasegameicons{bow-tie-ribbon}{534}
\showcasegameicons{bow-tie}{535}
\showcasegameicons{bowels}{536}
\showcasegameicons{bowen-knot}{537}
\showcasegameicons{bowie-knife}{538}
\showcasegameicons{bowl-of-rice}{539}
\showcasegameicons{bowl-spiral}{540}
\showcasegameicons{bowling-alley}{541}
\showcasegameicons{bowling-pin}{542}
\showcasegameicons{bowling-propulsion}{543}
\showcasegameicons{bowling-strike}{544}
\showcasegameicons{bowman}{545}
\showcasegameicons{box-cutter}{546}
\showcasegameicons{box-trap}{547}
\showcasegameicons{box-unpacking}{548}
\showcasegameicons{boxing-glove-surprise}{549}
\showcasegameicons{boxing-glove}{550}
\showcasegameicons{boxing-ring}{551}
\showcasegameicons{bracer}{552}
\showcasegameicons{bracers}{553}
\showcasegameicons{brain-dump}{554}
\showcasegameicons{brain-freeze}{555}
\showcasegameicons{brain-leak}{556}
\showcasegameicons{brain-stem}{557}
\showcasegameicons{brain-tentacle}{558}
\showcasegameicons{brain}{559}
\showcasegameicons{brainstorm}{560}
\showcasegameicons{branch-arrow}{561}
\showcasegameicons{brandy-bottle}{562}
\showcasegameicons{brasero}{563}
\showcasegameicons{brass-eye}{564}
\showcasegameicons{brass-knuckles}{565}
\showcasegameicons{brazil-flag}{566}
\showcasegameicons{brazil}{567}
\showcasegameicons{bread-slice}{568}
\showcasegameicons{bread}{569}
\showcasegameicons{breaking-chain}{570}
\showcasegameicons{breastplate}{571}
\showcasegameicons{brick-pile}{572}
\showcasegameicons{brick-wall}{573}
\showcasegameicons{bridge}{574}
\showcasegameicons{briefcase}{575}
\showcasegameicons{bright-explosion}{576}
\showcasegameicons{broad-dagger}{577}
\showcasegameicons{broadhead-arrow}{578}
\showcasegameicons{broadsword}{579}
\showcasegameicons{broccoli}{580}
\showcasegameicons{brodie-helmet}{581}
\showcasegameicons{broken-arrow}{582}
\showcasegameicons{broken-axe}{583}
\showcasegameicons{broken-bone}{584}
\showcasegameicons{broken-bottle}{585}
\showcasegameicons{broken-heart-zone}{586}
\showcasegameicons{broken-heart}{587}
\showcasegameicons{broken-pottery}{588}
\showcasegameicons{broken-ribbon}{589}
\showcasegameicons{broken-shield}{590}
\showcasegameicons{broken-skull}{591}
\showcasegameicons{broken-tablet}{592}
\showcasegameicons{broken-wall}{593}
\showcasegameicons{broom}{594}
\showcasegameicons{brutal-helm}{595}
\showcasegameicons{brute}{596}
\showcasegameicons{bubble-field}{597}
\showcasegameicons{bubbles}{598}
\showcasegameicons{bubbling-beam}{599}
\showcasegameicons{bubbling-bowl}{600}
\showcasegameicons{bubbling-flask}{601}
\showcasegameicons{bucket-wheel-excavator}{602}
\showcasegameicons{bud}{603}
\showcasegameicons{buffalo-head}{604}
\showcasegameicons{bug-net}{605}
\showcasegameicons{bugle-call}{606}
\showcasegameicons{bulb}{607}
\showcasegameicons{bulgaria}{608}
\showcasegameicons{bull-horns}{609}
\showcasegameicons{bull}{610}
\showcasegameicons{bulldozer}{611}
\showcasegameicons{bullet-bill}{612}
\showcasegameicons{bullet-impacts}{613}
\showcasegameicons{bullets}{614}
\showcasegameicons{bullseye}{615}
\showcasegameicons{bully-minion}{616}
\showcasegameicons{bundle-grenade}{617}
\showcasegameicons{bunk-beds}{618}
\showcasegameicons{bunker-assault}{619}
\showcasegameicons{bunker}{620}
\showcasegameicons{bunny-slippers}{621}
\showcasegameicons{buoy}{622}
\showcasegameicons{burn}{623}
\showcasegameicons{burning-blobs}{624}
\showcasegameicons{burning-book}{625}
\showcasegameicons{burning-dot}{626}
\showcasegameicons{burning-embers}{627}
\showcasegameicons{burning-eye}{628}
\showcasegameicons{burning-forest}{629}
\showcasegameicons{burning-meteor}{630}
\showcasegameicons{burning-passion}{631}
\showcasegameicons{burning-round-shot}{632}
\showcasegameicons{burning-skull}{633}
\showcasegameicons{burning-tree}{634}
\showcasegameicons{burst-blob}{635}
\showcasegameicons{bus-doors}{636}
\showcasegameicons{bus-stop}{637}
\showcasegameicons{bus}{638}
\showcasegameicons{butter-toast}{639}
\showcasegameicons{butter}{640}
\showcasegameicons{butterfly-flower}{641}
\showcasegameicons{butterfly-knife}{642}
\showcasegameicons{butterfly-warning}{643}
\showcasegameicons{butterfly}{644}
\showcasegameicons{button-finger}{645}
\showcasegameicons{buy-card}{646}
\showcasegameicons{byzantin-temple}{647}
\showcasegameicons{c96}{648}
\showcasegameicons{cabbage}{649}
\showcasegameicons{cable-stayed-bridge}{650}
\showcasegameicons{cactus-pot}{651}
\showcasegameicons{cactus-tap}{652}
\showcasegameicons{cactus}{653}
\showcasegameicons{cadillac-helm}{654}
\showcasegameicons{caduceus}{655}
\showcasegameicons{caesar}{656}
\showcasegameicons{cage}{657}
\showcasegameicons{caged-ball}{658}
\showcasegameicons{cake-slice}{659}
\showcasegameicons{calavera}{660}
\showcasegameicons{calculator}{661}
\showcasegameicons{caldera}{662}
\showcasegameicons{calendar-half-year}{663}
\showcasegameicons{calendar}{664}
\showcasegameicons{caltrops}{665}
\showcasegameicons{camargue-cross}{666}
\showcasegameicons{cambodia}{667}
\showcasegameicons{camel-head}{668}
\showcasegameicons{camel}{669}
\showcasegameicons{camp-cooking-pot}{670}
\showcasegameicons{campfire}{671}
\showcasegameicons{camping-tent}{672}
\showcasegameicons{canadian-goose}{673}
\showcasegameicons{cancel}{674}
\showcasegameicons{cancer}{675}
\showcasegameicons{candle-flame}{676}
\showcasegameicons{candle-holder}{677}
\showcasegameicons{candle-light}{678}
\showcasegameicons{candle-skull}{679}
\showcasegameicons{candlebright}{680}
\showcasegameicons{candles}{681}
\showcasegameicons{candlestick-phone}{682}
\showcasegameicons{candy-canes}{683}
\showcasegameicons{canned-fish}{684}
\showcasegameicons{cannister}{685}
\showcasegameicons{cannon-ball}{686}
\showcasegameicons{cannon-shot}{687}
\showcasegameicons{cannon}{688}
\showcasegameicons{canoe}{689}
\showcasegameicons{cantua}{690}
\showcasegameicons{cape-armor}{691}
\showcasegameicons{cape}{692}
\showcasegameicons{capitol}{693}
\showcasegameicons{capricorn}{694}
\showcasegameicons{captain-hat-profile}{695}
\showcasegameicons{capybara}{696}
\showcasegameicons{car-battery}{697}
\showcasegameicons{car-door}{698}
\showcasegameicons{car-key}{699}
\showcasegameicons{car-seat}{700}
\showcasegameicons{car-wheel}{701}
\showcasegameicons{carabiner}{702}
\showcasegameicons{carambola}{703}
\showcasegameicons{caravan}{704}
\showcasegameicons{caravel}{705}
\showcasegameicons{card-2-clubs}{706}
\showcasegameicons{card-2-diamonds}{707}
\showcasegameicons{card-2-hearts}{708}
\showcasegameicons{card-2-spades}{709}
\showcasegameicons{card-3-clubs}{710}
\showcasegameicons{card-3-diamonds}{711}
\showcasegameicons{card-3-hearts}{712}
\showcasegameicons{card-3-spades}{713}
\showcasegameicons{card-4-clubs}{714}
\showcasegameicons{card-4-diamonds}{715}
\showcasegameicons{card-4-hearts}{716}
\showcasegameicons{card-4-spades}{717}
\showcasegameicons{card-5-clubs}{718}
\showcasegameicons{card-5-diamonds}{719}
\showcasegameicons{card-5-hearts}{720}
\showcasegameicons{card-5-spades}{721}
\showcasegameicons{card-6-clubs}{722}
\showcasegameicons{card-6-diamonds}{723}
\showcasegameicons{card-6-hearts}{724}
\showcasegameicons{card-6-spades}{725}
\showcasegameicons{card-7-clubs}{726}
\showcasegameicons{card-7-diamonds}{727}
\showcasegameicons{card-7-hearts}{728}
\showcasegameicons{card-7-spades}{729}
\showcasegameicons{card-8-clubs}{730}
\showcasegameicons{card-8-diamonds}{731}
\showcasegameicons{card-8-hearts}{732}
\showcasegameicons{card-8-spades}{733}
\showcasegameicons{card-9-clubs}{734}
\showcasegameicons{card-9-diamonds}{735}
\showcasegameicons{card-9-hearts}{736}
\showcasegameicons{card-9-spades}{737}
\showcasegameicons{card-10-clubs}{738}
\showcasegameicons{card-10-diamonds}{739}
\showcasegameicons{card-10-hearts}{740}
\showcasegameicons{card-10-spades}{741}
\showcasegameicons{card-ace-clubs}{742}
\showcasegameicons{card-ace-diamonds}{743}
\showcasegameicons{card-ace-hearts}{744}
\showcasegameicons{card-ace-spades}{745}
\showcasegameicons{card-burn}{746}
\showcasegameicons{card-discard}{747}
\showcasegameicons{card-draw}{748}
\showcasegameicons{card-exchange}{749}
\showcasegameicons{card-jack-clubs}{750}
\showcasegameicons{card-jack-diamonds}{751}
\showcasegameicons{card-jack-hearts}{752}
\showcasegameicons{card-jack-spades}{753}
\showcasegameicons{card-joker}{754}
\showcasegameicons{card-king-clubs}{755}
\showcasegameicons{card-king-diamonds}{756}
\showcasegameicons{card-king-hearts}{757}
\showcasegameicons{card-king-spades}{758}
\showcasegameicons{card-pick}{759}
\showcasegameicons{card-pickup}{760}
\showcasegameicons{card-play}{761}
\showcasegameicons{card-queen-clubs}{762}
\showcasegameicons{card-queen-diamonds}{763}
\showcasegameicons{card-queen-hearts}{764}
\showcasegameicons{card-queen-spades}{765}
\showcasegameicons{card-random}{766}
\showcasegameicons{cardboard-box-closed}{767}
\showcasegameicons{cardboard-box}{768}
\showcasegameicons{cargo-crane}{769}
\showcasegameicons{cargo-crate}{770}
\showcasegameicons{cargo-ship}{771}
\showcasegameicons{carillon}{772}
\showcasegameicons{carnival-mask}{773}
\showcasegameicons{carnivore-mouth}{774}
\showcasegameicons{carnivorous-plant}{775}
\showcasegameicons{carnyx}{776}
\showcasegameicons{carousel}{777}
\showcasegameicons{carpet-bombing}{778}
\showcasegameicons{carrier}{779}
\showcasegameicons{carrion}{780}
\showcasegameicons{carrot}{781}
\showcasegameicons{cartwheel}{782}
\showcasegameicons{cash}{783}
\showcasegameicons{cassowary-head}{784}
\showcasegameicons{castle-ruins}{785}
\showcasegameicons{castle}{786}
\showcasegameicons{cat}{787}
\showcasegameicons{catapult}{788}
\showcasegameicons{catch}{789}
\showcasegameicons{caterpillar}{790}
\showcasegameicons{cauldron}{791}
\showcasegameicons{cavalry}{792}
\showcasegameicons{cave-entrance}{793}
\showcasegameicons{caveman}{794}
\showcasegameicons{cctv-camera}{795}
\showcasegameicons{ceiling-barnacle}{796}
\showcasegameicons{ceiling-light}{797}
\showcasegameicons{celebration-fire}{798}
\showcasegameicons{cellar-barrels}{799}
\showcasegameicons{cement-shoes}{800}
\showcasegameicons{centaur-heart}{801}
\showcasegameicons{centaur}{802}
\showcasegameicons{centipede}{803}
\showcasegameicons{centurion-helmet}{804}
\showcasegameicons{ceremonial-mask}{805}
\showcasegameicons{chain-lightning}{806}
\showcasegameicons{chain-mail}{807}
\showcasegameicons{chained-arrow-heads}{808}
\showcasegameicons{chained-heart}{809}
\showcasegameicons{chaingun}{810}
\showcasegameicons{chainsaw}{811}
\showcasegameicons{chakram}{812}
\showcasegameicons{chalice-drops}{813}
\showcasegameicons{chalk-outline-murder}{814}
\showcasegameicons{chameleon-glyph}{815}
\showcasegameicons{champagne-cork}{816}
\showcasegameicons{champions}{817}
\showcasegameicons{chanterelles}{818}
\showcasegameicons{character}{819}
\showcasegameicons{charcuterie}{820}
\showcasegameicons{charged-arrow}{821}
\showcasegameicons{charging-bull}{822}
\showcasegameicons{charging}{823}
\showcasegameicons{chariot}{824}
\showcasegameicons{charm}{825}
\showcasegameicons{chart}{826}
\showcasegameicons{chat-bubble}{827}
\showcasegameicons{check-mark}{828}
\showcasegameicons{checkbox-tree}{829}
\showcasegameicons{checked-shield}{830}
\showcasegameicons{checkered-diamond}{831}
\showcasegameicons{checkered-flag}{832}
\showcasegameicons{checklist}{833}
\showcasegameicons{cheerful}{834}
\showcasegameicons{cheese-wedge}{835}
\showcasegameicons{chef-toque}{836}
\showcasegameicons{chelsea-boot}{837}
\showcasegameicons{chemical-arrow}{838}
\showcasegameicons{chemical-bolt}{839}
\showcasegameicons{chemical-drop}{840}
\showcasegameicons{chemical-tank}{841}
\showcasegameicons{cherish}{842}
\showcasegameicons{cherry}{843}
\showcasegameicons{chess-bishop}{844}
\showcasegameicons{chess-king}{845}
\showcasegameicons{chess-knight}{846}
\showcasegameicons{chess-pawn}{847}
\showcasegameicons{chess-queen}{848}
\showcasegameicons{chess-rook}{849}
\showcasegameicons{chest-armor}{850}
\showcasegameicons{chest}{851}
\showcasegameicons{chestnut-leaf}{852}
\showcasegameicons{chewed-heart}{853}
\showcasegameicons{chewed-skull}{854}
\showcasegameicons{chicken-leg}{855}
\showcasegameicons{chicken-oven}{856}
\showcasegameicons{chicken}{857}
\showcasegameicons{chili-pepper}{858}
\showcasegameicons{chimney}{859}
\showcasegameicons{chips-bag}{860}
\showcasegameicons{chisel}{861}
\showcasegameicons{chocolate-bar}{862}
\showcasegameicons{choice}{863}
\showcasegameicons{chopped-skull}{864}
\showcasegameicons{chopsticks}{865}
\showcasegameicons{church}{866}
\showcasegameicons{cigale}{867}
\showcasegameicons{cigar}{868}
\showcasegameicons{cigarette}{869}
\showcasegameicons{circle-cage}{870}
\showcasegameicons{circle-claws}{871}
\showcasegameicons{circle-forest}{872}
\showcasegameicons{circle-sparks}{873}
\showcasegameicons{circle}{874}
\showcasegameicons{circling-fish}{875}
\showcasegameicons{circuitry}{876}
\showcasegameicons{circular-saw}{877}
\showcasegameicons{circular-sawblade}{878}
\showcasegameicons{city-car}{879}
\showcasegameicons{clamp}{880}
\showcasegameicons{clapperboard}{881}
\showcasegameicons{clarinet}{882}
\showcasegameicons{classical-knowledge}{883}
\showcasegameicons{claw-hammer}{884}
\showcasegameicons{claw-slashes}{885}
\showcasegameicons{claw-string}{886}
\showcasegameicons{claw}{887}
\showcasegameicons{claws}{888}
\showcasegameicons{clay-brick}{889}
\showcasegameicons{claymore-explosive}{890}
\showcasegameicons{cleaver}{891}
\showcasegameicons{cleopatra}{892}
\showcasegameicons{click}{893}
\showcasegameicons{cliff-crossing}{894}
\showcasegameicons{cloak-dagger}{895}
\showcasegameicons{cloak}{896}
\showcasegameicons{cloaked-figure-on-horseback}{897}
\showcasegameicons{clock-tower}{898}
\showcasegameicons{clockwise-rotation}{899}
\showcasegameicons{clockwork}{900}
\showcasegameicons{closed-barbute}{901}
\showcasegameicons{closed-doors}{902}
\showcasegameicons{cloth-jar}{903}
\showcasegameicons{clothes}{904}
\showcasegameicons{clothesline}{905}
\showcasegameicons{clothespin}{906}
\showcasegameicons{cloud-download}{907}
\showcasegameicons{cloud-ring}{908}
\showcasegameicons{cloud-upload}{909}
\showcasegameicons{cloudy-fork}{910}
\showcasegameicons{clout}{911}
\showcasegameicons{clover-spiked}{912}
\showcasegameicons{clover}{913}
\showcasegameicons{clown}{914}
\showcasegameicons{clownfish}{915}
\showcasegameicons{clubs}{916}
\showcasegameicons{cluster-bomb}{917}
\showcasegameicons{coa-de-jima}{918}
\showcasegameicons{coal-pile}{919}
\showcasegameicons{coal-wagon}{920}
\showcasegameicons{cobra}{921}
\showcasegameicons{cobweb}{922}
\showcasegameicons{coconuts}{923}
\showcasegameicons{coffee-beans}{924}
\showcasegameicons{coffee-cup}{925}
\showcasegameicons{coffee-mug}{926}
\showcasegameicons{coffee-pot}{927}
\showcasegameicons{coffin}{928}
\showcasegameicons{cog-lock}{929}
\showcasegameicons{cog}{930}
\showcasegameicons{cogsplosion}{931}
\showcasegameicons{coiled-nail}{932}
\showcasegameicons{coiling-curl}{933}
\showcasegameicons{coinflip}{934}
\showcasegameicons{coins-pile}{935}
\showcasegameicons{coins}{936}
\showcasegameicons{cold-heart}{937}
\showcasegameicons{coliseum}{938}
\showcasegameicons{colombia}{939}
\showcasegameicons{colombian-statue}{940}
\showcasegameicons{colt-m1911}{941}
\showcasegameicons{column-vase}{942}
\showcasegameicons{coma}{943}
\showcasegameicons{comb}{944}
\showcasegameicons{combination-lock}{945}
\showcasegameicons{comet-spark}{946}
\showcasegameicons{commercial-airplane}{947}
\showcasegameicons{compact-disc}{948}
\showcasegameicons{companion-cube}{949}
\showcasegameicons{compass}{950}
\showcasegameicons{computer-fan}{951}
\showcasegameicons{computing}{952}
\showcasegameicons{concentration-orb}{953}
\showcasegameicons{concentric-crescents}{954}
\showcasegameicons{concrete-bag}{955}
\showcasegameicons{condor-emblem}{956}
\showcasegameicons{condylura-skull}{957}
\showcasegameicons{confirmed}{958}
\showcasegameicons{confrontation}{959}
\showcasegameicons{congress}{960}
\showcasegameicons{conqueror}{961}
\showcasegameicons{console-controller}{962}
\showcasegameicons{contortionist}{963}
\showcasegameicons{contract}{964}
\showcasegameicons{control-tower}{965}
\showcasegameicons{convergence-target}{966}
\showcasegameicons{conversation}{967}
\showcasegameicons{converse-shoe}{968}
\showcasegameicons{convict}{969}
\showcasegameicons{convince}{970}
\showcasegameicons{conway-life-glider}{971}
\showcasegameicons{cook}{972}
\showcasegameicons{cookie}{973}
\showcasegameicons{cooking-glove}{974}
\showcasegameicons{cooking-pot}{975}
\showcasegameicons{cool-spices}{976}
\showcasegameicons{cooler}{977}
\showcasegameicons{cootie-catcher}{978}
\showcasegameicons{coral}{979}
\showcasegameicons{cork-hat}{980}
\showcasegameicons{corked-tube}{981}
\showcasegameicons{corkscrew}{982}
\showcasegameicons{corn}{983}
\showcasegameicons{corner-explosion}{984}
\showcasegameicons{corner-flag}{985}
\showcasegameicons{cornucopia}{986}
\showcasegameicons{coronation}{987}
\showcasegameicons{corporal}{988}
\showcasegameicons{corset}{989}
\showcasegameicons{corsica}{990}
\showcasegameicons{cosmic-egg}{991}
\showcasegameicons{cotton-flower}{992}
\showcasegameicons{covered-jar}{993}
\showcasegameicons{cow}{994}
\showcasegameicons{cowboy-boot}{995}
\showcasegameicons{cowboy-holster}{996}
\showcasegameicons{cowled}{997}
\showcasegameicons{cpu-shot}{998}
\showcasegameicons{cpu}{999}
\showcasegameicons{crab-claw}{1000}
\showcasegameicons{crab}{1001}
\showcasegameicons{cracked-alien-skull}{1002}
\showcasegameicons{cracked-ball-dunk}{1003}
\showcasegameicons{cracked-disc}{1004}
\showcasegameicons{cracked-glass}{1005}
\showcasegameicons{cracked-helm}{1006}
\showcasegameicons{cracked-mask}{1007}
\showcasegameicons{cracked-saber}{1008}
\showcasegameicons{cracked-shield}{1009}
\showcasegameicons{crafting}{1010}
\showcasegameicons{crags}{1011}
\showcasegameicons{crane}{1012}
\showcasegameicons{credits-currency}{1013}
\showcasegameicons{crenel-crown}{1014}
\showcasegameicons{crenulated-shield}{1015}
\showcasegameicons{crescent-blade}{1016}
\showcasegameicons{crescent-staff}{1017}
\showcasegameicons{crested-helmet}{1018}
\showcasegameicons{cricket-bat}{1019}
\showcasegameicons{cricket}{1020}
\showcasegameicons{crime-scene-tape}{1021}
\showcasegameicons{croc-jaws}{1022}
\showcasegameicons{croc-sword}{1023}
\showcasegameicons{croissant}{1024}
\showcasegameicons{croissants-pupil}{1025}
\showcasegameicons{crook-flail}{1026}
\showcasegameicons{cross-flare}{1027}
\showcasegameicons{cross-mark}{1028}
\showcasegameicons{cross-shield}{1029}
\showcasegameicons{crossbow}{1030}
\showcasegameicons{crosscut-saw}{1031}
\showcasegameicons{crossed-air-flows}{1032}
\showcasegameicons{crossed-axes}{1033}
\showcasegameicons{crossed-bones}{1034}
\showcasegameicons{crossed-chains}{1035}
\showcasegameicons{crossed-claws}{1036}
\showcasegameicons{crossed-pistols}{1037}
\showcasegameicons{crossed-sabres}{1038}
\showcasegameicons{crossed-slashes}{1039}
\showcasegameicons{crossed-swords}{1040}
\showcasegameicons{crosshair-arrow}{1041}
\showcasegameicons{crosshair}{1042}
\showcasegameicons{crossroad}{1043}
\showcasegameicons{crow-dive}{1044}
\showcasegameicons{crow-nest}{1045}
\showcasegameicons{crowbar}{1046}
\showcasegameicons{crown-coin}{1047}
\showcasegameicons{crown-of-thorns}{1048}
\showcasegameicons{crown}{1049}
\showcasegameicons{crowned-explosion}{1050}
\showcasegameicons{crowned-heart}{1051}
\showcasegameicons{crowned-skull}{1052}
\showcasegameicons{crucifix}{1053}
\showcasegameicons{cruiser}{1054}
\showcasegameicons{crumbling-ball}{1055}
\showcasegameicons{crush}{1056}
\showcasegameicons{cryo-chamber}{1057}
\showcasegameicons{crypt-entrance}{1058}
\showcasegameicons{crystal-ball}{1059}
\showcasegameicons{crystal-bars}{1060}
\showcasegameicons{crystal-cluster}{1061}
\showcasegameicons{crystal-earrings}{1062}
\showcasegameicons{crystal-eye}{1063}
\showcasegameicons{crystal-growth}{1064}
\showcasegameicons{crystal-shine}{1065}
\showcasegameicons{crystal-shrine}{1066}
\showcasegameicons{crystal-wand}{1067}
\showcasegameicons{crystalize}{1068}
\showcasegameicons{cuauhtli}{1069}
\showcasegameicons{cube}{1070}
\showcasegameicons{cubeforce}{1071}
\showcasegameicons{cubes}{1072}
\showcasegameicons{cuckoo-clock}{1073}
\showcasegameicons{cultist}{1074}
\showcasegameicons{cupcake}{1075}
\showcasegameicons{cupidon-arrow}{1076}
\showcasegameicons{curled-leaf}{1077}
\showcasegameicons{curled-tentacle}{1078}
\showcasegameicons{curling-stone}{1079}
\showcasegameicons{curling-vines}{1080}
\showcasegameicons{curly-mask}{1081}
\showcasegameicons{curly-wing}{1082}
\showcasegameicons{cursed-star}{1083}
\showcasegameicons{curvy-knife}{1084}
\showcasegameicons{custodian-helmet}{1085}
\showcasegameicons{cut-diamond}{1086}
\showcasegameicons{cut-lemon}{1087}
\showcasegameicons{cut-palm}{1088}
\showcasegameicons{cyber-eye}{1089}
\showcasegameicons{cyborg-face}{1090}
\showcasegameicons{cycle}{1091}
\showcasegameicons{cycling}{1092}
\showcasegameicons{cyclops}{1093}
\showcasegameicons{cz-skorpion}{1094}
\showcasegameicons{d4}{1095}
\showcasegameicons{d10}{1096}
\showcasegameicons{d12}{1097}
\showcasegameicons{daemon-pull}{1098}
\showcasegameicons{daemon-skull}{1099}
\showcasegameicons{dagger-rose}{1100}
\showcasegameicons{daggers}{1101}
\showcasegameicons{daisy}{1102}
\showcasegameicons{dam}{1103}
\showcasegameicons{damaged-house}{1104}
\showcasegameicons{dandelion-flower}{1105}
\showcasegameicons{dango}{1106}
\showcasegameicons{dark-squad}{1107}
\showcasegameicons{dart}{1108}
\showcasegameicons{database}{1109}
\showcasegameicons{dead-eye}{1110}
\showcasegameicons{dead-head}{1111}
\showcasegameicons{dead-wood}{1112}
\showcasegameicons{deadly-strike}{1113}
\showcasegameicons{death-juice}{1114}
\showcasegameicons{death-note}{1115}
\showcasegameicons{death-skull}{1116}
\showcasegameicons{death-star}{1117}
\showcasegameicons{death-zone}{1118}
\showcasegameicons{deathcab}{1119}
\showcasegameicons{decapitation}{1120}
\showcasegameicons{deer-head}{1121}
\showcasegameicons{deer-track}{1122}
\showcasegameicons{deer}{1123}
\showcasegameicons{defense-satellite}{1124}
\showcasegameicons{defensive-wall}{1125}
\showcasegameicons{defibrilate}{1126}
\showcasegameicons{deku-tree}{1127}
\showcasegameicons{delicate-perfume}{1128}
\showcasegameicons{delighted}{1129}
\showcasegameicons{delivery-drone}{1130}
\showcasegameicons{demolish}{1131}
\showcasegameicons{dervish-swords}{1132}
\showcasegameicons{desert-camp}{1133}
\showcasegameicons{desert-eagle}{1134}
\showcasegameicons{desert-skull}{1135}
\showcasegameicons{desert}{1136}
\showcasegameicons{deshret-red-crown}{1137}
\showcasegameicons{desk-lamp}{1138}
\showcasegameicons{desk}{1139}
\showcasegameicons{despair}{1140}
\showcasegameicons{detonator}{1141}
\showcasegameicons{detour}{1142}
\showcasegameicons{devil-mask}{1143}
\showcasegameicons{dew}{1144}
\showcasegameicons{diablo-skull}{1145}
\showcasegameicons{diagram}{1146}
\showcasegameicons{dial-padlock}{1147}
\showcasegameicons{diamond-hard}{1148}
\showcasegameicons{diamond-hilt}{1149}
\showcasegameicons{diamond-ring}{1150}
\showcasegameicons{diamond-trophy}{1151}
\showcasegameicons{diamonds-smile}{1152}
\showcasegameicons{diamonds}{1153}
\showcasegameicons{dice-eight-faces-eight}{1154}
\showcasegameicons{dice-fire}{1155}
\showcasegameicons{dice-shield}{1156}
\showcasegameicons{dice-six-faces-five}{1157}
\showcasegameicons{dice-six-faces-four}{1158}
\showcasegameicons{dice-six-faces-one}{1159}
\showcasegameicons{dice-six-faces-six}{1160}
\showcasegameicons{dice-six-faces-three}{1161}
\showcasegameicons{dice-six-faces-two}{1162}
\showcasegameicons{dice-target}{1163}
\showcasegameicons{dice-twenty-faces-one}{1164}
\showcasegameicons{dice-twenty-faces-twenty}{1165}
\showcasegameicons{dig-dug}{1166}
\showcasegameicons{dig-hole}{1167}
\showcasegameicons{digital-trace}{1168}
\showcasegameicons{dimetrodon}{1169}
\showcasegameicons{dinosaur-bones}{1170}
\showcasegameicons{dinosaur-egg}{1171}
\showcasegameicons{dinosaur-rex}{1172}
\showcasegameicons{diplodocus}{1173}
\showcasegameicons{diploma}{1174}
\showcasegameicons{direction-sign}{1175}
\showcasegameicons{direction-signs}{1176}
\showcasegameicons{director-chair}{1177}
\showcasegameicons{direwolf}{1178}
\showcasegameicons{disc-golf-bag}{1179}
\showcasegameicons{disc-golf-basket}{1180}
\showcasegameicons{disc}{1181}
\showcasegameicons{discobolus}{1182}
\showcasegameicons{discussion}{1183}
\showcasegameicons{disintegrate}{1184}
\showcasegameicons{distraction}{1185}
\showcasegameicons{distress-signal}{1186}
\showcasegameicons{divergence}{1187}
\showcasegameicons{divert}{1188}
\showcasegameicons{divided-spiral}{1189}
\showcasegameicons{divided-square}{1190}
\showcasegameicons{diving-dagger}{1191}
\showcasegameicons{diving-helmet}{1192}
\showcasegameicons{djed-pillar}{1193}
\showcasegameicons{djembe}{1194}
\showcasegameicons{djinn}{1195}
\showcasegameicons{dna1}{1196}
\showcasegameicons{dna2}{1197}
\showcasegameicons{doctor-face}{1198}
\showcasegameicons{dodge}{1199}
\showcasegameicons{dodging}{1200}
\showcasegameicons{dodo}{1201}
\showcasegameicons{dog-bowl}{1202}
\showcasegameicons{dog-house}{1203}
\showcasegameicons{dolmen}{1204}
\showcasegameicons{dolphin}{1205}
\showcasegameicons{domino-mask}{1206}
\showcasegameicons{domino-tiles}{1207}
\showcasegameicons{doner-kebab}{1208}
\showcasegameicons{donkey}{1209}
\showcasegameicons{donut}{1210}
\showcasegameicons{door-handle}{1211}
\showcasegameicons{door-ring-handle}{1212}
\showcasegameicons{door-watcher}{1213}
\showcasegameicons{door}{1214}
\showcasegameicons{doorway}{1215}
\showcasegameicons{dorsal-scales}{1216}
\showcasegameicons{double-diaphragm}{1217}
\showcasegameicons{double-dragon}{1218}
\showcasegameicons{double-face-mask}{1219}
\showcasegameicons{double-fish}{1220}
\showcasegameicons{double-necklace}{1221}
\showcasegameicons{double-quaver}{1222}
\showcasegameicons{double-ringed-orb}{1223}
\showcasegameicons{double-shot}{1224}
\showcasegameicons{double-street-lights}{1225}
\showcasegameicons{doubled}{1226}
\showcasegameicons{dough-roller}{1227}
\showcasegameicons{dove}{1228}
\showcasegameicons{dozen}{1229}
\showcasegameicons{dragon-balls}{1230}
\showcasegameicons{dragon-breath}{1231}
\showcasegameicons{dragon-head}{1232}
\showcasegameicons{dragon-orb}{1233}
\showcasegameicons{dragon-shield}{1234}
\showcasegameicons{dragon-spiral}{1235}
\showcasegameicons{dragonfly}{1236}
\showcasegameicons{drakkar-dragon}{1237}
\showcasegameicons{drakkar}{1238}
\showcasegameicons{drama-masks}{1239}
\showcasegameicons{drawbridge}{1240}
\showcasegameicons{dread-skull}{1241}
\showcasegameicons{dread}{1242}
\showcasegameicons{dreadnought}{1243}
\showcasegameicons{dream-catcher}{1244}
\showcasegameicons{dress}{1245}
\showcasegameicons{drill}{1246}
\showcasegameicons{drink-me}{1247}
\showcasegameicons{drinking}{1248}
\showcasegameicons{dripping-blade}{1249}
\showcasegameicons{dripping-goo}{1250}
\showcasegameicons{dripping-honey}{1251}
\showcasegameicons{dripping-knife}{1252}
\showcasegameicons{dripping-star}{1253}
\showcasegameicons{dripping-stone}{1254}
\showcasegameicons{dripping-sword}{1255}
\showcasegameicons{dripping-tube}{1256}
\showcasegameicons{drop-earrings}{1257}
\showcasegameicons{drop-weapon}{1258}
\showcasegameicons{drop}{1259}
\showcasegameicons{droplet-splash}{1260}
\showcasegameicons{droplets}{1261}
\showcasegameicons{drowning}{1262}
\showcasegameicons{drum-kit}{1263}
\showcasegameicons{drum}{1264}
\showcasegameicons{duality-mask}{1265}
\showcasegameicons{duality}{1266}
\showcasegameicons{duck-palm}{1267}
\showcasegameicons{duck}{1268}
\showcasegameicons{duel}{1269}
\showcasegameicons{duffel-bag}{1270}
\showcasegameicons{dumpling-bao}{1271}
\showcasegameicons{dumpling}{1272}
\showcasegameicons{dunce-cap}{1273}
\showcasegameicons{dungeon-gate}{1274}
\showcasegameicons{dungeon-light}{1275}
\showcasegameicons{duration}{1276}
\showcasegameicons{dust-cloud}{1277}
\showcasegameicons{dutch-bike}{1278}
\showcasegameicons{dwarf-face}{1279}
\showcasegameicons{dwarf-helmet}{1280}
\showcasegameicons{dwarf-king}{1281}
\showcasegameicons{dwennimmen}{1282}
\showcasegameicons{dynamite}{1283}
\showcasegameicons{eagle-emblem}{1284}
\showcasegameicons{eagle-head}{1285}
\showcasegameicons{earbuds}{1286}
\showcasegameicons{earrings}{1287}
\showcasegameicons{earth-africa-europe}{1288}
\showcasegameicons{earth-america}{1289}
\showcasegameicons{earth-asia-oceania}{1290}
\showcasegameicons{earth-crack}{1291}
\showcasegameicons{earth-spit}{1292}
\showcasegameicons{earth-worm}{1293}
\showcasegameicons{earwig}{1294}
\showcasegameicons{easel}{1295}
\showcasegameicons{easter-egg}{1296}
\showcasegameicons{eating-pelican}{1297}
\showcasegameicons{eating}{1298}
\showcasegameicons{echo-ripples}{1299}
\showcasegameicons{eclipse-flare}{1300}
\showcasegameicons{eclipse-saw}{1301}
\showcasegameicons{eclipse}{1302}
\showcasegameicons{ecology}{1303}
\showcasegameicons{edge-crack}{1304}
\showcasegameicons{edged-shield}{1305}
\showcasegameicons{eel}{1306}
\showcasegameicons{egg-clutch}{1307}
\showcasegameicons{egg-defense}{1308}
\showcasegameicons{egg-eye}{1309}
\showcasegameicons{egg-pod}{1310}
\showcasegameicons{egypt}{1311}
\showcasegameicons{egyptian-bird}{1312}
\showcasegameicons{egyptian-profile}{1313}
\showcasegameicons{egyptian-pyramids}{1314}
\showcasegameicons{egyptian-sphinx}{1315}
\showcasegameicons{egyptian-temple}{1316}
\showcasegameicons{egyptian-urns}{1317}
\showcasegameicons{egyptian-walk}{1318}
\showcasegameicons{eight-ball}{1319}
\showcasegameicons{elbow-pad}{1320}
\showcasegameicons{elderberry}{1321}
\showcasegameicons{electric-whip}{1322}
\showcasegameicons{electric}{1323}
\showcasegameicons{electrical-crescent}{1324}
\showcasegameicons{electrical-resistance}{1325}
\showcasegameicons{electrical-socket}{1326}
\showcasegameicons{elephant-head}{1327}
\showcasegameicons{elephant}{1328}
\showcasegameicons{elevator}{1329}
\showcasegameicons{elf-ear}{1330}
\showcasegameicons{elf-helmet}{1331}
\showcasegameicons{elven-castle}{1332}
\showcasegameicons{elysium-shade}{1333}
\showcasegameicons{ember-shot}{1334}
\showcasegameicons{embrassed-energy}{1335}
\showcasegameicons{embryo}{1336}
\showcasegameicons{emerald-necklace}{1337}
\showcasegameicons{emerald}{1338}
\showcasegameicons{empty-chessboard}{1339}
\showcasegameicons{empty-hourglass}{1340}
\showcasegameicons{empty-metal-bucket-handle}{1341}
\showcasegameicons{empty-metal-bucket}{1342}
\showcasegameicons{empty-wood-bucket-handle}{1343}
\showcasegameicons{empty-wood-bucket}{1344}
\showcasegameicons{encirclement}{1345}
\showcasegameicons{encrypted-channel}{1346}
\showcasegameicons{energise}{1347}
\showcasegameicons{energy-arrow}{1348}
\showcasegameicons{energy-breath}{1349}
\showcasegameicons{energy-shield}{1350}
\showcasegameicons{energy-sword}{1351}
\showcasegameicons{energy-tank}{1352}
\showcasegameicons{engagement-ring}{1353}
\showcasegameicons{enlightenment}{1354}
\showcasegameicons{enrage}{1355}
\showcasegameicons{ent-mouth}{1356}
\showcasegameicons{entangled-typhoon}{1357}
\showcasegameicons{entry-door}{1358}
\showcasegameicons{envelope}{1359}
\showcasegameicons{erlenmeyer}{1360}
\showcasegameicons{ermine}{1361}
\showcasegameicons{eruption}{1362}
\showcasegameicons{escalator}{1363}
\showcasegameicons{eternal-love}{1364}
\showcasegameicons{european-flag}{1365}
\showcasegameicons{evasion}{1366}
\showcasegameicons{evil-bat}{1367}
\showcasegameicons{evil-book}{1368}
\showcasegameicons{evil-bud}{1369}
\showcasegameicons{evil-comet}{1370}
\showcasegameicons{evil-eyes}{1371}
\showcasegameicons{evil-fork}{1372}
\showcasegameicons{evil-hand}{1373}
\showcasegameicons{evil-love}{1374}
\showcasegameicons{evil-minion}{1375}
\showcasegameicons{evil-moon}{1376}
\showcasegameicons{evil-tower}{1377}
\showcasegameicons{evil-tree}{1378}
\showcasegameicons{evil-wings}{1379}
\showcasegameicons{executioner-hood}{1380}
\showcasegameicons{exit-door}{1381}
\showcasegameicons{expand}{1382}
\showcasegameicons{expanded-rays}{1383}
\showcasegameicons{expander}{1384}
\showcasegameicons{expense}{1385}
\showcasegameicons{exploding-planet}{1386}
\showcasegameicons{explosion-rays}{1387}
\showcasegameicons{explosive-materials}{1388}
\showcasegameicons{explosive-meeting}{1389}
\showcasegameicons{extra-lucid}{1390}
\showcasegameicons{extra-time}{1391}
\showcasegameicons{extra-vision}{1392}
\showcasegameicons{extraction-orb}{1393}
\showcasegameicons{eye-of-horus}{1394}
\showcasegameicons{eye-shield}{1395}
\showcasegameicons{eye-target}{1396}
\showcasegameicons{eyeball}{1397}
\showcasegameicons{eyedropper}{1398}
\showcasegameicons{eyelashes}{1399}
\showcasegameicons{eyepatch}{1400}
\showcasegameicons{eyestalk}{1401}
\showcasegameicons{f-clef}{1402}
\showcasegameicons{f1-car}{1403}
\showcasegameicons{face-to-face}{1404}
\showcasegameicons{factory-arm}{1405}
\showcasegameicons{factory}{1406}
\showcasegameicons{fairy-wand}{1407}
\showcasegameicons{fairy-wings}{1408}
\showcasegameicons{fairy}{1409}
\showcasegameicons{falcon-moon}{1410}
\showcasegameicons{fall-down}{1411}
\showcasegameicons{falling-blob}{1412}
\showcasegameicons{falling-bomb}{1413}
\showcasegameicons{falling-boulder}{1414}
\showcasegameicons{falling-eye}{1415}
\showcasegameicons{falling-leaf}{1416}
\showcasegameicons{falling-ovoid}{1417}
\showcasegameicons{falling-rocks}{1418}
\showcasegameicons{falling-star}{1419}
\showcasegameicons{falling}{1420}
\showcasegameicons{fallout-shelter}{1421}
\showcasegameicons{famas}{1422}
\showcasegameicons{family-house}{1423}
\showcasegameicons{family-tree}{1424}
\showcasegameicons{fanged-skull}{1425}
\showcasegameicons{fangs-circle}{1426}
\showcasegameicons{fangs}{1427}
\showcasegameicons{farm-tractor}{1428}
\showcasegameicons{farmer}{1429}
\showcasegameicons{fast-arrow}{1430}
\showcasegameicons{fast-backward-button}{1431}
\showcasegameicons{fast-forward-button}{1432}
\showcasegameicons{fast-noodles}{1433}
\showcasegameicons{fat}{1434}
\showcasegameicons{feather-necklace}{1435}
\showcasegameicons{feather-wound}{1436}
\showcasegameicons{feather}{1437}
\showcasegameicons{feathered-wing}{1438}
\showcasegameicons{fedora}{1439}
\showcasegameicons{feline}{1440}
\showcasegameicons{female-legs}{1441}
\showcasegameicons{female}{1442}
\showcasegameicons{fencer}{1443}
\showcasegameicons{fern}{1444}
\showcasegameicons{fertilizer-bag}{1445}
\showcasegameicons{fetus}{1446}
\showcasegameicons{fez}{1447}
\showcasegameicons{field-gun}{1448}
\showcasegameicons{field}{1449}
\showcasegameicons{figurehead}{1450}
\showcasegameicons{files}{1451}
\showcasegameicons{film-projector}{1452}
\showcasegameicons{film-spool}{1453}
\showcasegameicons{film-strip}{1454}
\showcasegameicons{finch}{1455}
\showcasegameicons{finger-print}{1456}
\showcasegameicons{fingernail}{1457}
\showcasegameicons{fingers-crossed}{1458}
\showcasegameicons{finish-line}{1459}
\showcasegameicons{fire-ace}{1460}
\showcasegameicons{fire-axe}{1461}
\showcasegameicons{fire-bomb}{1462}
\showcasegameicons{fire-bottle}{1463}
\showcasegameicons{fire-bowl}{1464}
\showcasegameicons{fire-breath}{1465}
\showcasegameicons{fire-dash}{1466}
\showcasegameicons{fire-extinguisher}{1467}
\showcasegameicons{fire-flower}{1468}
\showcasegameicons{fire-gem}{1469}
\showcasegameicons{fire-iris}{1470}
\showcasegameicons{fire-punch}{1471}
\showcasegameicons{fire-ray}{1472}
\showcasegameicons{fire-ring}{1473}
\showcasegameicons{fire-shield}{1474}
\showcasegameicons{fire-shrine}{1475}
\showcasegameicons{fire-silhouette}{1476}
\showcasegameicons{fire-spell-cast}{1477}
\showcasegameicons{fire-tail}{1478}
\showcasegameicons{fire-wave}{1479}
\showcasegameicons{fire-zone}{1480}
\showcasegameicons{fire}{1481}
\showcasegameicons{fireball}{1482}
\showcasegameicons{fireflake}{1483}
\showcasegameicons{fireplace}{1484}
\showcasegameicons{firewall}{1485}
\showcasegameicons{firework-rocket}{1486}
\showcasegameicons{first-aid-kit}{1487}
\showcasegameicons{fish-bucket}{1488}
\showcasegameicons{fish-cooked}{1489}
\showcasegameicons{fish-corpse}{1490}
\showcasegameicons{fish-eggs}{1491}
\showcasegameicons{fish-escape}{1492}
\showcasegameicons{fish-monster}{1493}
\showcasegameicons{fish-scales}{1494}
\showcasegameicons{fish-smoking}{1495}
\showcasegameicons{fishbone}{1496}
\showcasegameicons{fishhook-fork}{1497}
\showcasegameicons{fishing-boat}{1498}
\showcasegameicons{fishing-hook}{1499}
\showcasegameicons{fishing-jig}{1500}
\showcasegameicons{fishing-lure}{1501}
\showcasegameicons{fishing-net}{1502}
\showcasegameicons{fishing-pole}{1503}
\showcasegameicons{fishing-spoon}{1504}
\showcasegameicons{fishing}{1505}
\showcasegameicons{fission}{1506}
\showcasegameicons{fist}{1507}
\showcasegameicons{fizzing-flask}{1508}
\showcasegameicons{flag-objective}{1509}
\showcasegameicons{flail}{1510}
\showcasegameicons{flake}{1511}
\showcasegameicons{flame-claws}{1512}
\showcasegameicons{flame-spin}{1513}
\showcasegameicons{flame-tunnel}{1514}
\showcasegameicons{flame}{1515}
\showcasegameicons{flamed-leaf}{1516}
\showcasegameicons{flamer}{1517}
\showcasegameicons{flamethrower-soldier}{1518}
\showcasegameicons{flamethrower}{1519}
\showcasegameicons{flaming-arrow}{1520}
\showcasegameicons{flaming-claw}{1521}
\showcasegameicons{flaming-sheet}{1522}
\showcasegameicons{flaming-trident}{1523}
\showcasegameicons{flamingo}{1524}
\showcasegameicons{flanged-mace}{1525}
\showcasegameicons{flash-grenade}{1526}
\showcasegameicons{flashlight}{1527}
\showcasegameicons{flat-hammer}{1528}
\showcasegameicons{flat-paw-print}{1529}
\showcasegameicons{flat-platform}{1530}
\showcasegameicons{flat-star}{1531}
\showcasegameicons{flat-tire}{1532}
\showcasegameicons{flatbed-covered}{1533}
\showcasegameicons{flatbed}{1534}
\showcasegameicons{flatfish}{1535}
\showcasegameicons{flax}{1536}
\showcasegameicons{fleshy-mass}{1537}
\showcasegameicons{fleur-de-lys}{1538}
\showcasegameicons{flexible-lamp}{1539}
\showcasegameicons{flexible-star}{1540}
\showcasegameicons{flint-spark}{1541}
\showcasegameicons{flip-flops}{1542}
\showcasegameicons{floating-crystal}{1543}
\showcasegameicons{floating-ghost}{1544}
\showcasegameicons{floating-platforms}{1545}
\showcasegameicons{floating-tentacles}{1546}
\showcasegameicons{flood}{1547}
\showcasegameicons{floor-hatch}{1548}
\showcasegameicons{floor-polisher}{1549}
\showcasegameicons{flour}{1550}
\showcasegameicons{flower-emblem}{1551}
\showcasegameicons{flower-hat}{1552}
\showcasegameicons{flower-pot}{1553}
\showcasegameicons{flower-star}{1554}
\showcasegameicons{flower-twirl}{1555}
\showcasegameicons{flowers}{1556}
\showcasegameicons{fluffy-cloud}{1557}
\showcasegameicons{fluffy-flame}{1558}
\showcasegameicons{fluffy-swirl}{1559}
\showcasegameicons{fluffy-trefoil}{1560}
\showcasegameicons{fluffy-wing}{1561}
\showcasegameicons{flute}{1562}
\showcasegameicons{fly}{1563}
\showcasegameicons{flying-beetle}{1564}
\showcasegameicons{flying-dagger}{1565}
\showcasegameicons{flying-flag}{1566}
\showcasegameicons{flying-fox}{1567}
\showcasegameicons{flying-shuriken}{1568}
\showcasegameicons{flying-target}{1569}
\showcasegameicons{flying-trout}{1570}
\showcasegameicons{fn-fal}{1571}
\showcasegameicons{foam}{1572}
\showcasegameicons{foamy-disc}{1573}
\showcasegameicons{focused-lightning}{1574}
\showcasegameicons{fog-light}{1575}
\showcasegameicons{fog}{1576}
\showcasegameicons{folded-paper}{1577}
\showcasegameicons{fomorian}{1578}
\showcasegameicons{food-chain}{1579}
\showcasegameicons{food-truck}{1580}
\showcasegameicons{foot-plaster}{1581}
\showcasegameicons{foot-trip}{1582}
\showcasegameicons{footprint}{1583}
\showcasegameicons{footsteps}{1584}
\showcasegameicons{footy-field}{1585}
\showcasegameicons{forearm}{1586}
\showcasegameicons{forest-camp}{1587}
\showcasegameicons{forest-entrance}{1588}
\showcasegameicons{forest}{1589}
\showcasegameicons{fork-knife-spoon}{1590}
\showcasegameicons{forklift}{1591}
\showcasegameicons{forward-field}{1592}
\showcasegameicons{forward-sun}{1593}
\showcasegameicons{fossil}{1594}
\showcasegameicons{foundry-bucket}{1595}
\showcasegameicons{fountain-pen}{1596}
\showcasegameicons{fountain}{1597}
\showcasegameicons{fox-head}{1598}
\showcasegameicons{fox-tail}{1599}
\showcasegameicons{fox}{1600}
\showcasegameicons{fragmented-meteor}{1601}
\showcasegameicons{fragmented-sword}{1602}
\showcasegameicons{fragrance}{1603}
\showcasegameicons{france}{1604}
\showcasegameicons{frankenstein-creature}{1605}
\showcasegameicons{frayed-arrow}{1606}
\showcasegameicons{freedom-dove}{1607}
\showcasegameicons{freemasonry}{1608}
\showcasegameicons{french-fries}{1609}
\showcasegameicons{french-horn}{1610}
\showcasegameicons{fridge}{1611}
\showcasegameicons{fried-eggs}{1612}
\showcasegameicons{fried-fish}{1613}
\showcasegameicons{frisbee}{1614}
\showcasegameicons{froe-and-mallet}{1615}
\showcasegameicons{frog-foot}{1616}
\showcasegameicons{frog-mouth-helm}{1617}
\showcasegameicons{frog-prince}{1618}
\showcasegameicons{frog}{1619}
\showcasegameicons{front-teeth}{1620}
\showcasegameicons{frontal-lobe}{1621}
\showcasegameicons{frostfire}{1622}
\showcasegameicons{frozen-arrow}{1623}
\showcasegameicons{frozen-block}{1624}
\showcasegameicons{frozen-body}{1625}
\showcasegameicons{frozen-orb}{1626}
\showcasegameicons{frozen-ring}{1627}
\showcasegameicons{fruit-bowl}{1628}
\showcasegameicons{fruit-tree}{1629}
\showcasegameicons{fruiting}{1630}
\showcasegameicons{fuel-tank}{1631}
\showcasegameicons{fuji}{1632}
\showcasegameicons{fulguro-punch}{1633}
\showcasegameicons{full-folder}{1634}
\showcasegameicons{full-metal-bucket-handle}{1635}
\showcasegameicons{full-metal-bucket}{1636}
\showcasegameicons{full-motorcycle-helmet}{1637}
\showcasegameicons{full-pizza}{1638}
\showcasegameicons{full-wood-bucket-handle}{1639}
\showcasegameicons{full-wood-bucket}{1640}
\showcasegameicons{funnel}{1641}
\showcasegameicons{fur-boot}{1642}
\showcasegameicons{fur-shirt}{1643}
\showcasegameicons{furnace}{1644}
\showcasegameicons{g-clef}{1645}
\showcasegameicons{galaxy}{1646}
\showcasegameicons{galea}{1647}
\showcasegameicons{galleon}{1648}
\showcasegameicons{galley}{1649}
\showcasegameicons{game-console}{1650}
\showcasegameicons{gamepad-cross}{1651}
\showcasegameicons{gamepad}{1652}
\showcasegameicons{gardening-shears}{1653}
\showcasegameicons{gargoyle}{1654}
\showcasegameicons{garlic}{1655}
\showcasegameicons{gas-mask}{1656}
\showcasegameicons{gas-pump}{1657}
\showcasegameicons{gas-stove}{1658}
\showcasegameicons{gate}{1659}
\showcasegameicons{gauls-helm}{1660}
\showcasegameicons{gauntlet}{1661}
\showcasegameicons{gavel}{1662}
\showcasegameicons{gaze}{1663}
\showcasegameicons{gear-hammer}{1664}
\showcasegameicons{gear-stick-pattern}{1665}
\showcasegameicons{gear-stick}{1666}
\showcasegameicons{gears}{1667}
\showcasegameicons{gecko}{1668}
\showcasegameicons{gem-chain}{1669}
\showcasegameicons{gem-necklace}{1670}
\showcasegameicons{gem-pendant}{1671}
\showcasegameicons{gemini}{1672}
\showcasegameicons{gems}{1673}
\showcasegameicons{ghost-ally}{1674}
\showcasegameicons{ghost}{1675}
\showcasegameicons{giant-squid}{1676}
\showcasegameicons{giant}{1677}
\showcasegameicons{gibbet}{1678}
\showcasegameicons{gift-of-knowledge}{1679}
\showcasegameicons{gift-trap}{1680}
\showcasegameicons{gingerbread-man}{1681}
\showcasegameicons{ginkgo-leaf}{1682}
\showcasegameicons{gladius}{1683}
\showcasegameicons{glaive}{1684}
\showcasegameicons{glass-ball}{1685}
\showcasegameicons{glass-celebration}{1686}
\showcasegameicons{glass-heart}{1687}
\showcasegameicons{glass-shot}{1688}
\showcasegameicons{glider}{1689}
\showcasegameicons{globe-ring}{1690}
\showcasegameicons{globe}{1691}
\showcasegameicons{glock}{1692}
\showcasegameicons{gloop}{1693}
\showcasegameicons{gloves}{1694}
\showcasegameicons{glowing-artifact}{1695}
\showcasegameicons{glowing-hands}{1696}
\showcasegameicons{gluttonous-smile}{1697}
\showcasegameicons{gluttony}{1698}
\showcasegameicons{goal-keeper}{1699}
\showcasegameicons{goat}{1700}
\showcasegameicons{goblin-camp}{1701}
\showcasegameicons{goblin-head}{1702}
\showcasegameicons{goblin}{1703}
\showcasegameicons{gold-bar}{1704}
\showcasegameicons{gold-mine}{1705}
\showcasegameicons{gold-nuggets}{1706}
\showcasegameicons{gold-scarab}{1707}
\showcasegameicons{gold-shell}{1708}
\showcasegameicons{gold-stack}{1709}
\showcasegameicons{golem-head}{1710}
\showcasegameicons{golf-flag}{1711}
\showcasegameicons{golf-tee}{1712}
\showcasegameicons{gong}{1713}
\showcasegameicons{goo-explosion}{1714}
\showcasegameicons{goo-skull}{1715}
\showcasegameicons{goo-spurt}{1716}
\showcasegameicons{gooey-daemon}{1717}
\showcasegameicons{gooey-eyed-sun}{1718}
\showcasegameicons{gooey-impact}{1719}
\showcasegameicons{gooey-molecule}{1720}
\showcasegameicons{gooey-sword}{1721}
\showcasegameicons{goose}{1722}
\showcasegameicons{gorilla}{1723}
\showcasegameicons{gothic-cross}{1724}
\showcasegameicons{gps}{1725}
\showcasegameicons{grab}{1726}
\showcasegameicons{graduate-cap}{1727}
\showcasegameicons{grain-bundle}{1728}
\showcasegameicons{grain}{1729}
\showcasegameicons{granary}{1730}
\showcasegameicons{grand-piano}{1731}
\showcasegameicons{grapes}{1732}
\showcasegameicons{grapple}{1733}
\showcasegameicons{grasping-claws}{1734}
\showcasegameicons{grasping-slug}{1735}
\showcasegameicons{grass-mushroom}{1736}
\showcasegameicons{grass}{1737}
\showcasegameicons{grave-flowers}{1738}
\showcasegameicons{graveyard}{1739}
\showcasegameicons{gravitation}{1740}
\showcasegameicons{grease-trap}{1741}
\showcasegameicons{great-pyramid}{1742}
\showcasegameicons{great-war-tank}{1743}
\showcasegameicons{greaves}{1744}
\showcasegameicons{greek-sphinx}{1745}
\showcasegameicons{greek-temple}{1746}
\showcasegameicons{green-power}{1747}
\showcasegameicons{greenhouse}{1748}
\showcasegameicons{grenade}{1749}
\showcasegameicons{griffin-shield}{1750}
\showcasegameicons{griffin-symbol}{1751}
\showcasegameicons{grim-reaper}{1752}
\showcasegameicons{ground-sprout}{1753}
\showcasegameicons{groundbreaker}{1754}
\showcasegameicons{grouped-drops}{1755}
\showcasegameicons{growth}{1756}
\showcasegameicons{guarded-tower}{1757}
\showcasegameicons{guards}{1758}
\showcasegameicons{guatemala}{1759}
\showcasegameicons{guillotine}{1760}
\showcasegameicons{guitar-bass-head}{1761}
\showcasegameicons{guitar-head}{1762}
\showcasegameicons{guitar}{1763}
\showcasegameicons{gun-rose}{1764}
\showcasegameicons{gun-stock}{1765}
\showcasegameicons{gunshot}{1766}
\showcasegameicons{gym-bag}{1767}
\showcasegameicons{gyroscope}{1768}
\showcasegameicons{h2o}{1769}
\showcasegameicons{habitat-dome}{1770}
\showcasegameicons{hades-symbol}{1771}
\showcasegameicons{hair-strands}{1772}
\showcasegameicons{halberd-shuriken}{1773}
\showcasegameicons{halberd}{1774}
\showcasegameicons{half-body-crawling}{1775}
\showcasegameicons{half-dead}{1776}
\showcasegameicons{half-heart}{1777}
\showcasegameicons{half-log}{1778}
\showcasegameicons{half-tornado}{1779}
\showcasegameicons{halt}{1780}
\showcasegameicons{ham-shank}{1781}
\showcasegameicons{hamburger-menu}{1782}
\showcasegameicons{hamburger}{1783}
\showcasegameicons{hammer-break}{1784}
\showcasegameicons{hammer-drop}{1785}
\showcasegameicons{hammer-nails}{1786}
\showcasegameicons{hammer-sickle}{1787}
\showcasegameicons{hand-bag}{1788}
\showcasegameicons{hand-bandage}{1789}
\showcasegameicons{hand-grip}{1790}
\showcasegameicons{hand-of-god}{1791}
\showcasegameicons{hand-ok}{1792}
\showcasegameicons{hand-saw}{1793}
\showcasegameicons{hand-truck}{1794}
\showcasegameicons{hand-wing}{1795}
\showcasegameicons{hand}{1796}
\showcasegameicons{handcuffed}{1797}
\showcasegameicons{handcuffs}{1798}
\showcasegameicons{handheld-fan}{1799}
\showcasegameicons{hang-glider}{1800}
\showcasegameicons{hanger}{1801}
\showcasegameicons{hanging-sign}{1802}
\showcasegameicons{hanging-spider}{1803}
\showcasegameicons{happy-skull}{1804}
\showcasegameicons{harbor-dock}{1805}
\showcasegameicons{harp}{1806}
\showcasegameicons{harpoon-chain}{1807}
\showcasegameicons{harpoon-trident}{1808}
\showcasegameicons{harpy}{1809}
\showcasegameicons{harry-potter-skull}{1810}
\showcasegameicons{hasty-grave}{1811}
\showcasegameicons{hatchet}{1812}
\showcasegameicons{hatchets}{1813}
\showcasegameicons{haunting}{1814}
\showcasegameicons{hawk-emblem}{1815}
\showcasegameicons{hazard-sign}{1816}
\showcasegameicons{hazmat-suit}{1817}
\showcasegameicons{head-shot}{1818}
\showcasegameicons{headband-knot}{1819}
\showcasegameicons{headphones}{1820}
\showcasegameicons{headshot}{1821}
\showcasegameicons{healing-shield}{1822}
\showcasegameicons{healing}{1823}
\showcasegameicons{health-capsule}{1824}
\showcasegameicons{health-decrease}{1825}
\showcasegameicons{health-increase}{1826}
\showcasegameicons{health-normal}{1827}
\showcasegameicons{health-potion}{1828}
\showcasegameicons{hearing-disabled}{1829}
\showcasegameicons{heart-armor}{1830}
\showcasegameicons{heart-battery}{1831}
\showcasegameicons{heart-beats}{1832}
\showcasegameicons{heart-bottle}{1833}
\showcasegameicons{heart-drop}{1834}
\showcasegameicons{heart-earrings}{1835}
\showcasegameicons{heart-inside}{1836}
\showcasegameicons{heart-key}{1837}
\showcasegameicons{heart-minus}{1838}
\showcasegameicons{heart-necklace}{1839}
\showcasegameicons{heart-organ}{1840}
\showcasegameicons{heart-plus}{1841}
\showcasegameicons{heart-shield}{1842}
\showcasegameicons{heart-stake}{1843}
\showcasegameicons{heart-tower}{1844}
\showcasegameicons{heart-wings}{1845}
\showcasegameicons{heartburn}{1846}
\showcasegameicons{hearts}{1847}
\showcasegameicons{heat-haze}{1848}
\showcasegameicons{heaven-gate}{1849}
\showcasegameicons{heavy-arrow}{1850}
\showcasegameicons{heavy-bullets}{1851}
\showcasegameicons{heavy-collar}{1852}
\showcasegameicons{heavy-fall}{1853}
\showcasegameicons{heavy-fighter}{1854}
\showcasegameicons{heavy-helm}{1855}
\showcasegameicons{heavy-lightning}{1856}
\showcasegameicons{heavy-rain}{1857}
\showcasegameicons{heavy-thorny-triskelion}{1858}
\showcasegameicons{heavy-timer}{1859}
\showcasegameicons{hedgehog}{1860}
\showcasegameicons{hedjet-white-crown}{1861}
\showcasegameicons{helicoprion}{1862}
\showcasegameicons{helicopter-tail}{1863}
\showcasegameicons{helicopter}{1864}
\showcasegameicons{hell-crosses}{1865}
\showcasegameicons{helmet-head-shot}{1866}
\showcasegameicons{helmet}{1867}
\showcasegameicons{help}{1868}
\showcasegameicons{hemp}{1869}
\showcasegameicons{heptagram}{1870}
\showcasegameicons{heraldic-sun}{1871}
\showcasegameicons{herbs-bundle}{1872}
\showcasegameicons{heron}{1873}
\showcasegameicons{hexagonal-nut}{1874}
\showcasegameicons{hexes}{1875}
\showcasegameicons{hidden}{1876}
\showcasegameicons{hieroglyph-legs}{1877}
\showcasegameicons{hieroglyph-y}{1878}
\showcasegameicons{high-five}{1879}
\showcasegameicons{high-grass}{1880}
\showcasegameicons{high-heel}{1881}
\showcasegameicons{high-kick}{1882}
\showcasegameicons{high-punch}{1883}
\showcasegameicons{high-shot}{1884}
\showcasegameicons{high-tide}{1885}
\showcasegameicons{highlighter}{1886}
\showcasegameicons{hiking}{1887}
\showcasegameicons{hill-conquest}{1888}
\showcasegameicons{hill-fort}{1889}
\showcasegameicons{hills}{1890}
\showcasegameicons{histogram}{1891}
\showcasegameicons{hive-mind}{1892}
\showcasegameicons{hive}{1893}
\showcasegameicons{hobbit-door}{1894}
\showcasegameicons{hobbit-dwelling}{1895}
\showcasegameicons{hockey}{1896}
\showcasegameicons{hole-ladder}{1897}
\showcasegameicons{hole}{1898}
\showcasegameicons{hollow-cat}{1899}
\showcasegameicons{holosphere}{1900}
\showcasegameicons{holy-grail}{1901}
\showcasegameicons{holy-hand-grenade}{1902}
\showcasegameicons{holy-oak}{1903}
\showcasegameicons{holy-symbol}{1904}
\showcasegameicons{holy-water}{1905}
\showcasegameicons{home-garage}{1906}
\showcasegameicons{honey-jar}{1907}
\showcasegameicons{honeycomb}{1908}
\showcasegameicons{honeypot}{1909}
\showcasegameicons{hood}{1910}
\showcasegameicons{hooded-assassin}{1911}
\showcasegameicons{hooded-figure}{1912}
\showcasegameicons{hoodie}{1913}
\showcasegameicons{hoof}{1914}
\showcasegameicons{hook-swords}{1915}
\showcasegameicons{hook}{1916}
\showcasegameicons{hops}{1917}
\showcasegameicons{horizon-road}{1918}
\showcasegameicons{horizontal-flip}{1919}
\showcasegameicons{horn-internal}{1920}
\showcasegameicons{horned-helm}{1921}
\showcasegameicons{horned-reptile}{1922}
\showcasegameicons{horned-skull}{1923}
\showcasegameicons{horse-head}{1924}
\showcasegameicons{horseshoe}{1925}
\showcasegameicons{horus}{1926}
\showcasegameicons{hospital-cross}{1927}
\showcasegameicons{hospital}{1928}
\showcasegameicons{hot-dog}{1929}
\showcasegameicons{hot-meal}{1930}
\showcasegameicons{hot-spices}{1931}
\showcasegameicons{hot-surface}{1932}
\showcasegameicons{hound}{1933}
\showcasegameicons{hourglass}{1934}
\showcasegameicons{house-keys}{1935}
\showcasegameicons{house}{1936}
\showcasegameicons{human-cannonball}{1937}
\showcasegameicons{human-ear}{1938}
\showcasegameicons{human-pyramid}{1939}
\showcasegameicons{human-target}{1940}
\showcasegameicons{hummingbird}{1941}
\showcasegameicons{hungary}{1942}
\showcasegameicons{hunter-eyes}{1943}
\showcasegameicons{hunting-bolas}{1944}
\showcasegameicons{hunting-horn}{1945}
\showcasegameicons{hut}{1946}
\showcasegameicons{huts-village}{1947}
\showcasegameicons{hydra-shot}{1948}
\showcasegameicons{hydra}{1949}
\showcasegameicons{hyena-head}{1950}
\showcasegameicons{hypersonic-bolt}{1951}
\showcasegameicons{hypersonic-melon}{1952}
\showcasegameicons{hypodermic-test}{1953}
\showcasegameicons{i-beam}{1954}
\showcasegameicons{i-brick}{1955}
\showcasegameicons{ibis}{1956}
\showcasegameicons{icarus}{1957}
\showcasegameicons{ice-bolt}{1958}
\showcasegameicons{ice-bomb}{1959}
\showcasegameicons{ice-cream-cone}{1960}
\showcasegameicons{ice-cream-scoop}{1961}
\showcasegameicons{ice-cube}{1962}
\showcasegameicons{ice-cubes}{1963}
\showcasegameicons{ice-golem}{1964}
\showcasegameicons{ice-iris}{1965}
\showcasegameicons{ice-pop}{1966}
\showcasegameicons{ice-shield}{1967}
\showcasegameicons{ice-skate}{1968}
\showcasegameicons{ice-spear}{1969}
\showcasegameicons{ice-spell-cast}{1970}
\showcasegameicons{iceberg}{1971}
\showcasegameicons{icebergs}{1972}
\showcasegameicons{iceland}{1973}
\showcasegameicons{icicles-aura}{1974}
\showcasegameicons{icicles-fence}{1975}
\showcasegameicons{id-card}{1976}
\showcasegameicons{idea}{1977}
\showcasegameicons{ifrit}{1978}
\showcasegameicons{igloo}{1979}
\showcasegameicons{imbricated-arrows}{1980}
\showcasegameicons{imp-laugh}{1981}
\showcasegameicons{imp}{1982}
\showcasegameicons{impact-point}{1983}
\showcasegameicons{imperial-crown}{1984}
\showcasegameicons{implosion}{1985}
\showcasegameicons{imprisoned}{1986}
\showcasegameicons{inauguration}{1987}
\showcasegameicons{incense}{1988}
\showcasegameicons{incisors}{1989}
\showcasegameicons{incoming-rocket}{1990}
\showcasegameicons{incubator}{1991}
\showcasegameicons{india-gate}{1992}
\showcasegameicons{indian-palace}{1993}
\showcasegameicons{inferno-bomb}{1994}
\showcasegameicons{infested-mass}{1995}
\showcasegameicons{infinity}{1996}
\showcasegameicons{info}{1997}
\showcasegameicons{injustice}{1998}
\showcasegameicons{ink-swirl}{1999}
\showcasegameicons{inner-self}{2000}
\showcasegameicons{insect-jaws}{2001}
\showcasegameicons{inspiration}{2002}
\showcasegameicons{interceptor-ship}{2003}
\showcasegameicons{interdiction}{2004}
\showcasegameicons{interlaced-tentacles}{2005}
\showcasegameicons{interleaved-arrows}{2006}
\showcasegameicons{interleaved-claws}{2007}
\showcasegameicons{internal-injury}{2008}
\showcasegameicons{internal-organ}{2009}
\showcasegameicons{interstellar-path}{2010}
\showcasegameicons{intricate-necklace}{2011}
\showcasegameicons{inverted-dice-1}{2012}
\showcasegameicons{inverted-dice-2}{2013}
\showcasegameicons{inverted-dice-3}{2014}
\showcasegameicons{inverted-dice-4}{2015}
\showcasegameicons{inverted-dice-5}{2016}
\showcasegameicons{inverted-dice-6}{2017}
\showcasegameicons{invisible-face}{2018}
\showcasegameicons{invisible}{2019}
\showcasegameicons{ion-cannon-blast}{2020}
\showcasegameicons{ionic-column}{2021}
\showcasegameicons{iraq}{2022}
\showcasegameicons{iron-cross}{2023}
\showcasegameicons{iron-hulled-warship}{2024}
\showcasegameicons{iron-mask}{2025}
\showcasegameicons{island}{2026}
\showcasegameicons{italia}{2027}
\showcasegameicons{ivory-tusks}{2028}
\showcasegameicons{j-brick}{2029}
\showcasegameicons{jack-plug}{2030}
\showcasegameicons{james-bond-aperture}{2031}
\showcasegameicons{japan}{2032}
\showcasegameicons{japanese-bridge}{2033}
\showcasegameicons{jasmine}{2034}
\showcasegameicons{jason-mask}{2035}
\showcasegameicons{jawbone}{2036}
\showcasegameicons{jawless-cyclop}{2037}
\showcasegameicons{jeep}{2038}
\showcasegameicons{jelly-beans}{2039}
\showcasegameicons{jelly}{2040}
\showcasegameicons{jellyfish}{2041}
\showcasegameicons{jerrycan}{2042}
\showcasegameicons{jerusalem-cross}{2043}
\showcasegameicons{jester-hat}{2044}
\showcasegameicons{jet-fighter}{2045}
\showcasegameicons{jet-pack}{2046}
\showcasegameicons{jetpack}{2047}
\showcasegameicons{jewel-crown}{2048}
\showcasegameicons{jeweled-chalice}{2049}
\showcasegameicons{jigsaw-box}{2050}
\showcasegameicons{jigsaw-piece}{2051}
\showcasegameicons{join}{2052}
\showcasegameicons{joint}{2053}
\showcasegameicons{journey}{2054}
\showcasegameicons{joystick}{2055}
\showcasegameicons{jug}{2056}
\showcasegameicons{juggler}{2057}
\showcasegameicons{juggling-clubs}{2058}
\showcasegameicons{juggling-seal}{2059}
\showcasegameicons{jump-across}{2060}
\showcasegameicons{jumping-dog}{2061}
\showcasegameicons{jumping-rope}{2062}
\showcasegameicons{jungle}{2063}
\showcasegameicons{jupiter}{2064}
\showcasegameicons{justice-star}{2065}
\showcasegameicons{kaleidoscope-pearls}{2066}
\showcasegameicons{kangaroo}{2067}
\showcasegameicons{katana}{2068}
\showcasegameicons{kebab-spit}{2069}
\showcasegameicons{kenku-head}{2070}
\showcasegameicons{kenya}{2071}
\showcasegameicons{ketchup}{2072}
\showcasegameicons{kevlar-vest}{2073}
\showcasegameicons{kevlar}{2074}
\showcasegameicons{key-card}{2075}
\showcasegameicons{key-lock}{2076}
\showcasegameicons{key}{2077}
\showcasegameicons{keyboard}{2078}
\showcasegameicons{keyring}{2079}
\showcasegameicons{kick-scooter}{2080}
\showcasegameicons{kid-slide}{2081}
\showcasegameicons{kidneys}{2082}
\showcasegameicons{kimono}{2083}
\showcasegameicons{kindle}{2084}
\showcasegameicons{king-ju-mask}{2085}
\showcasegameicons{king}{2086}
\showcasegameicons{kitchen-knives}{2087}
\showcasegameicons{kitchen-scale}{2088}
\showcasegameicons{kitchen-tap}{2089}
\showcasegameicons{kite}{2090}
\showcasegameicons{kiwi-bird}{2091}
\showcasegameicons{kiwi-fruit}{2092}
\showcasegameicons{klingon}{2093}
\showcasegameicons{knapsack}{2094}
\showcasegameicons{knee-bandage}{2095}
\showcasegameicons{knee-cap}{2096}
\showcasegameicons{knee-pad}{2097}
\showcasegameicons{kneeling}{2098}
\showcasegameicons{knife-fork}{2099}
\showcasegameicons{knife-thrust}{2100}
\showcasegameicons{knight-banner}{2101}
\showcasegameicons{knocked-out-stars}{2102}
\showcasegameicons{knockout}{2103}
\showcasegameicons{knot}{2104}
\showcasegameicons{koala}{2105}
\showcasegameicons{koholint-egg}{2106}
\showcasegameicons{kraken-tentacle}{2107}
\showcasegameicons{kusarigama}{2108}
\showcasegameicons{l-brick}{2109}
\showcasegameicons{lab-coat}{2110}
\showcasegameicons{labrador-head}{2111}
\showcasegameicons{ladder}{2112}
\showcasegameicons{ladders-platform}{2113}
\showcasegameicons{ladle}{2114}
\showcasegameicons{ladybug}{2115}
\showcasegameicons{lambda}{2116}
\showcasegameicons{lamellar}{2117}
\showcasegameicons{lamprey-mouth}{2118}
\showcasegameicons{land-mine}{2119}
\showcasegameicons{lantern-flame}{2120}
\showcasegameicons{lantern}{2121}
\showcasegameicons{laptop}{2122}
\showcasegameicons{large-dress}{2123}
\showcasegameicons{large-paint-brush}{2124}
\showcasegameicons{large-wound}{2125}
\showcasegameicons{laser-blast}{2126}
\showcasegameicons{laser-burst}{2127}
\showcasegameicons{laser-gun}{2128}
\showcasegameicons{laser-precision}{2129}
\showcasegameicons{laser-sparks}{2130}
\showcasegameicons{laser-turret}{2131}
\showcasegameicons{laser-warning}{2132}
\showcasegameicons{laserburn}{2133}
\showcasegameicons{lasso}{2134}
\showcasegameicons{latvia}{2135}
\showcasegameicons{laurel-crown}{2136}
\showcasegameicons{laurels-trophy}{2137}
\showcasegameicons{laurels}{2138}
\showcasegameicons{lava}{2139}
\showcasegameicons{law-star}{2140}
\showcasegameicons{layered-armor}{2141}
\showcasegameicons{lead-pipe}{2142}
\showcasegameicons{leaf-skeleton}{2143}
\showcasegameicons{leaf-swirl}{2144}
\showcasegameicons{leak}{2145}
\showcasegameicons{leaky-skull}{2146}
\showcasegameicons{leapfrog}{2147}
\showcasegameicons{leather-armor}{2148}
\showcasegameicons{leather-boot}{2149}
\showcasegameicons{leather-vest}{2150}
\showcasegameicons{led}{2151}
\showcasegameicons{lee-enfield}{2152}
\showcasegameicons{leeching-worm}{2153}
\showcasegameicons{leek}{2154}
\showcasegameicons{leg-armor}{2155}
\showcasegameicons{leg}{2156}
\showcasegameicons{lemon}{2157}
\showcasegameicons{leo}{2158}
\showcasegameicons{letter-bomb}{2159}
\showcasegameicons{level-crossing}{2160}
\showcasegameicons{level-end-flag}{2161}
\showcasegameicons{level-four-advanced}{2162}
\showcasegameicons{level-four}{2163}
\showcasegameicons{level-three-advanced}{2164}
\showcasegameicons{level-three}{2165}
\showcasegameicons{level-two-advanced}{2166}
\showcasegameicons{level-two}{2167}
\showcasegameicons{lever}{2168}
\showcasegameicons{liar}{2169}
\showcasegameicons{liberty-wing}{2170}
\showcasegameicons{libra}{2171}
\showcasegameicons{libya}{2172}
\showcasegameicons{life-bar}{2173}
\showcasegameicons{life-buoy}{2174}
\showcasegameicons{life-in-the-balance}{2175}
\showcasegameicons{life-jacket}{2176}
\showcasegameicons{life-support}{2177}
\showcasegameicons{life-tap}{2178}
\showcasegameicons{lift}{2179}
\showcasegameicons{light-backpack}{2180}
\showcasegameicons{light-bulb}{2181}
\showcasegameicons{light-fighter}{2182}
\showcasegameicons{light-helm}{2183}
\showcasegameicons{light-projector}{2184}
\showcasegameicons{light-sabers}{2185}
\showcasegameicons{light-thorny-triskelion}{2186}
\showcasegameicons{lighter}{2187}
\showcasegameicons{lighthouse}{2188}
\showcasegameicons{lightning-arc}{2189}
\showcasegameicons{lightning-bow}{2190}
\showcasegameicons{lightning-branches}{2191}
\showcasegameicons{lightning-dissipation}{2192}
\showcasegameicons{lightning-dome}{2193}
\showcasegameicons{lightning-electron}{2194}
\showcasegameicons{lightning-flame}{2195}
\showcasegameicons{lightning-frequency}{2196}
\showcasegameicons{lightning-helix}{2197}
\showcasegameicons{lightning-mask}{2198}
\showcasegameicons{lightning-saber}{2199}
\showcasegameicons{lightning-shadow}{2200}
\showcasegameicons{lightning-shield}{2201}
\showcasegameicons{lightning-shout}{2202}
\showcasegameicons{lightning-slashes}{2203}
\showcasegameicons{lightning-spanner}{2204}
\showcasegameicons{lightning-storm}{2205}
\showcasegameicons{lightning-tear}{2206}
\showcasegameicons{lightning-tree}{2207}
\showcasegameicons{lightning-trio}{2208}
\showcasegameicons{lily-pads}{2209}
\showcasegameicons{linden-leaf}{2210}
\showcasegameicons{linked-rings}{2211}
\showcasegameicons{lion}{2212}
\showcasegameicons{lips}{2213}
\showcasegameicons{lipstick}{2214}
\showcasegameicons{liquid-soap}{2215}
\showcasegameicons{lit-candelabra}{2216}
\showcasegameicons{liver}{2217}
\showcasegameicons{lizard-tongue}{2218}
\showcasegameicons{lizardman}{2219}
\showcasegameicons{load}{2220}
\showcasegameicons{lob-arrow}{2221}
\showcasegameicons{lock-picking}{2222}
\showcasegameicons{lock-spy}{2223}
\showcasegameicons{locked-box}{2224}
\showcasegameicons{locked-chest}{2225}
\showcasegameicons{locked-door}{2226}
\showcasegameicons{locked-fortress}{2227}
\showcasegameicons{locked-heart}{2228}
\showcasegameicons{lockers}{2229}
\showcasegameicons{lockpicks}{2230}
\showcasegameicons{log}{2231}
\showcasegameicons{logging}{2232}
\showcasegameicons{logic-gate-and}{2233}
\showcasegameicons{logic-gate-nand}{2234}
\showcasegameicons{logic-gate-nor}{2235}
\showcasegameicons{logic-gate-not}{2236}
\showcasegameicons{logic-gate-nxor}{2237}
\showcasegameicons{logic-gate-or}{2238}
\showcasegameicons{logic-gate-xor}{2239}
\showcasegameicons{loincloth}{2240}
\showcasegameicons{long-antennae-bug}{2241}
\showcasegameicons{long-legged-spider}{2242}
\showcasegameicons{look-at}{2243}
\showcasegameicons{lorgnette}{2244}
\showcasegameicons{lost-limb}{2245}
\showcasegameicons{lotus-flower}{2246}
\showcasegameicons{lotus}{2247}
\showcasegameicons{louvre-pyramid}{2248}
\showcasegameicons{love-howl}{2249}
\showcasegameicons{love-injection}{2250}
\showcasegameicons{love-letter}{2251}
\showcasegameicons{love-mystery}{2252}
\showcasegameicons{love-song}{2253}
\showcasegameicons{lovers}{2254}
\showcasegameicons{low-tide}{2255}
\showcasegameicons{luchador}{2256}
\showcasegameicons{lucifer-cannon}{2257}
\showcasegameicons{lucky-fisherman}{2258}
\showcasegameicons{luger}{2259}
\showcasegameicons{lunar-module}{2260}
\showcasegameicons{lunar-wand}{2261}
\showcasegameicons{lungs}{2262}
\showcasegameicons{lynx-head}{2263}
\showcasegameicons{lyre}{2264}
\showcasegameicons{m3-grease-gun}{2265}
\showcasegameicons{mac-10}{2266}
\showcasegameicons{mace-head}{2267}
\showcasegameicons{machete}{2268}
\showcasegameicons{machine-gun-magazine}{2269}
\showcasegameicons{machine-gun}{2270}
\showcasegameicons{mad-scientist}{2271}
\showcasegameicons{maggot}{2272}
\showcasegameicons{magic-axe}{2273}
\showcasegameicons{magic-broom}{2274}
\showcasegameicons{magic-gate}{2275}
\showcasegameicons{magic-hat}{2276}
\showcasegameicons{magic-lamp}{2277}
\showcasegameicons{magic-palm}{2278}
\showcasegameicons{magic-portal}{2279}
\showcasegameicons{magic-potion}{2280}
\showcasegameicons{magic-shield}{2281}
\showcasegameicons{magic-swirl}{2282}
\showcasegameicons{magic-trident}{2283}
\showcasegameicons{magick-trick}{2284}
\showcasegameicons{magnet-blast}{2285}
\showcasegameicons{magnet-man}{2286}
\showcasegameicons{magnet}{2287}
\showcasegameicons{magnifying-glass}{2288}
\showcasegameicons{mail-shirt}{2289}
\showcasegameicons{mailbox}{2290}
\showcasegameicons{mailed-fist}{2291}
\showcasegameicons{male}{2292}
\showcasegameicons{mammoth}{2293}
\showcasegameicons{manacles}{2294}
\showcasegameicons{mandrill-head}{2295}
\showcasegameicons{manta-ray}{2296}
\showcasegameicons{mantrap}{2297}
\showcasegameicons{manual-juicer}{2298}
\showcasegameicons{manual-meat-grinder}{2299}
\showcasegameicons{manual-water-pump}{2300}
\showcasegameicons{maple-leaf}{2301}
\showcasegameicons{maracas}{2302}
\showcasegameicons{marble-tap}{2303}
\showcasegameicons{marbles}{2304}
\showcasegameicons{marrow-drain}{2305}
\showcasegameicons{mars-curiosity}{2306}
\showcasegameicons{mars-pathfinder}{2307}
\showcasegameicons{marshmallows}{2308}
\showcasegameicons{martini}{2309}
\showcasegameicons{martyr-memorial}{2310}
\showcasegameicons{masked-spider}{2311}
\showcasegameicons{mason-jar}{2312}
\showcasegameicons{mass-driver}{2313}
\showcasegameicons{master-of-arms}{2314}
\showcasegameicons{match-head}{2315}
\showcasegameicons{match-tip}{2316}
\showcasegameicons{matchbox}{2317}
\showcasegameicons{materials-science}{2318}
\showcasegameicons{matryoshka-dolls}{2319}
\showcasegameicons{matter-states}{2320}
\showcasegameicons{mayan-pyramid}{2321}
\showcasegameicons{maze-cornea}{2322}
\showcasegameicons{maze-saw}{2323}
\showcasegameicons{maze}{2324}
\showcasegameicons{meal}{2325}
\showcasegameicons{measure-tape}{2326}
\showcasegameicons{meat-cleaver}{2327}
\showcasegameicons{meat-hook}{2328}
\showcasegameicons{meat}{2329}
\showcasegameicons{mecha-head}{2330}
\showcasegameicons{mecha-mask}{2331}
\showcasegameicons{mechanic-garage}{2332}
\showcasegameicons{mechanical-arm}{2333}
\showcasegameicons{medal-skull}{2334}
\showcasegameicons{medal}{2335}
\showcasegameicons{medallist}{2336}
\showcasegameicons{medical-drip}{2337}
\showcasegameicons{medical-pack-alt}{2338}
\showcasegameicons{medical-pack}{2339}
\showcasegameicons{medical-thermometer}{2340}
\showcasegameicons{medicine-pills}{2341}
\showcasegameicons{medicines}{2342}
\showcasegameicons{medieval-barracks}{2343}
\showcasegameicons{medieval-gate}{2344}
\showcasegameicons{medieval-pavilion}{2345}
\showcasegameicons{medieval-village-01}{2346}
\showcasegameicons{meditation}{2347}
\showcasegameicons{medusa-head}{2348}
\showcasegameicons{meeple-army}{2349}
\showcasegameicons{meeple-circle}{2350}
\showcasegameicons{meeple-group}{2351}
\showcasegameicons{meeple-king}{2352}
\showcasegameicons{meeple}{2353}
\showcasegameicons{megabot}{2354}
\showcasegameicons{megaphone}{2355}
\showcasegameicons{melting-ice-cube}{2356}
\showcasegameicons{melting-metal}{2357}
\showcasegameicons{menhir}{2358}
\showcasegameicons{mermaid}{2359}
\showcasegameicons{mesh-ball}{2360}
\showcasegameicons{mesh-network}{2361}
\showcasegameicons{metal-bar}{2362}
\showcasegameicons{metal-boot}{2363}
\showcasegameicons{metal-detector}{2364}
\showcasegameicons{metal-disc}{2365}
\showcasegameicons{metal-golem-head}{2366}
\showcasegameicons{metal-hand}{2367}
\showcasegameicons{metal-plate}{2368}
\showcasegameicons{metal-scales}{2369}
\showcasegameicons{metal-skirt}{2370}
\showcasegameicons{meteor-impact}{2371}
\showcasegameicons{metroid}{2372}
\showcasegameicons{metronome}{2373}
\showcasegameicons{mexico}{2374}
\showcasegameicons{microchip}{2375}
\showcasegameicons{microphone}{2376}
\showcasegameicons{microscope-lens}{2377}
\showcasegameicons{microscope}{2378}
\showcasegameicons{mid-summit}{2379}
\showcasegameicons{middle-arrow}{2380}
\showcasegameicons{midnight-claw}{2381}
\showcasegameicons{mighty-boosh}{2382}
\showcasegameicons{mighty-force}{2383}
\showcasegameicons{mighty-horn}{2384}
\showcasegameicons{mighty-spanner}{2385}
\showcasegameicons{military-ambulance}{2386}
\showcasegameicons{military-fort}{2387}
\showcasegameicons{milk-carton}{2388}
\showcasegameicons{millenium-key}{2389}
\showcasegameicons{mimic-chest}{2390}
\showcasegameicons{mine-explosion}{2391}
\showcasegameicons{mine-truck}{2392}
\showcasegameicons{mine-wagon}{2393}
\showcasegameicons{minefield}{2394}
\showcasegameicons{miner}{2395}
\showcasegameicons{mineral-heart}{2396}
\showcasegameicons{mineral-pearls}{2397}
\showcasegameicons{minerals}{2398}
\showcasegameicons{mini-submarine}{2399}
\showcasegameicons{minigun}{2400}
\showcasegameicons{mining-helmet}{2401}
\showcasegameicons{mining}{2402}
\showcasegameicons{minions}{2403}
\showcasegameicons{minotaur}{2404}
\showcasegameicons{miracle-medecine}{2405}
\showcasegameicons{mirror-mirror}{2406}
\showcasegameicons{misdirection}{2407}
\showcasegameicons{missile-launcher}{2408}
\showcasegameicons{missile-mech}{2409}
\showcasegameicons{missile-pod}{2410}
\showcasegameicons{missile-swarm}{2411}
\showcasegameicons{mite-alt}{2412}
\showcasegameicons{mite}{2413}
\showcasegameicons{moai}{2414}
\showcasegameicons{modern-city}{2415}
\showcasegameicons{moebius-star}{2416}
\showcasegameicons{moebius-trefoil}{2417}
\showcasegameicons{moebius-triangle}{2418}
\showcasegameicons{moka-pot}{2419}
\showcasegameicons{moldova}{2420}
\showcasegameicons{mole}{2421}
\showcasegameicons{molecule}{2422}
\showcasegameicons{molotov}{2423}
\showcasegameicons{mona-lisa}{2424}
\showcasegameicons{moncler-jacket}{2425}
\showcasegameicons{money-stack}{2426}
\showcasegameicons{mongolia}{2427}
\showcasegameicons{monk-face}{2428}
\showcasegameicons{monkey-wrench}{2429}
\showcasegameicons{monkey}{2430}
\showcasegameicons{mono-wheel-robot}{2431}
\showcasegameicons{monster-grasp}{2432}
\showcasegameicons{monstera-leaf}{2433}
\showcasegameicons{monument-valley}{2434}
\showcasegameicons{moon-bats}{2435}
\showcasegameicons{moon-claws}{2436}
\showcasegameicons{moon-orbit}{2437}
\showcasegameicons{moon}{2438}
\showcasegameicons{mooring-bollard}{2439}
\showcasegameicons{morbid-humour}{2440}
\showcasegameicons{morgue-feet}{2441}
\showcasegameicons{morph-ball}{2442}
\showcasegameicons{mortar}{2443}
\showcasegameicons{mountain-cave}{2444}
\showcasegameicons{mountain-climbing}{2445}
\showcasegameicons{mountain-road}{2446}
\showcasegameicons{mountains}{2447}
\showcasegameicons{mountaintop}{2448}
\showcasegameicons{mounted-knight}{2449}
\showcasegameicons{mouse}{2450}
\showcasegameicons{mouth-watering}{2451}
\showcasegameicons{move}{2452}
\showcasegameicons{movement-sensor}{2453}
\showcasegameicons{mp-40}{2454}
\showcasegameicons{mp5}{2455}
\showcasegameicons{mp5k}{2456}
\showcasegameicons{mucous-pillar}{2457}
\showcasegameicons{mug-shot}{2458}
\showcasegameicons{multi-directions}{2459}
\showcasegameicons{multiple-targets}{2460}
\showcasegameicons{mummy-head}{2461}
\showcasegameicons{muscle-fat}{2462}
\showcasegameicons{muscle-up}{2463}
\showcasegameicons{muscular-torso}{2464}
\showcasegameicons{mushroom-cloud}{2465}
\showcasegameicons{mushroom-gills}{2466}
\showcasegameicons{mushroom-house}{2467}
\showcasegameicons{mushroom}{2468}
\showcasegameicons{mushrooms-cluster}{2469}
\showcasegameicons{mushrooms}{2470}
\showcasegameicons{music-spell}{2471}
\showcasegameicons{musical-keyboard}{2472}
\showcasegameicons{musical-notes}{2473}
\showcasegameicons{musical-score}{2474}
\showcasegameicons{musket}{2475}
\showcasegameicons{mussel}{2476}
\showcasegameicons{mustache}{2477}
\showcasegameicons{mute}{2478}
\showcasegameicons{nachos}{2479}
\showcasegameicons{nailed-foot}{2480}
\showcasegameicons{nailed-head}{2481}
\showcasegameicons{nails}{2482}
\showcasegameicons{nano-bot}{2483}
\showcasegameicons{nautilus-shell}{2484}
\showcasegameicons{neck-bite}{2485}
\showcasegameicons{necklace-display}{2486}
\showcasegameicons{necklace}{2487}
\showcasegameicons{nectar}{2488}
\showcasegameicons{needle-drill}{2489}
\showcasegameicons{needle-jaws}{2490}
\showcasegameicons{nefertiti}{2491}
\showcasegameicons{nest-birds}{2492}
\showcasegameicons{nest-eggs}{2493}
\showcasegameicons{nested-eclipses}{2494}
\showcasegameicons{nested-hearts}{2495}
\showcasegameicons{nested-hexagons}{2496}
\showcasegameicons{network-bars}{2497}
\showcasegameicons{new-born}{2498}
\showcasegameicons{new-shoot}{2499}
\showcasegameicons{newspaper}{2500}
\showcasegameicons{next-button}{2501}
\showcasegameicons{nigeria}{2502}
\showcasegameicons{night-sky}{2503}
\showcasegameicons{night-sleep}{2504}
\showcasegameicons{night-vision}{2505}
\showcasegameicons{ninja-armor}{2506}
\showcasegameicons{ninja-head}{2507}
\showcasegameicons{ninja-heroic-stance}{2508}
\showcasegameicons{ninja-mask}{2509}
\showcasegameicons{ninja-star}{2510}
\showcasegameicons{ninja-velociraptor}{2511}
\showcasegameicons{nodular}{2512}
\showcasegameicons{noodle-ball}{2513}
\showcasegameicons{noodles}{2514}
\showcasegameicons{north-star-shuriken}{2515}
\showcasegameicons{nose-front}{2516}
\showcasegameicons{nose-side}{2517}
\showcasegameicons{notebook}{2518}
\showcasegameicons{nothing-to-say}{2519}
\showcasegameicons{nuclear-bomb}{2520}
\showcasegameicons{nuclear-plant}{2521}
\showcasegameicons{nuclear-waste}{2522}
\showcasegameicons{nuclear}{2523}
\showcasegameicons{nun-face}{2524}
\showcasegameicons{nunchaku}{2525}
\showcasegameicons{nurse-female}{2526}
\showcasegameicons{nurse-male}{2527}
\showcasegameicons{o-brick}{2528}
\showcasegameicons{oak-leaf}{2529}
\showcasegameicons{oak}{2530}
\showcasegameicons{oasis}{2531}
\showcasegameicons{oat}{2532}
\showcasegameicons{obelisk}{2533}
\showcasegameicons{observatory}{2534}
\showcasegameicons{ocarina}{2535}
\showcasegameicons{occupy}{2536}
\showcasegameicons{octogonal-eye}{2537}
\showcasegameicons{octoman}{2538}
\showcasegameicons{octopus}{2539}
\showcasegameicons{oden}{2540}
\showcasegameicons{office-chair}{2541}
\showcasegameicons{offshore-platform}{2542}
\showcasegameicons{ogre}{2543}
\showcasegameicons{oil-can}{2544}
\showcasegameicons{oil-drum}{2545}
\showcasegameicons{oil-pump}{2546}
\showcasegameicons{oil-rig}{2547}
\showcasegameicons{oily-spiral}{2548}
\showcasegameicons{old-king}{2549}
\showcasegameicons{old-lantern}{2550}
\showcasegameicons{old-microphone}{2551}
\showcasegameicons{old-wagon}{2552}
\showcasegameicons{olive}{2553}
\showcasegameicons{omega}{2554}
\showcasegameicons{on-sight}{2555}
\showcasegameicons{on-target}{2556}
\showcasegameicons{one-eyed}{2557}
\showcasegameicons{oni}{2558}
\showcasegameicons{onigori}{2559}
\showcasegameicons{open-book}{2560}
\showcasegameicons{open-chest}{2561}
\showcasegameicons{open-folder}{2562}
\showcasegameicons{open-gate}{2563}
\showcasegameicons{open-palm}{2564}
\showcasegameicons{open-treasure-chest}{2565}
\showcasegameicons{open-wound}{2566}
\showcasegameicons{opened-food-can}{2567}
\showcasegameicons{opening-shell}{2568}
\showcasegameicons{ophanim}{2569}
\showcasegameicons{ophiuchus}{2570}
\showcasegameicons{oppidum}{2571}
\showcasegameicons{opposite-hearts}{2572}
\showcasegameicons{oppression}{2573}
\showcasegameicons{orange-slice}{2574}
\showcasegameicons{orange}{2575}
\showcasegameicons{orb-direction}{2576}
\showcasegameicons{orb-wand}{2577}
\showcasegameicons{orbit}{2578}
\showcasegameicons{orbital-rays}{2579}
\showcasegameicons{orbital}{2580}
\showcasegameicons{orc-head}{2581}
\showcasegameicons{ore}{2582}
\showcasegameicons{organigram}{2583}
\showcasegameicons{orrery}{2584}
\showcasegameicons{ostrich}{2585}
\showcasegameicons{otamatone}{2586}
\showcasegameicons{ouroboros}{2587}
\showcasegameicons{outback-hat}{2588}
\showcasegameicons{over-infinity}{2589}
\showcasegameicons{overdose}{2590}
\showcasegameicons{overdrive}{2591}
\showcasegameicons{overhead}{2592}
\showcasegameicons{overkill}{2593}
\showcasegameicons{overlord-helm}{2594}
\showcasegameicons{overmind}{2595}
\showcasegameicons{owl}{2596}
\showcasegameicons{oyster-pearl}{2597}
\showcasegameicons{oyster}{2598}
\showcasegameicons{p90}{2599}
\showcasegameicons{packed-planks}{2600}
\showcasegameicons{paddle-steamer}{2601}
\showcasegameicons{paddles}{2602}
\showcasegameicons{padlock-open}{2603}
\showcasegameicons{padlock}{2604}
\showcasegameicons{pagoda}{2605}
\showcasegameicons{paint-brush}{2606}
\showcasegameicons{paint-bucket}{2607}
\showcasegameicons{paint-roller}{2608}
\showcasegameicons{painted-pottery}{2609}
\showcasegameicons{palette}{2610}
\showcasegameicons{palisade}{2611}
\showcasegameicons{palm-tree}{2612}
\showcasegameicons{palm}{2613}
\showcasegameicons{pan-flute}{2614}
\showcasegameicons{panda}{2615}
\showcasegameicons{pangolin}{2616}
\showcasegameicons{panzerfaust}{2617}
\showcasegameicons{paper-arrow}{2618}
\showcasegameicons{paper-bag-crumpled}{2619}
\showcasegameicons{paper-bag-folded}{2620}
\showcasegameicons{paper-bag-open}{2621}
\showcasegameicons{paper-boat}{2622}
\showcasegameicons{paper-bomb}{2623}
\showcasegameicons{paper-clip}{2624}
\showcasegameicons{paper-crane}{2625}
\showcasegameicons{paper-frog}{2626}
\showcasegameicons{paper-lantern}{2627}
\showcasegameicons{paper-plane}{2628}
\showcasegameicons{paper-tray}{2629}
\showcasegameicons{paper-windmill}{2630}
\showcasegameicons{paper}{2631}
\showcasegameicons{papers}{2632}
\showcasegameicons{papyrus}{2633}
\showcasegameicons{parachute}{2634}
\showcasegameicons{paraguay}{2635}
\showcasegameicons{paranoia}{2636}
\showcasegameicons{parasaurolophus}{2637}
\showcasegameicons{park-bench}{2638}
\showcasegameicons{parmecia}{2639}
\showcasegameicons{parrot-head}{2640}
\showcasegameicons{party-flags}{2641}
\showcasegameicons{party-hat}{2642}
\showcasegameicons{party-popper}{2643}
\showcasegameicons{passport}{2644}
\showcasegameicons{path-distance}{2645}
\showcasegameicons{path-tile}{2646}
\showcasegameicons{pauldrons}{2647}
\showcasegameicons{pause-button}{2648}
\showcasegameicons{paw-front}{2649}
\showcasegameicons{paw-heart}{2650}
\showcasegameicons{paw-print}{2651}
\showcasegameicons{paw}{2652}
\showcasegameicons{pawn}{2653}
\showcasegameicons{pay-money}{2654}
\showcasegameicons{pc}{2655}
\showcasegameicons{peace-dove}{2656}
\showcasegameicons{peach}{2657}
\showcasegameicons{peaks}{2658}
\showcasegameicons{peanut}{2659}
\showcasegameicons{pear}{2660}
\showcasegameicons{pearl-earring}{2661}
\showcasegameicons{pearl-necklace}{2662}
\showcasegameicons{peas}{2663}
\showcasegameicons{pegasus}{2664}
\showcasegameicons{pelvis-bone}{2665}
\showcasegameicons{pencil-brush}{2666}
\showcasegameicons{pencil-ruler}{2667}
\showcasegameicons{pencil}{2668}
\showcasegameicons{pendant-key}{2669}
\showcasegameicons{pendulum-swing}{2670}
\showcasegameicons{penguin}{2671}
\showcasegameicons{pentacle}{2672}
\showcasegameicons{pentagram-rose}{2673}
\showcasegameicons{pentarrows-tornado}{2674}
\showcasegameicons{perfume-bottle}{2675}
\showcasegameicons{periscope}{2676}
\showcasegameicons{perpendicular-rings}{2677}
\showcasegameicons{person-in-bed}{2678}
\showcasegameicons{person-in-blizzard}{2679}
\showcasegameicons{person}{2680}
\showcasegameicons{perspective-dice-five}{2681}
\showcasegameicons{perspective-dice-four}{2682}
\showcasegameicons{perspective-dice-one}{2683}
\showcasegameicons{perspective-dice-six-faces-five}{2684}
\showcasegameicons{perspective-dice-six-faces-four}{2685}
\showcasegameicons{perspective-dice-six-faces-one}{2686}
\showcasegameicons{perspective-dice-six-faces-random}{2687}
\showcasegameicons{perspective-dice-six-faces-six}{2688}
\showcasegameicons{perspective-dice-six-faces-three}{2689}
\showcasegameicons{perspective-dice-six-faces-two}{2690}
\showcasegameicons{perspective-dice-six}{2691}
\showcasegameicons{perspective-dice-three}{2692}
\showcasegameicons{perspective-dice-two}{2693}
\showcasegameicons{peru}{2694}
\showcasegameicons{pestle-mortar}{2695}
\showcasegameicons{pharoah}{2696}
\showcasegameicons{philosopher-bust}{2697}
\showcasegameicons{phone}{2698}
\showcasegameicons{photo-camera}{2699}
\showcasegameicons{phrygian-cap}{2700}
\showcasegameicons{pianist}{2701}
\showcasegameicons{piano-keys}{2702}
\showcasegameicons{pick-of-destiny}{2703}
\showcasegameicons{pickelhaube}{2704}
\showcasegameicons{pickle}{2705}
\showcasegameicons{pie-chart}{2706}
\showcasegameicons{pie-slice}{2707}
\showcasegameicons{piece-skull}{2708}
\showcasegameicons{pierced-body}{2709}
\showcasegameicons{pierced-heart}{2710}
\showcasegameicons{piercing-sword}{2711}
\showcasegameicons{pig-face}{2712}
\showcasegameicons{pig}{2713}
\showcasegameicons{piggy-bank}{2714}
\showcasegameicons{pikeman}{2715}
\showcasegameicons{pilgrim-hat}{2716}
\showcasegameicons{pill-drop}{2717}
\showcasegameicons{pill}{2718}
\showcasegameicons{pillow}{2719}
\showcasegameicons{pimiento}{2720}
\showcasegameicons{pin}{2721}
\showcasegameicons{pinata}{2722}
\showcasegameicons{pinball-flipper}{2723}
\showcasegameicons{pincers}{2724}
\showcasegameicons{pine-tree}{2725}
\showcasegameicons{pineapple}{2726}
\showcasegameicons{ping-pong-bat}{2727}
\showcasegameicons{pipe-organ}{2728}
\showcasegameicons{pipes}{2729}
\showcasegameicons{piranha}{2730}
\showcasegameicons{pirate-cannon}{2731}
\showcasegameicons{pirate-captain}{2732}
\showcasegameicons{pirate-coat}{2733}
\showcasegameicons{pirate-flag}{2734}
\showcasegameicons{pirate-grave}{2735}
\showcasegameicons{pirate-hat}{2736}
\showcasegameicons{pirate-hook}{2737}
\showcasegameicons{pirate-skull}{2738}
\showcasegameicons{pisa-tower}{2739}
\showcasegameicons{pisces}{2740}
\showcasegameicons{pistol-gun}{2741}
\showcasegameicons{pitchfork}{2742}
\showcasegameicons{pizza-cutter}{2743}
\showcasegameicons{pizza-slice}{2744}
\showcasegameicons{plague-doctor-profile}{2745}
\showcasegameicons{plain-arrow}{2746}
\showcasegameicons{plain-circle}{2747}
\showcasegameicons{plain-dagger}{2748}
\showcasegameicons{plain-padlock}{2749}
\showcasegameicons{plain-square}{2750}
\showcasegameicons{plane-pilot}{2751}
\showcasegameicons{plane-wing}{2752}
\showcasegameicons{planet-conquest}{2753}
\showcasegameicons{planet-core}{2754}
\showcasegameicons{planks}{2755}
\showcasegameicons{plant-roots}{2756}
\showcasegameicons{plant-seed}{2757}
\showcasegameicons{plant-watering}{2758}
\showcasegameicons{plants-and-animals}{2759}
\showcasegameicons{plasma-bolt}{2760}
\showcasegameicons{plastic-duck}{2761}
\showcasegameicons{plastron}{2762}
\showcasegameicons{plate-claw}{2763}
\showcasegameicons{platform}{2764}
\showcasegameicons{play-button}{2765}
\showcasegameicons{player-base}{2766}
\showcasegameicons{player-next}{2767}
\showcasegameicons{player-previous}{2768}
\showcasegameicons{player-time}{2769}
\showcasegameicons{plesiosaurus}{2770}
\showcasegameicons{plow}{2771}
\showcasegameicons{plug}{2772}
\showcasegameicons{plum}{2773}
\showcasegameicons{plunger}{2774}
\showcasegameicons{pocket-bow}{2775}
\showcasegameicons{pocket-radio}{2776}
\showcasegameicons{pocket-watch}{2777}
\showcasegameicons{podium-second}{2778}
\showcasegameicons{podium-third}{2779}
\showcasegameicons{podium-winner}{2780}
\showcasegameicons{podium}{2781}
\showcasegameicons{pointing}{2782}
\showcasegameicons{pointy-hat}{2783}
\showcasegameicons{pointy-sword}{2784}
\showcasegameicons{poison-bottle}{2785}
\showcasegameicons{poison-cloud}{2786}
\showcasegameicons{poison-gas}{2787}
\showcasegameicons{poison}{2788}
\showcasegameicons{pokecog}{2789}
\showcasegameicons{poker-hand}{2790}
\showcasegameicons{poland}{2791}
\showcasegameicons{polar-bear}{2792}
\showcasegameicons{polar-star}{2793}
\showcasegameicons{police-badge}{2794}
\showcasegameicons{police-car}{2795}
\showcasegameicons{police-officer-head}{2796}
\showcasegameicons{police-target}{2797}
\showcasegameicons{pollen-dust}{2798}
\showcasegameicons{polo-shirt}{2799}
\showcasegameicons{poncho}{2800}
\showcasegameicons{pool-dive}{2801}
\showcasegameicons{pool-table-corner}{2802}
\showcasegameicons{pool-triangle}{2803}
\showcasegameicons{popcorn}{2804}
\showcasegameicons{pope-crown}{2805}
\showcasegameicons{poppy}{2806}
\showcasegameicons{porcelain-vase}{2807}
\showcasegameicons{porcupine}{2808}
\showcasegameicons{porcupinefish}{2809}
\showcasegameicons{portal}{2810}
\showcasegameicons{portculis}{2811}
\showcasegameicons{portrait}{2812}
\showcasegameicons{portugal}{2813}
\showcasegameicons{position-marker}{2814}
\showcasegameicons{post-office}{2815}
\showcasegameicons{post-stamp}{2816}
\showcasegameicons{potato}{2817}
\showcasegameicons{potion-ball}{2818}
\showcasegameicons{potion-of-madness}{2819}
\showcasegameicons{pouch-with-beads}{2820}
\showcasegameicons{pounce}{2821}
\showcasegameicons{pouring-chalice}{2822}
\showcasegameicons{pouring-pot}{2823}
\showcasegameicons{powder-bag}{2824}
\showcasegameicons{powder}{2825}
\showcasegameicons{power-button}{2826}
\showcasegameicons{power-generator}{2827}
\showcasegameicons{power-lightning}{2828}
\showcasegameicons{power-ring}{2829}
\showcasegameicons{prank-glasses}{2830}
\showcasegameicons{prayer-beads}{2831}
\showcasegameicons{prayer}{2832}
\showcasegameicons{praying-mantis}{2833}
\showcasegameicons{present}{2834}
\showcasegameicons{pressure-cooker}{2835}
\showcasegameicons{pretty-fangs}{2836}
\showcasegameicons{pretzel}{2837}
\showcasegameicons{previous-button}{2838}
\showcasegameicons{price-tag}{2839}
\showcasegameicons{primitive-necklace}{2840}
\showcasegameicons{primitive-torch}{2841}
\showcasegameicons{prism}{2842}
\showcasegameicons{prisoner}{2843}
\showcasegameicons{private-first-class}{2844}
\showcasegameicons{private}{2845}
\showcasegameicons{processor}{2846}
\showcasegameicons{profit}{2847}
\showcasegameicons{progression}{2848}
\showcasegameicons{propeller-beanie}{2849}
\showcasegameicons{protection-glasses}{2850}
\showcasegameicons{pschent-double-crown}{2851}
\showcasegameicons{psychic-waves}{2852}
\showcasegameicons{pterodactylus}{2853}
\showcasegameicons{pteruges}{2854}
\showcasegameicons{public-speaker}{2855}
\showcasegameicons{pull}{2856}
\showcasegameicons{pulley-hook}{2857}
\showcasegameicons{pulse}{2858}
\showcasegameicons{pummeled}{2859}
\showcasegameicons{pumpkin-lantern}{2860}
\showcasegameicons{pumpkin-mask}{2861}
\showcasegameicons{pumpkin}{2862}
\showcasegameicons{punch-blast}{2863}
\showcasegameicons{punch}{2864}
\showcasegameicons{punching-bag}{2865}
\showcasegameicons{puppet}{2866}
\showcasegameicons{purple-tentacle}{2867}
\showcasegameicons{push}{2868}
\showcasegameicons{puzzle}{2869}
\showcasegameicons{pylon}{2870}
\showcasegameicons{pyre}{2871}
\showcasegameicons{pyromaniac}{2872}
\showcasegameicons{qaitbay-citadel}{2873}
\showcasegameicons{quake-stomp}{2874}
\showcasegameicons{queen-crown}{2875}
\showcasegameicons{quick-man}{2876}
\showcasegameicons{quick-slash}{2877}
\showcasegameicons{quicksand}{2878}
\showcasegameicons{quill-ink}{2879}
\showcasegameicons{quill}{2880}
\showcasegameicons{quiver}{2881}
\showcasegameicons{rabbit-head}{2882}
\showcasegameicons{rabbit}{2883}
\showcasegameicons{raccoon-head}{2884}
\showcasegameicons{race-car}{2885}
\showcasegameicons{radar-cross-section}{2886}
\showcasegameicons{radar-dish}{2887}
\showcasegameicons{radar-sweep}{2888}
\showcasegameicons{raddish}{2889}
\showcasegameicons{radial-balance}{2890}
\showcasegameicons{radiations}{2891}
\showcasegameicons{radio-tower}{2892}
\showcasegameicons{radioactive}{2893}
\showcasegameicons{raft}{2894}
\showcasegameicons{ragged-wound}{2895}
\showcasegameicons{rail-road}{2896}
\showcasegameicons{railway}{2897}
\showcasegameicons{rainbow-star}{2898}
\showcasegameicons{raining}{2899}
\showcasegameicons{raise-skeleton}{2900}
\showcasegameicons{raise-zombie}{2901}
\showcasegameicons{rake}{2902}
\showcasegameicons{rally-the-troops}{2903}
\showcasegameicons{ram-profile}{2904}
\showcasegameicons{ram}{2905}
\showcasegameicons{ranch-gate}{2906}
\showcasegameicons{rank-1}{2907}
\showcasegameicons{rank-2}{2908}
\showcasegameicons{rank-3}{2909}
\showcasegameicons{rapidshare-arrow}{2910}
\showcasegameicons{raspberry}{2911}
\showcasegameicons{rat}{2912}
\showcasegameicons{rattlesnake}{2913}
\showcasegameicons{raven}{2914}
\showcasegameicons{raw-egg}{2915}
\showcasegameicons{ray-gun}{2916}
\showcasegameicons{razor-blade}{2917}
\showcasegameicons{razor}{2918}
\showcasegameicons{reactor}{2919}
\showcasegameicons{read}{2920}
\showcasegameicons{reaper-scythe}{2921}
\showcasegameicons{rear-aura}{2922}
\showcasegameicons{receive-money}{2923}
\showcasegameicons{recycle}{2924}
\showcasegameicons{red-carpet}{2925}
\showcasegameicons{reed}{2926}
\showcasegameicons{refinery}{2927}
\showcasegameicons{regeneration}{2928}
\showcasegameicons{relationship-bounds}{2929}
\showcasegameicons{relic-blade}{2930}
\showcasegameicons{reload-gun-barrel}{2931}
\showcasegameicons{remedy}{2932}
\showcasegameicons{rempart}{2933}
\showcasegameicons{reptile-tail}{2934}
\showcasegameicons{resize}{2935}
\showcasegameicons{resonance}{2936}
\showcasegameicons{resting-vampire}{2937}
\showcasegameicons{reticule}{2938}
\showcasegameicons{retro-controller}{2939}
\showcasegameicons{return-arrow}{2940}
\showcasegameicons{revolt}{2941}
\showcasegameicons{revolver}{2942}
\showcasegameicons{rhinoceros-horn}{2943}
\showcasegameicons{rialto-bridge}{2944}
\showcasegameicons{ribbon-medal}{2945}
\showcasegameicons{ribbon-shield}{2946}
\showcasegameicons{ribbon}{2947}
\showcasegameicons{ribcage}{2948}
\showcasegameicons{rice-cooker}{2949}
\showcasegameicons{ricochet}{2950}
\showcasegameicons{rifle}{2951}
\showcasegameicons{ring-box}{2952}
\showcasegameicons{ring-mould}{2953}
\showcasegameicons{ring}{2954}
\showcasegameicons{ringed-beam}{2955}
\showcasegameicons{ringed-planet}{2956}
\showcasegameicons{ringed-tentacle}{2957}
\showcasegameicons{ringing-alarm}{2958}
\showcasegameicons{ringing-bell}{2959}
\showcasegameicons{ringmaster}{2960}
\showcasegameicons{riot-shield}{2961}
\showcasegameicons{riposte}{2962}
\showcasegameicons{river}{2963}
\showcasegameicons{road}{2964}
\showcasegameicons{roast-chicken}{2965}
\showcasegameicons{robber-hand}{2966}
\showcasegameicons{robber-mask}{2967}
\showcasegameicons{robber}{2968}
\showcasegameicons{robe}{2969}
\showcasegameicons{robin-hood-hat}{2970}
\showcasegameicons{robot-antennas}{2971}
\showcasegameicons{robot-golem}{2972}
\showcasegameicons{robot-grab}{2973}
\showcasegameicons{robot-helmet}{2974}
\showcasegameicons{robot-leg}{2975}
\showcasegameicons{rock-golem}{2976}
\showcasegameicons{rock}{2977}
\showcasegameicons{rocket-flight}{2978}
\showcasegameicons{rocket-thruster}{2979}
\showcasegameicons{rocket}{2980}
\showcasegameicons{rocking-chair}{2981}
\showcasegameicons{rod-of-asclepius}{2982}
\showcasegameicons{rogue}{2983}
\showcasegameicons{rolled-cloth}{2984}
\showcasegameicons{roller-skate}{2985}
\showcasegameicons{rolling-bomb}{2986}
\showcasegameicons{rolling-dice-cup}{2987}
\showcasegameicons{rolling-dices}{2988}
\showcasegameicons{rolling-energy}{2989}
\showcasegameicons{rolling-suitcase}{2990}
\showcasegameicons{roman-shield}{2991}
\showcasegameicons{roman-toga}{2992}
\showcasegameicons{rooster}{2993}
\showcasegameicons{root-tip}{2994}
\showcasegameicons{rope-bridge}{2995}
\showcasegameicons{rope-coil}{2996}
\showcasegameicons{rope-dart}{2997}
\showcasegameicons{ropeway}{2998}
\showcasegameicons{rosa-shield}{2999}
\showcasegameicons{rose}{3000}
\showcasegameicons{rotary-phone}{3001}
\showcasegameicons{rough-wound}{3002}
\showcasegameicons{round-bottom-flask}{3003}
\showcasegameicons{round-knob}{3004}
\showcasegameicons{round-potion}{3005}
\showcasegameicons{round-shield}{3006}
\showcasegameicons{round-silo}{3007}
\showcasegameicons{round-star}{3008}
\showcasegameicons{round-straw-bale}{3009}
\showcasegameicons{round-struck}{3010}
\showcasegameicons{round-table}{3011}
\showcasegameicons{royal-love}{3012}
\showcasegameicons{rss}{3013}
\showcasegameicons{rub-el-hizb}{3014}
\showcasegameicons{rubber-boot}{3015}
\showcasegameicons{rugby-conversion}{3016}
\showcasegameicons{rule-book}{3017}
\showcasegameicons{run}{3018}
\showcasegameicons{rune-stone}{3019}
\showcasegameicons{rune-sword}{3020}
\showcasegameicons{running-ninja}{3021}
\showcasegameicons{running-shoe}{3022}
\showcasegameicons{rupee}{3023}
\showcasegameicons{rusty-sword}{3024}
\showcasegameicons{s-brick}{3025}
\showcasegameicons{saber-and-pistol}{3026}
\showcasegameicons{saber-slash}{3027}
\showcasegameicons{saber-tooth}{3028}
\showcasegameicons{saber-toothed-cat-head}{3029}
\showcasegameicons{sabers-choc}{3030}
\showcasegameicons{sacrificial-dagger}{3031}
\showcasegameicons{sad-crab}{3032}
\showcasegameicons{saddle}{3033}
\showcasegameicons{safety-pin}{3034}
\showcasegameicons{sagittarius}{3035}
\showcasegameicons{sai}{3036}
\showcasegameicons{sail}{3037}
\showcasegameicons{sailboat}{3038}
\showcasegameicons{saint-basil-cathedral}{3039}
\showcasegameicons{saiyan-suit}{3040}
\showcasegameicons{sake-bottle}{3041}
\showcasegameicons{salamander}{3042}
\showcasegameicons{salmon}{3043}
\showcasegameicons{saloon-doors}{3044}
\showcasegameicons{saloon}{3045}
\showcasegameicons{salt-shaker}{3046}
\showcasegameicons{samara-mosque}{3047}
\showcasegameicons{samurai-helmet}{3048}
\showcasegameicons{samus-helmet}{3049}
\showcasegameicons{sand-castle}{3050}
\showcasegameicons{sand-snake}{3051}
\showcasegameicons{sandal}{3052}
\showcasegameicons{sands-of-time}{3053}
\showcasegameicons{sandstorm}{3054}
\showcasegameicons{sandwich}{3055}
\showcasegameicons{santa-hat}{3056}
\showcasegameicons{saphir}{3057}
\showcasegameicons{sarcophagus}{3058}
\showcasegameicons{sasquatch}{3059}
\showcasegameicons{satellite-communication}{3060}
\showcasegameicons{satellite}{3061}
\showcasegameicons{saucepan}{3062}
\showcasegameicons{sauropod-head}{3063}
\showcasegameicons{sauropod-skeleton}{3064}
\showcasegameicons{sausage}{3065}
\showcasegameicons{sausages-ribbon}{3066}
\showcasegameicons{save-arrow}{3067}
\showcasegameicons{save}{3068}
\showcasegameicons{saw-claw}{3069}
\showcasegameicons{sawed-off-shotgun}{3070}
\showcasegameicons{saxophone}{3071}
\showcasegameicons{scabbard}{3072}
\showcasegameicons{scale-mail}{3073}
\showcasegameicons{scales}{3074}
\showcasegameicons{scallop}{3075}
\showcasegameicons{scalpel-strike}{3076}
\showcasegameicons{scalpel}{3077}
\showcasegameicons{scar-wound}{3078}
\showcasegameicons{scarab-beetle}{3079}
\showcasegameicons{scarecrow}{3080}
\showcasegameicons{school-bag}{3081}
\showcasegameicons{school-of-fish}{3082}
\showcasegameicons{schrodingers-cat-alive}{3083}
\showcasegameicons{schrodingers-cat-dead}{3084}
\showcasegameicons{scissors}{3085}
\showcasegameicons{scooter}{3086}
\showcasegameicons{scorpio}{3087}
\showcasegameicons{scorpion-tail}{3088}
\showcasegameicons{scorpion}{3089}
\showcasegameicons{scout-ship}{3090}
\showcasegameicons{screaming}{3091}
\showcasegameicons{screen-impact}{3092}
\showcasegameicons{screw}{3093}
\showcasegameicons{screwdriver}{3094}
\showcasegameicons{scroll-quill}{3095}
\showcasegameicons{scroll-unfurled}{3096}
\showcasegameicons{scuba-mask}{3097}
\showcasegameicons{scuba-tanks}{3098}
\showcasegameicons{scythe}{3099}
\showcasegameicons{sea-cliff}{3100}
\showcasegameicons{sea-creature}{3101}
\showcasegameicons{sea-dragon}{3102}
\showcasegameicons{sea-serpent}{3103}
\showcasegameicons{sea-star}{3104}
\showcasegameicons{sea-turtle}{3105}
\showcasegameicons{sea-urchin}{3106}
\showcasegameicons{seagull}{3107}
\showcasegameicons{seahorse}{3108}
\showcasegameicons{seated-mouse}{3109}
\showcasegameicons{secret-book}{3110}
\showcasegameicons{secret-door}{3111}
\showcasegameicons{security-gate}{3112}
\showcasegameicons{seedling}{3113}
\showcasegameicons{select}{3114}
\showcasegameicons{self-love}{3115}
\showcasegameicons{sell-card}{3116}
\showcasegameicons{semi-closed-eye}{3117}
\showcasegameicons{sensuousness}{3118}
\showcasegameicons{sentry-gun}{3119}
\showcasegameicons{sergeant}{3120}
\showcasegameicons{serrated-slash}{3121}
\showcasegameicons{server-rack}{3122}
\showcasegameicons{sesame}{3123}
\showcasegameicons{settings-knobs}{3124}
\showcasegameicons{seven-pointed-star}{3125}
\showcasegameicons{severed-hand}{3126}
\showcasegameicons{sewed-shell}{3127}
\showcasegameicons{sewing-machine}{3128}
\showcasegameicons{sewing-needle}{3129}
\showcasegameicons{sewing-string}{3130}
\showcasegameicons{sextant}{3131}
\showcasegameicons{shadow-follower}{3132}
\showcasegameicons{shadow-grasp}{3133}
\showcasegameicons{shaking-hands}{3134}
\showcasegameicons{shambling-mound}{3135}
\showcasegameicons{shambling-zombie}{3136}
\showcasegameicons{shamrock}{3137}
\showcasegameicons{shard-sword}{3138}
\showcasegameicons{share}{3139}
\showcasegameicons{shark-bite}{3140}
\showcasegameicons{shark-fin}{3141}
\showcasegameicons{shark-jaws}{3142}
\showcasegameicons{sharp-axe}{3143}
\showcasegameicons{sharp-crown}{3144}
\showcasegameicons{sharp-halberd}{3145}
\showcasegameicons{sharp-lips}{3146}
\showcasegameicons{sharp-shuriken}{3147}
\showcasegameicons{sharp-smile}{3148}
\showcasegameicons{sharped-teeth-skull}{3149}
\showcasegameicons{shatter}{3150}
\showcasegameicons{shattered-glass}{3151}
\showcasegameicons{shattered-heart}{3152}
\showcasegameicons{shattered-sword}{3153}
\showcasegameicons{shears}{3154}
\showcasegameicons{sheep}{3155}
\showcasegameicons{sheikah-eye}{3156}
\showcasegameicons{shepherds-crook}{3157}
\showcasegameicons{sherlock-holmes}{3158}
\showcasegameicons{shield-bash}{3159}
\showcasegameicons{shield-bounces}{3160}
\showcasegameicons{shield-disabled}{3161}
\showcasegameicons{shield-echoes}{3162}
\showcasegameicons{shield-impact}{3163}
\showcasegameicons{shield-opposition}{3164}
\showcasegameicons{shield-reflect}{3165}
\showcasegameicons{shield}{3166}
\showcasegameicons{shieldcomb}{3167}
\showcasegameicons{shining-claw}{3168}
\showcasegameicons{shining-heart}{3169}
\showcasegameicons{shining-sword}{3170}
\showcasegameicons{shinto-shrine-mirror}{3171}
\showcasegameicons{shinto-shrine}{3172}
\showcasegameicons{shiny-apple}{3173}
\showcasegameicons{shiny-entrance}{3174}
\showcasegameicons{shiny-iris}{3175}
\showcasegameicons{shiny-omega}{3176}
\showcasegameicons{shiny-purse}{3177}
\showcasegameicons{ship-bow}{3178}
\showcasegameicons{ship-wheel}{3179}
\showcasegameicons{ship-wreck}{3180}
\showcasegameicons{shirt-button}{3181}
\showcasegameicons{shirt}{3182}
\showcasegameicons{shoebill-stork}{3183}
\showcasegameicons{shooner-sailboat}{3184}
\showcasegameicons{shop}{3185}
\showcasegameicons{shopping-bag}{3186}
\showcasegameicons{shopping-cart}{3187}
\showcasegameicons{shorts}{3188}
\showcasegameicons{shotgun-rounds}{3189}
\showcasegameicons{shotgun}{3190}
\showcasegameicons{shoulder-armor}{3191}
\showcasegameicons{shoulder-bag}{3192}
\showcasegameicons{shoulder-scales}{3193}
\showcasegameicons{shouting}{3194}
\showcasegameicons{shower}{3195}
\showcasegameicons{shrimp}{3196}
\showcasegameicons{shrug}{3197}
\showcasegameicons{shuriken-aperture}{3198}
\showcasegameicons{shuriken}{3199}
\showcasegameicons{shut-rose}{3200}
\showcasegameicons{shuttlecock}{3201}
\showcasegameicons{sickle}{3202}
\showcasegameicons{sideswipe}{3203}
\showcasegameicons{siege-ram}{3204}
\showcasegameicons{siege-tower}{3205}
\showcasegameicons{sight-disabled}{3206}
\showcasegameicons{silence}{3207}
\showcasegameicons{silenced}{3208}
\showcasegameicons{silex}{3209}
\showcasegameicons{silver-bullet}{3210}
\showcasegameicons{sinagot}{3211}
\showcasegameicons{sing}{3212}
\showcasegameicons{sink}{3213}
\showcasegameicons{sinking-ship}{3214}
\showcasegameicons{sinking-trap}{3215}
\showcasegameicons{sinusoidal-beam}{3216}
\showcasegameicons{siren}{3217}
\showcasegameicons{sitting-dog}{3218}
\showcasegameicons{six-eyes}{3219}
\showcasegameicons{skateboard}{3220}
\showcasegameicons{skateboarding-01}{3221}
\showcasegameicons{skateboarding-02}{3222}
\showcasegameicons{skeletal-hand}{3223}
\showcasegameicons{skeleton-inside}{3224}
\showcasegameicons{skeleton-key}{3225}
\showcasegameicons{skeleton}{3226}
\showcasegameicons{ski-boot}{3227}
\showcasegameicons{skid-mark}{3228}
\showcasegameicons{skier}{3229}
\showcasegameicons{skills}{3230}
\showcasegameicons{skimmer-hat}{3231}
\showcasegameicons{skipping-rope}{3232}
\showcasegameicons{skirt}{3233}
\showcasegameicons{skis}{3234}
\showcasegameicons{skull-bolt}{3235}
\showcasegameicons{skull-crack}{3236}
\showcasegameicons{skull-crossed-bones}{3237}
\showcasegameicons{skull-in-jar}{3238}
\showcasegameicons{skull-mask}{3239}
\showcasegameicons{skull-ring}{3240}
\showcasegameicons{skull-sabertooth}{3241}
\showcasegameicons{skull-shield}{3242}
\showcasegameicons{skull-signet}{3243}
\showcasegameicons{skull-slices}{3244}
\showcasegameicons{skull-staff}{3245}
\showcasegameicons{skull-with-syringe}{3246}
\showcasegameicons{slalom}{3247}
\showcasegameicons{slap}{3248}
\showcasegameicons{slashed-shield}{3249}
\showcasegameicons{slavery-whip}{3250}
\showcasegameicons{sleeping-bag}{3251}
\showcasegameicons{sleepy}{3252}
\showcasegameicons{sleeveless-jacket}{3253}
\showcasegameicons{sleeveless-top}{3254}
\showcasegameicons{sliced-bread}{3255}
\showcasegameicons{sliced-mushroom}{3256}
\showcasegameicons{sliced-sausage}{3257}
\showcasegameicons{slicing-arrow}{3258}
\showcasegameicons{slime}{3259}
\showcasegameicons{sling}{3260}
\showcasegameicons{slingshot}{3261}
\showcasegameicons{slipknot}{3262}
\showcasegameicons{slippers}{3263}
\showcasegameicons{slot-machine}{3264}
\showcasegameicons{sloth}{3265}
\showcasegameicons{slow-blob}{3266}
\showcasegameicons{slumbering-sanctuary}{3267}
\showcasegameicons{sly}{3268}
\showcasegameicons{small-fire}{3269}
\showcasegameicons{small-fishing-sailboat}{3270}
\showcasegameicons{smart}{3271}
\showcasegameicons{smartphone}{3272}
\showcasegameicons{smash-arrows}{3273}
\showcasegameicons{smitten}{3274}
\showcasegameicons{smoke-bomb}{3275}
\showcasegameicons{smoking-finger}{3276}
\showcasegameicons{smoking-orb}{3277}
\showcasegameicons{smoking-pipe}{3278}
\showcasegameicons{smoking-volcano}{3279}
\showcasegameicons{snail-eyes}{3280}
\showcasegameicons{snail}{3281}
\showcasegameicons{snake-bite}{3282}
\showcasegameicons{snake-egg}{3283}
\showcasegameicons{snake-jar}{3284}
\showcasegameicons{snake-spiral}{3285}
\showcasegameicons{snake-tongue}{3286}
\showcasegameicons{snake-totem}{3287}
\showcasegameicons{snake}{3288}
\showcasegameicons{snatch}{3289}
\showcasegameicons{sniffing-dog}{3290}
\showcasegameicons{snitch-quidditch-ball}{3291}
\showcasegameicons{snorkel}{3292}
\showcasegameicons{snout}{3293}
\showcasegameicons{snow-bottle}{3294}
\showcasegameicons{snowboard}{3295}
\showcasegameicons{snowflake-1}{3296}
\showcasegameicons{snowflake-2}{3297}
\showcasegameicons{snowing}{3298}
\showcasegameicons{snowman}{3299}
\showcasegameicons{soap-experiment}{3300}
\showcasegameicons{soap}{3301}
\showcasegameicons{soccer-ball}{3302}
\showcasegameicons{soccer-field}{3303}
\showcasegameicons{soccer-kick}{3304}
\showcasegameicons{socks}{3305}
\showcasegameicons{soda-bottle}{3306}
\showcasegameicons{soda-can}{3307}
\showcasegameicons{sofa}{3308}
\showcasegameicons{solar-power}{3309}
\showcasegameicons{solar-system}{3310}
\showcasegameicons{solar-time}{3311}
\showcasegameicons{soldering-iron}{3312}
\showcasegameicons{solid-leaf}{3313}
\showcasegameicons{sombrero}{3314}
\showcasegameicons{sonic-boom}{3315}
\showcasegameicons{sonic-lightning}{3316}
\showcasegameicons{sonic-screech}{3317}
\showcasegameicons{sonic-shoes}{3318}
\showcasegameicons{sonic-shout}{3319}
\showcasegameicons{soul-vessel}{3320}
\showcasegameicons{soul}{3321}
\showcasegameicons{sound-off}{3322}
\showcasegameicons{sound-on}{3323}
\showcasegameicons{sound-waves}{3324}
\showcasegameicons{south-africa-flag}{3325}
\showcasegameicons{south-africa}{3326}
\showcasegameicons{south-america}{3327}
\showcasegameicons{south-korea}{3328}
\showcasegameicons{space-needle}{3329}
\showcasegameicons{space-shuttle}{3330}
\showcasegameicons{space-suit}{3331}
\showcasegameicons{spaceship}{3332}
\showcasegameicons{spade-skull}{3333}
\showcasegameicons{spade}{3334}
\showcasegameicons{spades}{3335}
\showcasegameicons{spain}{3336}
\showcasegameicons{spanner}{3337}
\showcasegameicons{spark-plug}{3338}
\showcasegameicons{spark-spirit}{3339}
\showcasegameicons{sparkles}{3340}
\showcasegameicons{sparkling-sabre}{3341}
\showcasegameicons{sparky-bomb}{3342}
\showcasegameicons{sparrow}{3343}
\showcasegameicons{spartan-helmet}{3344}
\showcasegameicons{spartan}{3345}
\showcasegameicons{spatter}{3346}
\showcasegameicons{spawn-node}{3347}
\showcasegameicons{speaker-off}{3348}
\showcasegameicons{speaker}{3349}
\showcasegameicons{spear-feather}{3350}
\showcasegameicons{spear-hook}{3351}
\showcasegameicons{spearfishing}{3352}
\showcasegameicons{spears}{3353}
\showcasegameicons{spectacle-lenses}{3354}
\showcasegameicons{spectacles}{3355}
\showcasegameicons{spectre-m4}{3356}
\showcasegameicons{spectre}{3357}
\showcasegameicons{speed-boat}{3358}
\showcasegameicons{speedometer}{3359}
\showcasegameicons{spell-book}{3360}
\showcasegameicons{sperm-whale}{3361}
\showcasegameicons{spider-alt}{3362}
\showcasegameicons{spider-bot}{3363}
\showcasegameicons{spider-eye}{3364}
\showcasegameicons{spider-face}{3365}
\showcasegameicons{spider-mask}{3366}
\showcasegameicons{spider-web}{3367}
\showcasegameicons{spikeball}{3368}
\showcasegameicons{spiked-armor}{3369}
\showcasegameicons{spiked-ball}{3370}
\showcasegameicons{spiked-bat}{3371}
\showcasegameicons{spiked-collar}{3372}
\showcasegameicons{spiked-dragon-head}{3373}
\showcasegameicons{spiked-fence}{3374}
\showcasegameicons{spiked-halo}{3375}
\showcasegameicons{spiked-mace}{3376}
\showcasegameicons{spiked-shell}{3377}
\showcasegameicons{spiked-shield}{3378}
\showcasegameicons{spiked-shoulder-armor}{3379}
\showcasegameicons{spiked-snail}{3380}
\showcasegameicons{spiked-tail}{3381}
\showcasegameicons{spiked-tentacle}{3382}
\showcasegameicons{spiked-trunk}{3383}
\showcasegameicons{spiked-wall}{3384}
\showcasegameicons{spikes-full}{3385}
\showcasegameicons{spikes-half}{3386}
\showcasegameicons{spikes-init}{3387}
\showcasegameicons{spikes}{3388}
\showcasegameicons{spiky-eclipse}{3389}
\showcasegameicons{spiky-explosion}{3390}
\showcasegameicons{spiky-field}{3391}
\showcasegameicons{spiky-pit}{3392}
\showcasegameicons{spiky-wing}{3393}
\showcasegameicons{spill}{3394}
\showcasegameicons{spinal-coil}{3395}
\showcasegameicons{spine-arrow}{3396}
\showcasegameicons{spinning-blades}{3397}
\showcasegameicons{spinning-ribbons}{3398}
\showcasegameicons{spinning-sword}{3399}
\showcasegameicons{spinning-top}{3400}
\showcasegameicons{spinning-wheel}{3401}
\showcasegameicons{spiral-arrow}{3402}
\showcasegameicons{spiral-bloom}{3403}
\showcasegameicons{spiral-bottle}{3404}
\showcasegameicons{spiral-hilt}{3405}
\showcasegameicons{spiral-lollipop}{3406}
\showcasegameicons{spiral-shell}{3407}
\showcasegameicons{spiral-tentacle}{3408}
\showcasegameicons{spiral-thrust}{3409}
\showcasegameicons{splash}{3410}
\showcasegameicons{splashy-stream}{3411}
\showcasegameicons{split-arrows}{3412}
\showcasegameicons{split-body}{3413}
\showcasegameicons{split-cross}{3414}
\showcasegameicons{splurt}{3415}
\showcasegameicons{spock-hand}{3416}
\showcasegameicons{spooky-house}{3417}
\showcasegameicons{spoon}{3418}
\showcasegameicons{sport-medal}{3419}
\showcasegameicons{spoted-flower}{3420}
\showcasegameicons{spotted-arrowhead}{3421}
\showcasegameicons{spotted-bug}{3422}
\showcasegameicons{spotted-mushroom}{3423}
\showcasegameicons{spotted-wound}{3424}
\showcasegameicons{spoutnik}{3425}
\showcasegameicons{spray}{3426}
\showcasegameicons{spring}{3427}
\showcasegameicons{sprint}{3428}
\showcasegameicons{sprout-disc}{3429}
\showcasegameicons{sprout}{3430}
\showcasegameicons{spy}{3431}
\showcasegameicons{spyglass}{3432}
\showcasegameicons{square-bottle}{3433}
\showcasegameicons{square}{3434}
\showcasegameicons{squib}{3435}
\showcasegameicons{squid-head}{3436}
\showcasegameicons{squid}{3437}
\showcasegameicons{squirrel}{3438}
\showcasegameicons{sri-lanka}{3439}
\showcasegameicons{stabbed-note}{3440}
\showcasegameicons{stable}{3441}
\showcasegameicons{stack}{3442}
\showcasegameicons{stag-head}{3443}
\showcasegameicons{stahlhelm}{3444}
\showcasegameicons{stairs-cake}{3445}
\showcasegameicons{stairs-goal}{3446}
\showcasegameicons{stairs}{3447}
\showcasegameicons{stake-hammer}{3448}
\showcasegameicons{stakes-fence}{3449}
\showcasegameicons{stalactites}{3450}
\showcasegameicons{stalagtite}{3451}
\showcasegameicons{stamper}{3452}
\showcasegameicons{standing-potion}{3453}
\showcasegameicons{stapler-heavy-duty}{3454}
\showcasegameicons{stapler-pneumatic}{3455}
\showcasegameicons{stapler}{3456}
\showcasegameicons{star-altar}{3457}
\showcasegameicons{star-cycle}{3458}
\showcasegameicons{star-flag}{3459}
\showcasegameicons{star-formation}{3460}
\showcasegameicons{star-gate}{3461}
\showcasegameicons{star-key}{3462}
\showcasegameicons{star-medal}{3463}
\showcasegameicons{star-prominences}{3464}
\showcasegameicons{star-pupil}{3465}
\showcasegameicons{star-satellites}{3466}
\showcasegameicons{star-shuriken}{3467}
\showcasegameicons{star-skull}{3468}
\showcasegameicons{star-struck}{3469}
\showcasegameicons{star-swirl}{3470}
\showcasegameicons{starfighter}{3471}
\showcasegameicons{stars-stack}{3472}
\showcasegameicons{staryu}{3473}
\showcasegameicons{static-guard}{3474}
\showcasegameicons{static-waves}{3475}
\showcasegameicons{static}{3476}
\showcasegameicons{steak}{3477}
\showcasegameicons{stealth-bomber}{3478}
\showcasegameicons{steam-blast}{3479}
\showcasegameicons{steam-locomotive}{3480}
\showcasegameicons{steam}{3481}
\showcasegameicons{steampunk-goggles}{3482}
\showcasegameicons{steamroller}{3483}
\showcasegameicons{steel-claws}{3484}
\showcasegameicons{steel-door}{3485}
\showcasegameicons{steeltoe-boots}{3486}
\showcasegameicons{steelwing-emblem}{3487}
\showcasegameicons{steering-wheel}{3488}
\showcasegameicons{stegosaurus-scales}{3489}
\showcasegameicons{stethoscope}{3490}
\showcasegameicons{steyr-aug}{3491}
\showcasegameicons{stick-frame}{3492}
\showcasegameicons{stick-grenade}{3493}
\showcasegameicons{stick-splitting}{3494}
\showcasegameicons{sticking-plaster}{3495}
\showcasegameicons{sticky-boot}{3496}
\showcasegameicons{stigmata}{3497}
\showcasegameicons{stiletto}{3498}
\showcasegameicons{stitched-wound}{3499}
\showcasegameicons{stockpiles}{3500}
\showcasegameicons{stomach}{3501}
\showcasegameicons{stomp-tornado}{3502}
\showcasegameicons{stomp}{3503}
\showcasegameicons{stone-axe}{3504}
\showcasegameicons{stone-block}{3505}
\showcasegameicons{stone-bridge}{3506}
\showcasegameicons{stone-bust}{3507}
\showcasegameicons{stone-crafting}{3508}
\showcasegameicons{stone-path}{3509}
\showcasegameicons{stone-pile}{3510}
\showcasegameicons{stone-spear}{3511}
\showcasegameicons{stone-sphere}{3512}
\showcasegameicons{stone-stack}{3513}
\showcasegameicons{stone-tablet}{3514}
\showcasegameicons{stone-throne}{3515}
\showcasegameicons{stone-tower}{3516}
\showcasegameicons{stone-wall}{3517}
\showcasegameicons{stone-wheel}{3518}
\showcasegameicons{stoned-skull}{3519}
\showcasegameicons{stop-sign}{3520}
\showcasegameicons{stopwatch}{3521}
\showcasegameicons{stork-delivery}{3522}
\showcasegameicons{strafe}{3523}
\showcasegameicons{straight-pipe}{3524}
\showcasegameicons{strawberry}{3525}
\showcasegameicons{street-light}{3526}
\showcasegameicons{striking-arrows}{3527}
\showcasegameicons{striking-balls}{3528}
\showcasegameicons{striking-clamps}{3529}
\showcasegameicons{striking-diamonds}{3530}
\showcasegameicons{striking-splinter}{3531}
\showcasegameicons{striped-sun}{3532}
\showcasegameicons{striped-sword}{3533}
\showcasegameicons{strong-man}{3534}
\showcasegameicons{strong}{3535}
\showcasegameicons{strongbox}{3536}
\showcasegameicons{stump-regrowth}{3537}
\showcasegameicons{stun-grenade}{3538}
\showcasegameicons{submarine-missile}{3539}
\showcasegameicons{submarine}{3540}
\showcasegameicons{subway-pass}{3541}
\showcasegameicons{subway-train}{3542}
\showcasegameicons{subway}{3543}
\showcasegameicons{suckered-tentacle}{3544}
\showcasegameicons{sugar-cane}{3545}
\showcasegameicons{suicide}{3546}
\showcasegameicons{suitcase}{3547}
\showcasegameicons{suits}{3548}
\showcasegameicons{summits}{3549}
\showcasegameicons{sun-cloud}{3550}
\showcasegameicons{sun-priest}{3551}
\showcasegameicons{sun-radiations}{3552}
\showcasegameicons{sun-spear}{3553}
\showcasegameicons{sun}{3554}
\showcasegameicons{sunbeams}{3555}
\showcasegameicons{sundial}{3556}
\showcasegameicons{sunflower}{3557}
\showcasegameicons{sunglasses}{3558}
\showcasegameicons{sunken-eye}{3559}
\showcasegameicons{sunrise}{3560}
\showcasegameicons{sunset}{3561}
\showcasegameicons{super-mushroom}{3562}
\showcasegameicons{supersonic-arrow}{3563}
\showcasegameicons{supersonic-bullet}{3564}
\showcasegameicons{surf-board}{3565}
\showcasegameicons{surfer-van}{3566}
\showcasegameicons{surprised-skull}{3567}
\showcasegameicons{surprised}{3568}
\showcasegameicons{surrounded-eye}{3569}
\showcasegameicons{surrounded-shield}{3570}
\showcasegameicons{sushis}{3571}
\showcasegameicons{suspension-bridge}{3572}
\showcasegameicons{suspicious}{3573}
\showcasegameicons{sverd-i-fjell}{3574}
\showcasegameicons{swallow}{3575}
\showcasegameicons{swallower}{3576}
\showcasegameicons{swamp-bat}{3577}
\showcasegameicons{swamp}{3578}
\showcasegameicons{swan-breeze}{3579}
\showcasegameicons{swan}{3580}
\showcasegameicons{swap-bag}{3581}
\showcasegameicons{swimfins}{3582}
\showcasegameicons{swipe-card}{3583}
\showcasegameicons{swirl-ring}{3584}
\showcasegameicons{swirl-string}{3585}
\showcasegameicons{swirled-shell}{3586}
\showcasegameicons{swiss-army-knife}{3587}
\showcasegameicons{switch-weapon}{3588}
\showcasegameicons{switchblade}{3589}
\showcasegameicons{switzerland}{3590}
\showcasegameicons{sword-altar}{3591}
\showcasegameicons{sword-array}{3592}
\showcasegameicons{sword-brandish}{3593}
\showcasegameicons{sword-break}{3594}
\showcasegameicons{sword-clash}{3595}
\showcasegameicons{sword-hilt}{3596}
\showcasegameicons{sword-in-stone}{3597}
\showcasegameicons{sword-mold}{3598}
\showcasegameicons{sword-slice}{3599}
\showcasegameicons{sword-smithing}{3600}
\showcasegameicons{sword-spade}{3601}
\showcasegameicons{sword-spin}{3602}
\showcasegameicons{sword-tie}{3603}
\showcasegameicons{sword-wound}{3604}
\showcasegameicons{swordman}{3605}
\showcasegameicons{swords-emblem}{3606}
\showcasegameicons{swords-power}{3607}
\showcasegameicons{swordwoman}{3608}
\showcasegameicons{sydney-opera-house}{3609}
\showcasegameicons{syringe}{3610}
\showcasegameicons{t-brick}{3611}
\showcasegameicons{t-rex-skull}{3612}
\showcasegameicons{t-shirt}{3613}
\showcasegameicons{tabi-boot}{3614}
\showcasegameicons{table}{3615}
\showcasegameicons{tablet}{3616}
\showcasegameicons{tabletop-players}{3617}
\showcasegameicons{tacos}{3618}
\showcasegameicons{tadpole}{3619}
\showcasegameicons{take-my-money}{3620}
\showcasegameicons{talk}{3621}
\showcasegameicons{tall-bridge}{3622}
\showcasegameicons{tally-mark-5}{3623}
\showcasegameicons{tambourine}{3624}
\showcasegameicons{tangerine}{3625}
\showcasegameicons{tank-top}{3626}
\showcasegameicons{tank-tread}{3627}
\showcasegameicons{tank}{3628}
\showcasegameicons{tanzania}{3629}
\showcasegameicons{tap}{3630}
\showcasegameicons{tapir}{3631}
\showcasegameicons{target-arrows}{3632}
\showcasegameicons{target-dummy}{3633}
\showcasegameicons{target-laser}{3634}
\showcasegameicons{target-poster}{3635}
\showcasegameicons{target-prize}{3636}
\showcasegameicons{target-shot}{3637}
\showcasegameicons{targeted}{3638}
\showcasegameicons{targeting}{3639}
\showcasegameicons{tarot-00-the-fool}{3640}
\showcasegameicons{tarot-01-the-magician}{3641}
\showcasegameicons{tarot-02-the-high-priestess}{3642}
\showcasegameicons{tarot-03-the-empress}{3643}
\showcasegameicons{tarot-04-the-emperor}{3644}
\showcasegameicons{tarot-05-the-hierophant}{3645}
\showcasegameicons{tarot-06-the-lovers}{3646}
\showcasegameicons{tarot-07-the-chariot}{3647}
\showcasegameicons{tarot-08-strength}{3648}
\showcasegameicons{tarot-09-the-hermit}{3649}
\showcasegameicons{tarot-10-wheel-of-fortune}{3650}
\showcasegameicons{tarot-11-justice}{3651}
\showcasegameicons{tarot-12-the-hanged-man}{3652}
\showcasegameicons{tarot-13-death}{3653}
\showcasegameicons{tarot-14-temperance}{3654}
\showcasegameicons{tarot-15-the-devil}{3655}
\showcasegameicons{tarot-16-the-tower}{3656}
\showcasegameicons{tarot-17-the-star}{3657}
\showcasegameicons{tarot-18-the-moon}{3658}
\showcasegameicons{tarot-19-the-sun}{3659}
\showcasegameicons{tarot-20-judgement}{3660}
\showcasegameicons{tarot-21-the-world}{3661}
\showcasegameicons{tattered-banner}{3662}
\showcasegameicons{taurus}{3663}
\showcasegameicons{tavern-sign}{3664}
\showcasegameicons{teacher}{3665}
\showcasegameicons{team-downgrade}{3666}
\showcasegameicons{team-idea}{3667}
\showcasegameicons{team-upgrade}{3668}
\showcasegameicons{teapot-leaves}{3669}
\showcasegameicons{teapot}{3670}
\showcasegameicons{tear-tracks}{3671}
\showcasegameicons{tearing}{3672}
\showcasegameicons{tec-9}{3673}
\showcasegameicons{techno-heart}{3674}
\showcasegameicons{tee-pipe}{3675}
\showcasegameicons{telefrag}{3676}
\showcasegameicons{telepathy}{3677}
\showcasegameicons{teleport}{3678}
\showcasegameicons{telescope}{3679}
\showcasegameicons{telescopic-baton}{3680}
\showcasegameicons{teller-mine}{3681}
\showcasegameicons{templar-eye}{3682}
\showcasegameicons{templar-heart}{3683}
\showcasegameicons{templar-shield}{3684}
\showcasegameicons{temple-door}{3685}
\showcasegameicons{temple-gate}{3686}
\showcasegameicons{temporary-shield}{3687}
\showcasegameicons{temptation}{3688}
\showcasegameicons{tennis-ball}{3689}
\showcasegameicons{tennis-court}{3690}
\showcasegameicons{tennis-racket}{3691}
\showcasegameicons{tension-snowflake}{3692}
\showcasegameicons{tentacle-heart}{3693}
\showcasegameicons{tentacle-strike}{3694}
\showcasegameicons{tentacles-barrier}{3695}
\showcasegameicons{tentacles-skull}{3696}
\showcasegameicons{tentacurl}{3697}
\showcasegameicons{terror}{3698}
\showcasegameicons{tesla-coil}{3699}
\showcasegameicons{tesla-turret}{3700}
\showcasegameicons{tesla}{3701}
\showcasegameicons{test-tube-held}{3702}
\showcasegameicons{test-tube-rack}{3703}
\showcasegameicons{test-tubes}{3704}
\showcasegameicons{texas}{3705}
\showcasegameicons{theater-curtains}{3706}
\showcasegameicons{theater}{3707}
\showcasegameicons{thermometer-cold}{3708}
\showcasegameicons{thermometer-hot}{3709}
\showcasegameicons{thermometer-scale}{3710}
\showcasegameicons{think}{3711}
\showcasegameicons{third-eye}{3712}
\showcasegameicons{thompson-m1}{3713}
\showcasegameicons{thompson-m1928}{3714}
\showcasegameicons{thor-fist}{3715}
\showcasegameicons{thor-hammer}{3716}
\showcasegameicons{thorn-helix}{3717}
\showcasegameicons{thorned-arrow}{3718}
\showcasegameicons{thorny-tentacle}{3719}
\showcasegameicons{thorny-vine}{3720}
\showcasegameicons{thought-bubble}{3721}
\showcasegameicons{three-burning-balls}{3722}
\showcasegameicons{three-friends}{3723}
\showcasegameicons{three-keys}{3724}
\showcasegameicons{three-leaves}{3725}
\showcasegameicons{three-pointed-shuriken}{3726}
\showcasegameicons{throne-king}{3727}
\showcasegameicons{throwing-ball}{3728}
\showcasegameicons{thrown-charcoal}{3729}
\showcasegameicons{thrown-daggers}{3730}
\showcasegameicons{thrown-knife}{3731}
\showcasegameicons{thrown-spear}{3732}
\showcasegameicons{thrust-bend}{3733}
\showcasegameicons{thrust}{3734}
\showcasegameicons{thumb-down}{3735}
\showcasegameicons{thumb-up}{3736}
\showcasegameicons{thunder-blade}{3737}
\showcasegameicons{thunder-skull}{3738}
\showcasegameicons{thunder-struck}{3739}
\showcasegameicons{thunderball}{3740}
\showcasegameicons{thwomp}{3741}
\showcasegameicons{tiara}{3742}
\showcasegameicons{tic-tac-toe}{3743}
\showcasegameicons{tick}{3744}
\showcasegameicons{ticket}{3745}
\showcasegameicons{tie}{3746}
\showcasegameicons{tied-scroll}{3747}
\showcasegameicons{tiger-head}{3748}
\showcasegameicons{tiger}{3749}
\showcasegameicons{tightrope}{3750}
\showcasegameicons{time-bomb}{3751}
\showcasegameicons{time-dynamite}{3752}
\showcasegameicons{time-synchronization}{3753}
\showcasegameicons{time-trap}{3754}
\showcasegameicons{tinker}{3755}
\showcasegameicons{tipi}{3756}
\showcasegameicons{tire-iron-cross}{3757}
\showcasegameicons{tire-iron}{3758}
\showcasegameicons{tire-tracks}{3759}
\showcasegameicons{tired-eye}{3760}
\showcasegameicons{toad-teeth}{3761}
\showcasegameicons{toaster}{3762}
\showcasegameicons{toggles}{3763}
\showcasegameicons{token}{3764}
\showcasegameicons{tomahawk}{3765}
\showcasegameicons{tomato}{3766}
\showcasegameicons{tombstone}{3767}
\showcasegameicons{tongue}{3768}
\showcasegameicons{toolbox}{3769}
\showcasegameicons{tooth}{3770}
\showcasegameicons{toothbrush}{3771}
\showcasegameicons{top-hat}{3772}
\showcasegameicons{top-paw}{3773}
\showcasegameicons{topaz}{3774}
\showcasegameicons{torc}{3775}
\showcasegameicons{torch}{3776}
\showcasegameicons{tornado-discs}{3777}
\showcasegameicons{tornado}{3778}
\showcasegameicons{torpedo}{3779}
\showcasegameicons{tortoise}{3780}
\showcasegameicons{totem-head}{3781}
\showcasegameicons{totem-mask}{3782}
\showcasegameicons{totem}{3783}
\showcasegameicons{toucan}{3784}
\showcasegameicons{tow-truck}{3785}
\showcasegameicons{towel}{3786}
\showcasegameicons{tower-bridge}{3787}
\showcasegameicons{tower-fall}{3788}
\showcasegameicons{tower-flag}{3789}
\showcasegameicons{toy-mallet}{3790}
\showcasegameicons{tracked-robot}{3791}
\showcasegameicons{trade}{3792}
\showcasegameicons{traffic-cone}{3793}
\showcasegameicons{traffic-lights-green}{3794}
\showcasegameicons{traffic-lights-orange}{3795}
\showcasegameicons{traffic-lights-ready-to-go}{3796}
\showcasegameicons{traffic-lights-red}{3797}
\showcasegameicons{trail}{3798}
\showcasegameicons{trample}{3799}
\showcasegameicons{transform}{3800}
\showcasegameicons{transfuse}{3801}
\showcasegameicons{transparent-slime}{3802}
\showcasegameicons{transparent-tubes}{3803}
\showcasegameicons{transportation-rings}{3804}
\showcasegameicons{trap-mask}{3805}
\showcasegameicons{trash-can}{3806}
\showcasegameicons{travel-dress}{3807}
\showcasegameicons{tread}{3808}
\showcasegameicons{treasure-map}{3809}
\showcasegameicons{trebuchet}{3810}
\showcasegameicons{tree-beehive}{3811}
\showcasegameicons{tree-branch}{3812}
\showcasegameicons{tree-door}{3813}
\showcasegameicons{tree-face}{3814}
\showcasegameicons{tree-growth}{3815}
\showcasegameicons{tree-roots}{3816}
\showcasegameicons{tree-swing}{3817}
\showcasegameicons{treehouse}{3818}
\showcasegameicons{trefoil-lily}{3819}
\showcasegameicons{trefoil-shuriken}{3820}
\showcasegameicons{trench-assault}{3821}
\showcasegameicons{trench-body-armor}{3822}
\showcasegameicons{trench-knife}{3823}
\showcasegameicons{trench-spade}{3824}
\showcasegameicons{triangle-target}{3825}
\showcasegameicons{tribal-gear}{3826}
\showcasegameicons{tribal-mask}{3827}
\showcasegameicons{tribal-pendant}{3828}
\showcasegameicons{tribal-shield}{3829}
\showcasegameicons{tribunal-jury}{3830}
\showcasegameicons{triceratops-head}{3831}
\showcasegameicons{trident-shield}{3832}
\showcasegameicons{trident}{3833}
\showcasegameicons{triforce}{3834}
\showcasegameicons{trigger-hurt}{3835}
\showcasegameicons{trilobite}{3836}
\showcasegameicons{trinacria}{3837}
\showcasegameicons{triorb}{3838}
\showcasegameicons{triple-beak}{3839}
\showcasegameicons{triple-claws}{3840}
\showcasegameicons{triple-corn}{3841}
\showcasegameicons{triple-gate}{3842}
\showcasegameicons{triple-lock}{3843}
\showcasegameicons{triple-needle}{3844}
\showcasegameicons{triple-plier}{3845}
\showcasegameicons{triple-scratches}{3846}
\showcasegameicons{triple-shells}{3847}
\showcasegameicons{triple-skulls}{3848}
\showcasegameicons{triple-yin}{3849}
\showcasegameicons{tripwire}{3850}
\showcasegameicons{triquetra}{3851}
\showcasegameicons{trireme}{3852}
\showcasegameicons{triton-head}{3853}
\showcasegameicons{troglodyte}{3854}
\showcasegameicons{trojan-horse}{3855}
\showcasegameicons{troll}{3856}
\showcasegameicons{trombone}{3857}
\showcasegameicons{tron-arrow}{3858}
\showcasegameicons{trophies-shelf}{3859}
\showcasegameicons{trophy-cup}{3860}
\showcasegameicons{trophy}{3861}
\showcasegameicons{tropical-fish}{3862}
\showcasegameicons{trousers}{3863}
\showcasegameicons{trowel}{3864}
\showcasegameicons{truce}{3865}
\showcasegameicons{truck}{3866}
\showcasegameicons{trumpet-flag}{3867}
\showcasegameicons{trumpet}{3868}
\showcasegameicons{trunk-mushroom}{3869}
\showcasegameicons{tuba}{3870}
\showcasegameicons{tumbleweed}{3871}
\showcasegameicons{tumor}{3872}
\showcasegameicons{tumulus}{3873}
\showcasegameicons{tune-pitch}{3874}
\showcasegameicons{tunisia}{3875}
\showcasegameicons{turban}{3876}
\showcasegameicons{turbine}{3877}
\showcasegameicons{turd}{3878}
\showcasegameicons{turnstile}{3879}
\showcasegameicons{turret}{3880}
\showcasegameicons{turtle-shell}{3881}
\showcasegameicons{turtle}{3882}
\showcasegameicons{tusks-flag}{3883}
\showcasegameicons{tv-remote}{3884}
\showcasegameicons{tv-tower}{3885}
\showcasegameicons{tv}{3886}
\showcasegameicons{twin-shell}{3887}
\showcasegameicons{twirl-center}{3888}
\showcasegameicons{twirly-flower}{3889}
\showcasegameicons{twister}{3890}
\showcasegameicons{two-coins}{3891}
\showcasegameicons{two-feathers}{3892}
\showcasegameicons{two-handed-sword}{3893}
\showcasegameicons{two-shadows}{3894}
\showcasegameicons{tyre}{3895}
\showcasegameicons{ubisoft-sun}{3896}
\showcasegameicons{udder}{3897}
\showcasegameicons{ufo}{3898}
\showcasegameicons{ultrasound}{3899}
\showcasegameicons{uluru}{3900}
\showcasegameicons{umbrella-bayonet}{3901}
\showcasegameicons{umbrella}{3902}
\showcasegameicons{unbalanced}{3903}
\showcasegameicons{uncertainty}{3904}
\showcasegameicons{underground-cave}{3905}
\showcasegameicons{underhand}{3906}
\showcasegameicons{underwear-shorts}{3907}
\showcasegameicons{underwear}{3908}
\showcasegameicons{unfriendly-fire}{3909}
\showcasegameicons{unicorn}{3910}
\showcasegameicons{unicycle}{3911}
\showcasegameicons{union-jack}{3912}
\showcasegameicons{unlit-bomb}{3913}
\showcasegameicons{unlit-candelabra}{3914}
\showcasegameicons{unlocking}{3915}
\showcasegameicons{unplugged}{3916}
\showcasegameicons{unstable-orb}{3917}
\showcasegameicons{unstable-projectile}{3918}
\showcasegameicons{up-card}{3919}
\showcasegameicons{upgrade}{3920}
\showcasegameicons{uprising}{3921}
\showcasegameicons{ursa-major}{3922}
\showcasegameicons{uruguay}{3923}
\showcasegameicons{usa-flag}{3924}
\showcasegameicons{usable}{3925}
\showcasegameicons{usb-key}{3926}
\showcasegameicons{ushanka}{3927}
\showcasegameicons{uzi}{3928}
\showcasegameicons{vacuum-cleaner}{3929}
\showcasegameicons{valley}{3930}
\showcasegameicons{valve}{3931}
\showcasegameicons{vampire-cape}{3932}
\showcasegameicons{vampire-dracula}{3933}
\showcasegameicons{van-damme-split}{3934}
\showcasegameicons{vanilla-flower}{3935}
\showcasegameicons{velocipede}{3936}
\showcasegameicons{velociraptor-tracks}{3937}
\showcasegameicons{velociraptor}{3938}
\showcasegameicons{vending-machine}{3939}
\showcasegameicons{venezuela}{3940}
\showcasegameicons{venus-flytrap}{3941}
\showcasegameicons{venus-of-willendorf}{3942}
\showcasegameicons{vertical-banner}{3943}
\showcasegameicons{vertical-flip}{3944}
\showcasegameicons{vhs}{3945}
\showcasegameicons{vial}{3946}
\showcasegameicons{vibrating-ball}{3947}
\showcasegameicons{vibrating-shield}{3948}
\showcasegameicons{vibrating-smartphone}{3949}
\showcasegameicons{video-camera}{3950}
\showcasegameicons{video-conference}{3951}
\showcasegameicons{viking-church}{3952}
\showcasegameicons{viking-head}{3953}
\showcasegameicons{viking-helmet}{3954}
\showcasegameicons{viking-longhouse}{3955}
\showcasegameicons{viking-shield}{3956}
\showcasegameicons{vile-fluid}{3957}
\showcasegameicons{village}{3958}
\showcasegameicons{vine-flower}{3959}
\showcasegameicons{vine-leaf}{3960}
\showcasegameicons{vine-whip}{3961}
\showcasegameicons{vines}{3962}
\showcasegameicons{vintage-robot}{3963}
\showcasegameicons{viola}{3964}
\showcasegameicons{violin}{3965}
\showcasegameicons{virgo}{3966}
\showcasegameicons{virtual-marker}{3967}
\showcasegameicons{virus}{3968}
\showcasegameicons{visored-helm}{3969}
\showcasegameicons{vitruvian-man}{3970}
\showcasegameicons{volcano}{3971}
\showcasegameicons{volleyball-ball}{3972}
\showcasegameicons{vomiting}{3973}
\showcasegameicons{voodoo-doll}{3974}
\showcasegameicons{vortex}{3975}
\showcasegameicons{vote}{3976}
\showcasegameicons{vr-headset}{3977}
\showcasegameicons{vulture}{3978}
\showcasegameicons{vuvuzelas}{3979}
\showcasegameicons{walk}{3980}
\showcasegameicons{walkie-talkie}{3981}
\showcasegameicons{walking-boot}{3982}
\showcasegameicons{walking-scout}{3983}
\showcasegameicons{walking-turret}{3984}
\showcasegameicons{wall-light}{3985}
\showcasegameicons{wallet}{3986}
\showcasegameicons{walrus-head}{3987}
\showcasegameicons{walther-ppk}{3988}
\showcasegameicons{wanted-reward}{3989}
\showcasegameicons{war-axe}{3990}
\showcasegameicons{war-bonnet}{3991}
\showcasegameicons{war-pick}{3992}
\showcasegameicons{warehouse}{3993}
\showcasegameicons{warhammer}{3994}
\showcasegameicons{warlock-eye}{3995}
\showcasegameicons{warlock-hood}{3996}
\showcasegameicons{warlord-helmet}{3997}
\showcasegameicons{warp-pipe}{3998}
\showcasegameicons{washing-machine}{3999}
\showcasegameicons{wasp-sting}{4000}
\showcasegameicons{watch}{4001}
\showcasegameicons{watchtower}{4002}
\showcasegameicons{water-bolt}{4003}
\showcasegameicons{water-bottle}{4004}
\showcasegameicons{water-diviner-stick}{4005}
\showcasegameicons{water-drop}{4006}
\showcasegameicons{water-flask}{4007}
\showcasegameicons{water-fountain}{4008}
\showcasegameicons{water-gallon}{4009}
\showcasegameicons{water-gun}{4010}
\showcasegameicons{water-mill}{4011}
\showcasegameicons{water-polo}{4012}
\showcasegameicons{water-recycling}{4013}
\showcasegameicons{water-splash}{4014}
\showcasegameicons{water-tank}{4015}
\showcasegameicons{water-tower}{4016}
\showcasegameicons{waterfall}{4017}
\showcasegameicons{watering-can}{4018}
\showcasegameicons{watermelon}{4019}
\showcasegameicons{waterskin}{4020}
\showcasegameicons{wave-crest}{4021}
\showcasegameicons{wave-strike}{4022}
\showcasegameicons{wave-surfer}{4023}
\showcasegameicons{waves}{4024}
\showcasegameicons{wavy-chains}{4025}
\showcasegameicons{wavy-itinerary}{4026}
\showcasegameicons{wax-seal}{4027}
\showcasegameicons{wax-tablet}{4028}
\showcasegameicons{web-spit}{4029}
\showcasegameicons{weight-crush}{4030}
\showcasegameicons{weight-lifting-down}{4031}
\showcasegameicons{weight-lifting-up}{4032}
\showcasegameicons{weight-scale}{4033}
\showcasegameicons{weight}{4034}
\showcasegameicons{well}{4035}
\showcasegameicons{werewolf}{4036}
\showcasegameicons{western-hat}{4037}
\showcasegameicons{whale-tail}{4038}
\showcasegameicons{wheat}{4039}
\showcasegameicons{wheelbarrow}{4040}
\showcasegameicons{whip}{4041}
\showcasegameicons{whiplash}{4042}
\showcasegameicons{whirlpool-shuriken}{4043}
\showcasegameicons{whirlwind}{4044}
\showcasegameicons{whisk}{4045}
\showcasegameicons{whistle}{4046}
\showcasegameicons{white-book}{4047}
\showcasegameicons{white-tower}{4048}
\showcasegameicons{wide-arrow-dunk}{4049}
\showcasegameicons{wifi-router}{4050}
\showcasegameicons{wildfires}{4051}
\showcasegameicons{william-tell-skull}{4052}
\showcasegameicons{william-tell}{4053}
\showcasegameicons{willow-tree}{4054}
\showcasegameicons{winchester-rifle}{4055}
\showcasegameicons{wind-hole}{4056}
\showcasegameicons{wind-slap}{4057}
\showcasegameicons{wind-turbine}{4058}
\showcasegameicons{windchimes}{4059}
\showcasegameicons{windmill}{4060}
\showcasegameicons{window-bars}{4061}
\showcasegameicons{window}{4062}
\showcasegameicons{windpump}{4063}
\showcasegameicons{windsock}{4064}
\showcasegameicons{windy-stripes}{4065}
\showcasegameicons{wine-bottle}{4066}
\showcasegameicons{wine-glass}{4067}
\showcasegameicons{wing-cloak}{4068}
\showcasegameicons{winged-arrow}{4069}
\showcasegameicons{winged-emblem}{4070}
\showcasegameicons{winged-leg}{4071}
\showcasegameicons{winged-scepter}{4072}
\showcasegameicons{winged-shield}{4073}
\showcasegameicons{winged-sword}{4074}
\showcasegameicons{wingfoot}{4075}
\showcasegameicons{winter-gloves}{4076}
\showcasegameicons{winter-hat}{4077}
\showcasegameicons{wire-coil}{4078}
\showcasegameicons{wireframe-globe}{4079}
\showcasegameicons{wisdom}{4080}
\showcasegameicons{witch-face}{4081}
\showcasegameicons{witch-flight}{4082}
\showcasegameicons{wizard-face}{4083}
\showcasegameicons{wizard-staff}{4084}
\showcasegameicons{wok}{4085}
\showcasegameicons{wolf-head}{4086}
\showcasegameicons{wolf-howl}{4087}
\showcasegameicons{wolf-trap}{4088}
\showcasegameicons{wolverine-claws}{4089}
\showcasegameicons{woman-elf-face}{4090}
\showcasegameicons{wood-axe}{4091}
\showcasegameicons{wood-beam}{4092}
\showcasegameicons{wood-cabin}{4093}
\showcasegameicons{wood-canoe}{4094}
\showcasegameicons{wood-club}{4095}
\showcasegameicons{wood-frame}{4096}
\showcasegameicons{wood-pile}{4097}
\showcasegameicons{wood-stick}{4098}
\showcasegameicons{wooden-chair}{4099}
\showcasegameicons{wooden-clogs}{4100}
\showcasegameicons{wooden-crate}{4101}
\showcasegameicons{wooden-door}{4102}
\showcasegameicons{wooden-fence}{4103}
\showcasegameicons{wooden-helmet}{4104}
\showcasegameicons{wooden-pegleg}{4105}
\showcasegameicons{wooden-pier}{4106}
\showcasegameicons{wooden-sign}{4107}
\showcasegameicons{wool}{4108}
\showcasegameicons{world}{4109}
\showcasegameicons{worm-mouth}{4110}
\showcasegameicons{worms}{4111}
\showcasegameicons{worried-eyes}{4112}
\showcasegameicons{wrapped-heart}{4113}
\showcasegameicons{wrapped-sweet}{4114}
\showcasegameicons{wrapping-star}{4115}
\showcasegameicons{wrecking-ball}{4116}
\showcasegameicons{wrench}{4117}
\showcasegameicons{wyvern}{4118}
\showcasegameicons{xylophone}{4119}
\showcasegameicons{yarn}{4120}
\showcasegameicons{yin-yang}{4121}
\showcasegameicons{yunluo}{4122}
\showcasegameicons{z-brick}{4123}
\showcasegameicons{zat-gun}{4124}
\showcasegameicons{zebra-shield}{4125}
\showcasegameicons{zeppelin}{4126}
\showcasegameicons{zeus-sword}{4127}
\showcasegameicons{zig-arrow}{4128}
\showcasegameicons{zigzag-cage}{4129}
\showcasegameicons{zigzag-hieroglyph}{4130}
\showcasegameicons{zigzag-leaf}{4131}
\showcasegameicons{zigzag-tune}{4132}
\showcasegameicons{zipper}{4133}
\end{gameiconcase}


\newpage

\section{List of badges}

The \textsf{PDF} file is \texttt{gameicons-badges-all}.

\begin{gameiconbadgecase}
\showcasegameiconsbadge{anvil}{1}
\showcasegameiconsbadge{arrow-down}{2}
\showcasegameiconsbadge{arrow-up}{3}
\showcasegameiconsbadge{blank}{4}
\showcasegameiconsbadge{body}{5}
\showcasegameiconsbadge{bolt}{6}
\showcasegameiconsbadge{building}{7}
\showcasegameiconsbadge{cloud}{8}
\showcasegameiconsbadge{club}{9}
\showcasegameiconsbadge{cog}{10}
\showcasegameiconsbadge{coins}{11}
\showcasegameiconsbadge{corporal}{12}
\showcasegameiconsbadge{crosshair}{13}
\showcasegameiconsbadge{crown}{14}
\showcasegameiconsbadge{diamond}{15}
\showcasegameiconsbadge{exclamation}{16}
\showcasegameiconsbadge{eye}{17}
\showcasegameiconsbadge{fire}{18}
\showcasegameiconsbadge{flask}{19}
\showcasegameiconsbadge{happy}{20}
\showcasegameiconsbadge{heart}{21}
\showcasegameiconsbadge{ice}{22}
\showcasegameiconsbadge{justice}{23}
\showcasegameiconsbadge{leaf}{24}
\showcasegameiconsbadge{lock}{25}
\showcasegameiconsbadge{minus}{26}
\showcasegameiconsbadge{moon}{27}
\showcasegameiconsbadge{mountain}{28}
\showcasegameiconsbadge{multiply}{29}
\showcasegameiconsbadge{music}{30}
\showcasegameiconsbadge{number-0}{31}
\showcasegameiconsbadge{number-1}{32}
\showcasegameiconsbadge{number-2}{33}
\showcasegameiconsbadge{number-3}{34}
\showcasegameiconsbadge{number-4}{35}
\showcasegameiconsbadge{number-5}{36}
\showcasegameiconsbadge{number-6}{37}
\showcasegameiconsbadge{number-7}{38}
\showcasegameiconsbadge{number-8}{39}
\showcasegameiconsbadge{number-9}{40}
\showcasegameiconsbadge{paw}{41}
\showcasegameiconsbadge{pencil}{42}
\showcasegameiconsbadge{percentage}{43}
\showcasegameiconsbadge{plus}{44}
\showcasegameiconsbadge{private}{45}
\showcasegameiconsbadge{purse}{46}
\showcasegameiconsbadge{question}{47}
\showcasegameiconsbadge{radioactive}{48}
\showcasegameiconsbadge{sad}{49}
\showcasegameiconsbadge{sergeant}{50}
\showcasegameiconsbadge{shield}{51}
\showcasegameiconsbadge{skull}{52}
\showcasegameiconsbadge{spade}{53}
\showcasegameiconsbadge{star}{54}
\showcasegameiconsbadge{strength}{55}
\showcasegameiconsbadge{summon}{56}
\showcasegameiconsbadge{sun}{57}
\showcasegameiconsbadge{time}{58}
\showcasegameiconsbadge{water}{59}
\end{gameiconbadgecase}


\end{document}
